\documentclass[
    a4paper,
    %uninadraft,
    %colorlinks,
    blacklinks,
    %palatino,
    libertine,
    %mathpazo,
    %tgschola,
    %tgscholamathpazo,
    %alegreya
]{uninathesis}

%\AtBeginDocument{\addtocontents{toc}{\protect\thispagestyle{empty}}} %remove page number from toc

% ---------------------------------------------------------------------------------------
% Lingua: la tesi e' in italiano.
% ---------------------------------------------------------------------------------------
\usepackage[italian]{babel}
\usepackage[T1]{fontenc}

% Math
\RequirePackage{mathtools}
\RequirePackage{mathrsfs} %need this for mathscr
\let\openbox\relax
\RequirePackage{amsthm}
\RequirePackage{stmaryrd} %for \llbracket

\theoremstyle{plain}
\newtheorem{theorem}{Teorema}
\newtheorem{proposition}[theorem]{Proposizione}
\newtheorem{corollary}[theorem]{Corollario}
\newtheorem{lemma}[theorem]{Lemma}

\theoremstyle{definition}
\newtheorem{definition}[theorem]{Definizione}
\newtheorem{example}[theorem]{Esempio}
\newtheorem{observation}[theorem]{Osservazione}

\theoremstyle{remark}
\newtheorem{remark}[theorem]{Nota}

% Tikz - may be removed if unneeded
\RequirePackage{tikz}
\RequirePackage{pgfplots}
\pgfplotsset{compat=1.16}
\usetikzlibrary{positioning,automata,calc,fit,shapes,backgrounds,arrows.meta}
\usepackage{makecell}
\usepackage{multirow}   % celle su piu' righe nelle tabelle

% Miscellanea Packages
\RequirePackage[utf8]{inputenc}
\usepackage{ marvosym } %/Lightning
\usepackage{subcaption} %subfigures
\usepackage{booktabs}   % tabelle professionali: \toprule \midrule \bottomrule
\usepackage{longtable}  % tabelle lunghe, che si spezzano su piu' pagine (appendice B)
\usepackage{xurl}      % permette di spezzare le URL lunghe in bibliografia
\usepackage{siunitx}    % numeri allineati e unita' di misura
\sisetup{group-separator={\,},group-minimum-digits=4,output-decimal-marker={,}}
\usepackage{xcolor}
\usepackage{enumitem}

% read configuration file
%double quote command
\newcommand\dquote[1]{``#1''}
\newcommand\dquoteit[1]{``\emph{#1}''}

%single quote command
\newcommand\squote[1]{`#1'}

% Segno di percentuale fuori da \SI (siunitx non lo definisce da solo)
\providecommand{\percent}{\%}

% =======================================================================================
% Macro per la notazione del lavoro.
% =======================================================================================

% RFD, LHS, RHS
\newcommand{\rfd}{\ensuremath{\mathit{rfd}}}
\newcommand{\LHS}{\ensuremath{\mathrm{LHS}}}
\newcommand{\RHS}{\ensuremath{\mathrm{RHS}}}
\newcommand{\NullCols}{\ensuremath{\mathrm{NullCols}}}
\newcommand{\thr}{\ensuremath{\theta}}

% Le tre condizioni derivate dal bytecode
\newcommand{\Vone}{\ensuremath{[\mathrm{V1}]}}
\newcommand{\Vtwo}{\ensuremath{[\mathrm{V2}]}}
\newcommand{\Vthree}{\ensuremath{[\mathrm{V3}]}}

% Strategie citate nel testo: monospazio per restare coerenti con il codice
\newcommand{\str}[1]{\texttt{#1}}

% Termini tecnici ricorrenti
\newcommand{\renuver}{\textsc{Renuver}}
\newcommand{\jar}{\texttt{RENUVER\_EDBT.jar}}

% Tabelle: intestazione compatta
\newcommand{\hd}[1]{\textbf{#1}}
 %can write your macros here

% Bibliography
\RequirePackage{csquotes}
\RequirePackage[sorting=none,maxbibnames=4,backend=biber]{biblatex}
\addbibresource{_bib/bibliography.bib}
% biblatex: consente di spezzare le URL e i DOI anche in corrispondenza di cifre e lettere,
% evitando che una voce lunga sbordi dal margine destro.
\setcounter{biburlnumpenalty}{100}
\setcounter{biburlucpenalty}{100}
\setcounter{biburllcpenalty}{100}

% Front page (package uninafrontespizio)
% ---------------------------------------------------------------------------------------
% ATTENZIONE: i campi contrassegnati con  <<<  vanno completati a cura del candidato.
% ---------------------------------------------------------------------------------------
\Universita{Università degli Studi di Napoli Federico II}
\Facolta{Scuola Politecnica e delle Scienze di Base}
\Dipartimento{Dipartimento di Ingegneria Elettrica e Tecnologie dell'Informazione}
\CorsoDiLaurea{Corso di Laurea in Informatica}                       % <<< triennale
\AnnoAccademico{2025--2026}                                          % <<< da confermare
\Titolo{Semantic Pruning di metadati per l'imputazione di valori mancanti:\\ analisi, metodologia e valutazione su RENUVER}
\Relatore{Prof.\ \textsc{Gennaro Ventrone}}                              % <<< da completare
\RelatoreLabel{Relatore}
\Candidato{\textsc{Gennaro Ventrone}}                                    % <<< da completare
\Matricola{N86004074}                                                   % <<< da completare
\Logo{logo-federico-II.pdf}
\LogoWidth{3.5cm}
\LogoPosition{below-uni}

\usepackage[title,titletoc]{appendix}

% ---------------------------------------------------------------------------------------
% Riferimenti incrociati in italiano.
% hyperref nomina le unita' in inglese ("Table", "section", "chapter", "Appendix",
% "paragraph") e le traduzioni di babel non le coprono tutte: senza queste ridefinizioni
% il testo italiano conterrebbe "la Table 5.2" o "nel chapter 6".
% ---------------------------------------------------------------------------------------
\addto\captionsitalian{%
    \def\chapterautorefname{capitolo}%
    \def\sectionautorefname{sezione}%
    \def\subsectionautorefname{sottosezione}%
    \def\paragraphautorefname{paragrafo}%
    \def\tableautorefname{Tabella}%
    \def\figureautorefname{Figura}%
    \def\appendixautorefname{Appendice}%
}

% Metadati del PDF
\hypersetup{
    pdftitle={Semantic Pruning di metadati per l'imputazione di valori mancanti},
    pdfauthor={Nome Cognome},
    pdfsubject={Tesi di Laurea in Informatica}
}

\begin{document}

    \pagestyle{empty}
    \frontmatter
    \makefrontespizio~\clearpage

    \pagestyle{backmatter}
    % =======================================================================================
% Sommario
% =======================================================================================
\chapter*{Sommario}

L'imputazione di valori mancanti basata su \emph{Relaxed Functional Dependencies} (RFD)
affronta il problema dell'incompletezza dei dati sfruttando regole estratte dal dataset stesso.
Il costo computazionale di questi approcci cresce con il numero di regole disponibili, e una
parte di esse è ridondante: da qui l'interesse per il \emph{semantic pruning}, la selezione
delle RFD effettivamente rilevanti.

Questo lavoro studia il pruning semantico di un insieme di metadati applicato a \renuver{}, un
algoritmo di imputazione basato su RFD. Dall'analisi del bytecode si derivano tre condizioni
esatte: \Vone{}, che caratterizza le regole che possono esercitare un veto sulle imputazioni;
\Vtwo{}, che individua una classe di regole la cui potatura è dimostrabilmente priva di effetti;
\Vthree{}, che isola le regole provabilmente inerti. Da queste discende un teorema negativo: con
potatura esterna, nessuna riduzione non banale può essere contemporaneamente sicura rispetto alla
generazione dei candidati e al potere di veto.

La condizione \Vtwo{} è stata verificata empiricamente su \num{1242321} terne con zero
discrepanze, e la sua applicazione produce un output identico a quello del sistema originale su
\num{113} configurazioni su \num{113}, su quattro famiglie di dataset (\num{111} delle tre
famiglie principali e \num{2} di una famiglia indipendente). Il teorema separa ciò che è
dimostrato --- la caratterizzazione delle rimozioni certificabili --- da ciò che è misurato, e
mostra che la riduzione garantita non supera il \SI{20.9}{\percent} sui dati esaminati. Un'euristica di potatura
aggressiva, denominata \str{A5b}, riduce le regole dell'\SI{84.2}{\percent} su Bridges migliorando
al contempo il risultato di \num{69} imputazioni corrette e \num{72} errori in meno; sugli altri
dataset la stessa strategia danneggia il risultato. Il lavoro identifica il meccanismo che governa
questa differenza --- il guadagno è il prodotto della quota di imputazioni ottenute per via esatta
e del numero di imputazioni effettuate --- e ne ricava un criterio diagnostico misurabile prima di
potare. Su Glass e Restaurant una taratura più conservativa della medesima strategia non perde che
una imputazione corretta su \num{97} configurazioni. Il lavoro mostra infine che la riduzione
delle regole non implica un risparmio di tempo: la strategia che ne rimuove di più è anche quella
che rallenta di più, e il tempo dipende dalla composizione dell'insieme potato più che dalla sua
dimensione.


    \clearpage
    \tableofcontents

    \mainmatter

    \newpage\pagestyle{introduction}

\unnumberedchapter{Introduzione}

\unnumberedsection{Contesto e motivazioni}

La qualità dei dati è un requisito operativo prima che teorico: un'analisi condotta su
informazioni incomplete produce conclusioni che nessun metodo statistico può recuperare. Quando
un attributo di una tupla è privo di valore, la scelta non è fra ``imputare'' e ``non imputare'',
ma fra imputare in modo informato e scartare il dato. È questa la ragione dell'interesse per la
\emph{data imputation}: ricostruire il valore mancante a partire dal resto dell'informazione
disponibile.

Tra gli approcci recenti, quelli basati su \emph{Relaxed Functional Dependencies} (RFD) hanno una
proprietà che li distingue: le regole di imputazione non vengono fornite dall'esterno, ma estratte
dal dataset stesso. Una RFD esprime un vincolo di dipendenza approssimato fra attributi ---
\emph{se due tuple concordano entro certe soglie sugli attributi di sinistra, allora concordano
entro una soglia sull'attributo di destra} --- e può quindi essere usata sia per generare
candidati all'imputazione sia per rifiutare quelli incoerenti. Il sistema \renuver{} \cite{breve2022renuver}
è un rappresentante di questa famiglia.

Il prezzo di questa ricchezza è computazionale: la fase di estrazione produce tipicamente migliaia
di regole, e il costo dell'imputazione cresce con il loro numero. Molte di esse, però, non
influenzano il risultato. Da qui il problema affrontato in questo lavoro: \emph{si possono
selezionare le RFD rilevanti senza alterare il comportamento del sistema?} È la domanda del
\emph{semantic pruning} applicato a un insieme di metadati.

\unnumberedsection{Problema affrontato}

Il sistema \renuver{} riceve in ingresso, oltre al dataset, un file di RFD. L'ipotesi di partenza
--- suggerita dalla struttura del sistema --- è che una parte di quel file sia ridondante e che
rimuoverla riduca il tempo di esecuzione senza degradare la qualità delle imputazioni.

Il problema si rivela però più sottile di così, per una ragione che questo lavoro porta alla luce
in modo formale: nel sistema studiato le RFD svolgono \emph{due ruoli distinti}, e i due ruoli
impongono requisiti opposti. Una RFD serve a \emph{generare} candidati all'imputazione, ma serve
anche a \emph{vetare} imputazioni che violerebbero il vincolo. Rimuovere una regola inutile per il
primo ruolo non garantisce che essa sia inutile per il secondo, e nulla nel formato del file
distingue i due casi.

Da questa osservazione discendono le due domande che strutturano la tesi:

\begin{enumerate}[label=\textbf{D\arabic*.},leftmargin=2.2em]
    \item \textbf{Che cosa si può potare con garanzia?} Esiste una classe di regole la cui
    rimozione è dimostrabilmente priva di conseguenze, e quanto è grande?
    \item \textbf{Che cosa si guadagna quando la garanzia si abbandona?} Se si rimuovono regole
    anche a costo di alterare il comportamento, quali sono i guadagni e a quali condizioni?
\end{enumerate}

\unnumberedsection{Domanda di ricerca}\label{sec:research-question}

Il lavoro muove da una domanda formulata in termini verificabili. Dato l'insieme di RFD che il
sistema riceve in ingresso, \emph{è possibile ridurne il numero almeno della metà mantenendo
l'accuratezza e la copertura entro cinque punti percentuali dalla versione originale, e riducendo
il tempo di esecuzione?} La domanda ha tre esiti distinti --- la riduzione, la qualità delle
imputazioni e il costo --- e la risposta del lavoro non è uniforme sui tre.

Sul primo esito la risposta è \emph{negativa}, e non per insufficienza delle strategie: il
\autoref{chap:theory} dimostra che una potatura che operi sul solo file delle regole e non rinunci
a nessuna garanzia non può superare, sui dataset esaminati, il \SI{20.9}{\percent} di riduzione.
Sul secondo esito la risposta è \emph{condizionata}: la potatura può migliorare il risultato --- su
Bridges lo migliora in modo statisticamente significativo --- ma soltanto quando la composizione
delle imputazioni lascia spazio al miglioramento, e sugli altri dataset danneggia. Sul terzo esito
la risposta è \emph{negativa nel caso generale e positiva in un caso}: la riduzione delle regole
non implica un risparmio di tempo, e la strategia che rimuove di più è quella che rallenta di più.
Il contributo del lavoro è di stabilire questi tre confini e di fornire il criterio con cui
riconoscere, su un dataset nuovo, in quale dei tre regimi ci si trovi.

La \autoref{tab:research-question} anticipa la forma delle tre risposte, con il rimando alle
sezioni in cui ciascuna è argomentata: serve a rendere esplicito fin dall'inizio che cosa il
lavoro sosterrà e che cosa non sosterrà.

\begin{table}[htbp]
\centering
\caption{Risposta alla domanda di ricerca, per esito. La tabella anticipa la struttura
dell'argomentazione; le grandezze sono discusse nelle sezioni indicate.}
\label{tab:research-question}
\footnotesize
\begin{tabular}{@{}p{2.5cm}p{7.3cm}p{2.4cm}@{}}
\toprule
esito & risposta & dove \\
\midrule
riduzione $\ge \SI{50}{\percent}$ \emph{con garanzia} & \textbf{impossibile}: la classe di regole
la cui rimozione è certificabile non supera il \SI{20.9}{\percent} sui dataset esaminati, e il
limite è del criterio, non dell'implementazione & \autoref{chap:theory}, \autoref{sec:maxextra} \\
accuratezza e copertura entro cinque punti & \textbf{condizionata}: il risultato migliora dove la
composizione delle imputazioni ha spazio per migliorare e la copertura non cala, e peggiora dove
una delle due condizioni manca & \autoref{sec:res-a5b}, \autoref{sec:res-boundary} \\
riduzione del tempo di esecuzione & \textbf{non garantita}: il tempo dipende dalla composizione
dell'insieme potato, non dalla sua dimensione; il caso migliore accelera di \num{11} volte, il
peggiore rallenta di \num{251} & \autoref{sec:res-cost} \\
\midrule
\multicolumn{3}{@{}p{12.2cm}@{}}{\footnotesize Il criterio diagnostico della
\autoref{sec:diagnostic} rende operativa la seconda risposta: decide a costo nullo se la leva del
miglioramento esiste, e impone di misurare la copertura prima di adottare la potatura.} \\
\bottomrule
\end{tabular}
\end{table}

\unnumberedsection{Obiettivi del lavoro}

Gli obiettivi, in ordine di priorità:

\begin{itemize}[leftmargin=1.4em]
    \item \textbf{Caratterizzare il sistema.} Derivare dal comportamento effettivo di \renuver{}
    --- non dalla sua documentazione --- le condizioni sotto le quali una RFD può influenzare
    un'imputazione, e tradurle in criteri di potatura verificabili.
    \item \textbf{Definire una metodologia di pruning.} Progettare e implementare strategie di
    selezione, distinguendo esplicitamente quelle che offrono una garanzia da quelle euristiche.
    \item \textbf{Valutare l'impatto.} Misurare l'effetto della potatura sulle prestazioni
    dell'imputazione --- accuratezza, copertura, precisione e tempi --- con un protocollo
    appaiato e test statistici, su più famiglie di dataset.
    \item \textbf{Spiegare i risultati.} Non limitarsi a misurare: individuare il meccanismo che
    determina quando la potatura conviene e quando danneggia.
\end{itemize}

\unnumberedsection{Contributi della tesi}

I contributi originali di questo lavoro sono i seguenti.

\begin{enumerate}[leftmargin=1.6em]
    \item \textbf{Tre condizioni esatte derivate dal bytecode} \Vone{}, \Vtwo{}, \Vthree{},
    che classificano le RFD in base al loro ruolo effettivo nel sistema. \Vtwo{} è la più
    rilevante: identifica le regole il cui potere di veto è implicato da un'altra regola
    sopravvissuta, e la loro rimozione è quindi priva di effetti.

    \item \textbf{Un teorema negativo} (\autoref{chap:theory}) che delimita il campo: se la
    potatura opera sul solo file delle regole e ogni rimozione è certificata con l'informazione
    ivi disponibile, allora le regole rimovibili sono quelle inerti oppure quelle il cui veto è
    implicato da una regola sopravvissuta. Il risultato è quantificato: la classe potabile con
    garanzia non supera il \SI{20.9}{\percent} sui dataset studiati, e si dimostra che tale tetto
    è \emph{intrinseco al criterio} e non un artefatto della sua implementazione
    (\autoref{sec:maxextra}). Il teorema separa ciò che è dimostrato --- la caratterizzazione
    delle rimozioni certificabili --- da ciò che è misurato, e la distinzione è mantenuta in
    tutto il lavoro.

    \item \textbf{Una verifica empirica della condizione \Vtwo{}} (\autoref{sec:res-vetocover})
    su \num{1242321} terne con zero
    discrepanze, e la conferma che la sua applicazione lascia l'output identico al sistema
    originale su \num{113} configurazioni su \num{113}, su quattro famiglie di dataset.

    \item \textbf{Un'euristica di potatura aggressiva} (\str{A5b}, \autoref{sec:res-a5b}) che su Bridges rimuove
    l'\SI{84.2}{\percent} delle regole \emph{migliorando} il risultato (\num{69} imputazioni
    corrette in più, \num{72} errori in meno, \num{11} volte più veloce), e la misurazione del
    fatto che la stessa strategia danneggia gli altri dataset.

    \item \textbf{La spiegazione del perché} (\autoref{sec:mechanism}). Il guadagno della potatura aggressiva si scompone
    nel prodotto di due fattori --- la quota di imputazioni ottenute per via esatta e il numero
    di imputazioni effettuate. La potatura conviene solo quando entrambi crescono. Ne discende
    un criterio diagnostico misurabile prima di potare, che sostituisce la soglia euristica
    precedentemente adottata.

    \item \textbf{Sei difetti documentati} (\autoref{chap:defects}), quattro del materiale
    d'archivio e due del codice sviluppato, con impatto quantificato e procedura di
    riproduzione: con procedura di
    riproduzione: un ramo irraggiungibile nel codice di calcolo della similarità, un
    disallineamento di schema nei dataset Glass, indici di ground truth non ricostruibili nei
    dataset Physician, e una build con istruzioni di debug. A questi si aggiunge un difetto di
    correttezza nel codice sviluppato in questo lavoro, individuato e corretto.

    \item \textbf{La caratterizzazione del costo computazionale} del sistema
    (\autoref{sec:cost}), ottenuta ricompilando il codice originale da sorgenti decompilati ---
    con equivalenza verificata rispetto all'eseguibile di partenza --- e strumentandolo: il
    \SI{99}{\percent} del tempo di esecuzione è nella verifica dei veti, non nella generazione
    dei candidati.

    \item \textbf{La confutazione della proporzionalità fra riduzione e tempo}
    (\autoref{sec:res-cost}). Il lavoro mostra che il tempo di esecuzione dipende dalla
    \emph{composizione} dell'insieme potato e non dalla sua dimensione: la strategia che rimuove
    fino al \SI{98}{\percent} delle regole rallenta l'esecuzione fino a \num{251} volte, e una
    variante che conserva più regole di un'altra risulta circa nove volte più lenta. La
    spiegazione candidata --- la perdita della scorciatoia esatta per le tuple con una sola cella
    nulla --- è dichiarata come non verificata.
\end{enumerate}

\unnumberedsection{Struttura dell'elaborato}

Il \autoref{chap:background} introduce i concetti preliminari e colloca il lavoro rispetto alla
letteratura: qualità dei dati, valori mancanti, tecniche di imputazione, \emph{Relaxed Functional
Dependencies}, riduzione di un insieme di regole, misure di similarità e la nozione di semantic
pruning adottata.

Il \autoref{chap:renuver} analizza il sistema \renuver{} per disassemblaggio: architettura, flusso
di imputazione, dataset impiegati, il doppio ruolo delle RFD, le due procedure di riempimento e la
caratterizzazione del costo computazionale.

Il \autoref{chap:methodology} presenta la metodologia: motivazioni, spazio delle strategie, le
strategie implementate, le condizioni esatte di potatura derivate dal comportamento del sistema e
il contratto di integrazione.

Il \autoref{chap:theory} formula e dimostra il risultato negativo che delimita il campo di
applicabilità, ne trae la conseguenza sul target di riduzione e ne precisa la portata.

Il \autoref{chap:hypotheses} espone come sono nate le ipotesi, le sette ipotesi primarie con il
loro criterio di falsificazione e le domande esplorative.

Il \autoref{chap:experiments} descrive la valutazione sperimentale: il protocollo di verifica e
la metodologia statistica, l'infrastruttura, i dataset e le configurazioni, le metriche, il
termine di paragone banale, gli esperimenti condotti con il collegamento alle ipotesi e lo studio
di ablazione.

Il \autoref{chap:results} analizza i risultati: la potatura che non perde nulla, il confronto
con la riduzione di ridondanza della teoria classica, la potatura che produce il guadagno e il
confine in cui il guadagno non si realizza, il meccanismo che li spiega, il criterio diagnostico,
i risultati negativi e il costo.

Il \autoref{chap:discussion} discute gli esiti delle ipotesi, il quadro complessivo, i limiti
dell'evidenza e la collocazione del lavoro rispetto alla letteratura.

Il \autoref{chap:defects} raccoglie i difetti e le criticità: quelli del sistema studiato, quelli
del codice sviluppato in questo lavoro e le conseguenze generali che se ne traggono.

Il \autoref{chap:conclusions} sintetizza il lavoro e ne indica gli sviluppi possibili. L'\autoref{app:reproducibility} documenta la riproducibilità dei risultati e
l'\autoref{app:results} raccoglie le tabelle complete delle campagne.


    \newpage\pagestyle{main}

    \chapter{Background e stato dell'arte}\label{chap:background}


\noindent Questo capitolo introduce i concetti su cui poggia il lavoro e li colloca rispetto alla
letteratura. Il problema affrontato nasce da una proprietà del dato --- l'incompletezza --- e dalla
possibilità di rimediarvi sfruttando regolarità osservate; il capitolo ricostruisce perciò il
quadro della qualità dei dati e delle sue dimensioni, la natura dei valori mancanti e i meccanismi
che li producono, le famiglie di tecniche di imputazione con le ipotesi che ciascuna richiede, le
dipendenze funzionali rilassate e le misure di similarità che ne definiscono la semantica, la
riduzione di un insieme di regole nella sua accezione classica, e infine la nozione di
\emph{semantic pruning} adottata in questo lavoro, con la sua collocazione rispetto alle linee di
ricerca affini.

% =======================================================================================
\section{Qualità dei dati e sue dimensioni}\label{sec:dq}

La qualità dei dati non è una proprietà intrinseca del dato ma una relazione fra il dato e l'uso
che ne viene fatto: un dataset è di qualità sufficiente rispetto a un compito, e la stessa
raccolta può risultare adeguata per un'analisi descrittiva e inadeguata per una decisione
operativa. La letteratura organizza questa relazione in \emph{dimensioni}, e la classificazione
più diffusa, derivata dallo studio delle aspettative degli utilizzatori, comprende accuratezza,
completezza, consistenza, tempestività e unicità \cite{wang1996beyond}. Le dimensioni non sono
indipendenti --- un dataset completo ma inconsistente è di uso difficile quanto uno accurato ma
incompleto --- e la loro rilevanza dipende dal compito.

\begin{table}[htbp]
\centering
\caption{Dimensioni della qualità dei dati rilevanti per il lavoro. Per ciascuna sono indicati il
difetto corrispondente e il rapporto con il problema affrontato.}
\label{tab:dq-dimensions}
\footnotesize
\begin{tabular}{@{}lp{3.6cm}p{4.6cm}@{}}
\toprule
dimensione & difetto corrispondente & rapporto con questo lavoro \\
\midrule
completezza & valori assenti in celle che dovrebbero essere valorizzate & è la dimensione di
partenza: l'imputazione esiste per rimediare alla sua violazione \\
accuratezza & valore presente ma diverso da quello corretto & è la dimensione che l'imputazione
può migliorare o peggiorare, ed è misurata dal ground truth delle configurazioni remodulate \\
consistenza & violazione di un vincolo fra attributi & è la proprietà che le RFD esprimono in
forma rilassata, e che il sistema usa per rifiutare le imputazioni incoerenti \\
ridondanza & informazione duplicata o non necessaria & è la dimensione su cui agisce la potatura:
un insieme di regole ridondante è un costo senza contropartita informativa \\
\bottomrule
\end{tabular}
\end{table}

Formalmente, la completezza di un dataset $D$ con schema $\mathcal{A} = \{A_1, \dots, A_m\}$ si
misura come la frazione di celle valorizzate sull'insieme delle celle attese,
$\mathrm{comp}(D) = 1 - |\{(t, A_i) : t[A_i] = \bot\}| \,/\, (|D| \cdot m)$; la sua violazione è
un \emph{valore mancante}, che è l'oggetto della sezione seguente. Nel lavoro la completezza è la
dimensione di partenza --- il sistema studiato esiste per rimediare alla sua violazione --- mentre
l'accuratezza è la dimensione lungo la quale si misura se il rimedio migliora o peggiora il dato.
La ridondanza dei metadati, che non compare nelle classificazioni classiche perché riguarda il
processo di arricchimento e non il dato, è il bersaglio della potatura ed è trattata nella
\autoref{sec:rule-reduction}.

Il quadro complessivo è quello dei sistemi di data quality, in cui l'imputazione è una delle
attività di miglioramento e va valutata insieme ai costi che introduce \cite{batini2016data}. È
questa duplice natura --- miglioramento del dato e costo di elaborazione --- a rendere sensato
chiedersi se una parte dell'informazione impiegata sia superflua.

% =======================================================================================
\section{Valori mancanti}\label{sec:missing}

Un valore mancante non è un valore: è l'assenza di un'informazione. La distinzione classica
riguarda il \emph{meccanismo} che lo ha prodotto, poiché da esso dipende la legittimità di
qualunque tecnica di ricostruzione \cite{rubin1976inference}.

\begin{definition}[Meccanismi di mancanza]
Sia $X$ una variabile e $M$ l'indicatore della sua mancanza. Si distinguono:
\begin{itemize}[leftmargin=1.4em]
    \item \textbf{MCAR} (\emph{Missing Completely At Random}): $M$ è indipendente da $X$ e da
    ogni altra variabile osservata;
    \item \textbf{MAR} (\emph{Missing At Random}): $M$ può dipendere dalle variabili osservate,
    ma non dal valore mancante $X$;
    \item \textbf{MNAR} (\emph{Missing Not At Random}): $M$ dipende dal valore mancante stesso.
\end{itemize}
\end{definition}

La distinzione ha una conseguenza pratica diretta: solo sotto MCAR e MAR la distribuzione dei
dati osservati è informativa rispetto al valore mancante, e solo in quei regimi un metodo di
imputazione basato sulle correlazioni osservate può essere corretto in senso statistico
\cite{little2019statistical}. Sotto MNAR la ricostruzione richiede di modellare il meccanismo di
mancanza, e un metodo che non lo faccia produce stime distorte in modo non controllabile dai dati
disponibili.

Nel contesto di questo lavoro la distinzione ha un ruolo operativo preciso. I dataset impiegati
sono \emph{remodulati}: i valori mancanti sono introdotti artificialmente a partire da un dataset
completo, e il valore originale è conservato come ground truth. Il meccanismo è quindi, per
costruzione, noto e controllato. Va però osservato che la rimozione artificiale è progettata per
essere casuale, e ricade pertanto nel caso MCAR: i risultati del lavoro valgono per questo regime,
e non si estendono automaticamente a mancanze strutturali, in cui l'assenza di un valore è essa
stessa un'informazione --- come accade quando un attributo non è applicabile a una parte delle
tuple.

% =======================================================================================
\section{Data imputation}\label{sec:imputation}

L'imputazione consiste nel sostituire un valore mancante con una stima ottenuta dal resto
dell'informazione disponibile. Le famiglie di tecniche si distinguono per la natura
dell'informazione sfruttata e per le ipotesi che devono valere perché la stima sia sensata.

\begin{table}[htbp]
\centering
\caption{Famiglie di tecniche di imputazione. Per ciascuna sono indicati l'informazione
impiegata, l'ipotesi richiesta e il comportamento rispetto alla copertura.}
\label{tab:imputation-families}
\footnotesize
\begin{tabular}{@{}lp{2.6cm}p{3.2cm}p{2.5cm}@{}}
\toprule
famiglia & informazione impiegata & ipotesi & copertura \\
\midrule
statistica & distribuzione marginale o congiunta di una variabile & il modello globale è corretto;
sotto MAR, stima non distorta & totale: un valore viene sempre prodotto \\
apprendimento automatico & relazioni fra attributi apprese dai dati osservati & i dati mancanti
seguono il medesimo processo dei dati osservati & totale \\
similarità & somiglianza fra tuple (\emph{k}-NN) & tuple vicine hanno valori simili & totale o
quasi, secondo la soglia adottata \\
vincoli & dipendenze fra attributi, espresse come regole & il vincolo è valido sul dato reale &
\emph{parziale}: l'imputazione può essere rifiutata \\
\bottomrule
\end{tabular}
\end{table}

I metodi statistici classici --- media, mediana, moda, regressione --- sono semplici e veloci e
assumono una struttura globale che raramente sussiste; l'imputazione multipla per equazioni
concatenate ne rappresenta la versione moderna e condizionata, in cui ogni variabile è stimata a
partire dalle altre in modo iterativo \cite{vanbuuren2011mice}. I metodi di apprendimento
automatico, come le foreste casuali applicate all'imputazione di dati misti, rinunciano a
un'assunzione distributiva e catturano relazioni non lineari, al prezzo di un costo di
addestramento e di una minore interpretabilità \cite{stekhoven2012missforest}. I metodi basati su
similarità copiano il valore da tuple ritenute vicine, con la scelta della metrica e del numero di
vicini come parametri critici \cite{troyanskaya2001missing}. I metodi basati su vincoli
impiegano dipendenze fra attributi --- funzionali, di inclusione, di ordine --- e possono essere
formulati come un problema di riparazione che massimizza la coerenza complessiva
\cite{rekatsinas2017holoclean, fan2012foundations}; è la famiglia a cui appartiene \renuver{}.

La scelta della famiglia non è neutrale rispetto alla valutazione. Un metodo statistico produce
sempre un valore: la sua copertura è per definizione totale, e la sua qualità si misura
sull'accuratezza. Un metodo basato su vincoli può invece \emph{rifiutare} di imputare, quando
nessun candidato soddisfa i vincoli: in quel caso la copertura è una metrica tanto rilevante
quanto l'accuratezza, e le due vanno lette insieme. Questa osservazione è alla base delle metriche
adottate nel \autoref{chap:experiments}, ed è anche la ragione per cui una modifica dell'insieme
di vincoli --- come la potatura studiata in questo lavoro --- può cambiare il risultato in due
direzioni opposte, migliorando la qualità delle imputazioni prodotte e riducendo il numero di
celle che ricevono un valore.

% =======================================================================================
\section{Relaxed Functional Dependencies}\label{sec:rfd}

Una dipendenza funzionale (FD) classica asserisce che due tuple concordanti su un insieme di
attributi $X$ concordano necessariamente su un attributo $A$. Su dati reali l'asserzione è troppo
rigida: il rumore e le eccezioni la falsificano quasi ovunque. Le dipendenze funzionali
\emph{rilassate} introducono forme di tolleranza, tipicamente sulle soglie di confronto e sulle
eccezioni ammesse \cite{caruccio2016rfd}.

\begin{definition}[Relaxed Functional Dependency]\label{def:rfd}
Una RFD è un'espressione della forma
\[
    (A_1 \le \theta_1,\ \dots,\ A_k \le \theta_k) \;\rightarrow\; (A_{\RHS} \le \theta_{\RHS})
\]
dove gli $A_i$ sono attributi del dataset, $\theta_i \ge 0$ sono soglie di tolleranza e
$A_{\RHS}$ è l'attributo di destra. La parte sinistra è detta \LHS{}, la destra \RHS{}.
\end{definition}

La semantica della tolleranza dipende dal tipo dell'attributo: per un attributo numerico
$A_i \le \theta_i$ significa che due tuple possono differire di al più $\theta_i$ su $A_i$; per un
attributo categorico significa che i valori devono coincidere, oppure che la loro distanza di
stringa non supera $\theta_i$. La distinzione è rilevante nel seguito, poiché il sistema studiato
tratta i due casi con procedure diverse.

Due osservazioni sono necessarie per la lettura del lavoro.

\begin{observation}[Ruolo doppio]\label{obs:dual}
Una RFD non è soltanto un descrittore del dataset: è un \emph{oggetto operativo}. Nel sistema
studiato essa partecipa a due procedure distinte, la generazione di candidati per una cella nulla
e la verifica di coerenza di un'imputazione. Le due procedure hanno requisiti opposti rispetto
alla tolleranza: una soglia larga rende la regola permissiva nella generazione e severa nella
verifica. Questo doppio ruolo è il vincolo che determina l'intero problema affrontato
(\autoref{sec:dual-role}).
\end{observation}

\begin{observation}[Soglia nulla]
Il caso $\theta_{\RHS} = 0$ merita attenzione separata. Una RFD con soglia di destra nulla
richiede coincidenza esatta sull'attributo di destra, ed è quella che il sistema tratta con una
scorciatoia dedicata. Nel seguito la distinzione fra regole a soglia nulla e regole a soglia
positiva ricorrerà come variabile esplicativa.
\end{observation}

% =======================================================================================
\section{Ridondanza e riduzione di un insieme di regole}\label{sec:rule-reduction}

L'insieme delle RFD estratte da un dataset è, come ogni insieme di dipendenze, ridondante: alcune
regole sono implicati da altre, e la loro presenza non aggiunge informazione sui vincoli del dato.
La teoria classica delle dipendenze funzionali affronta il problema con la nozione di
\emph{copertura}: due insiemi di dipendenze sono equivalenti se hanno la stessa chiusura, e un
insieme è \emph{minimale} se nessuna delle sue dipendenze è derivabile dalle altre per mezzo degli
assiomi di Armstrong \cite{armstrong1974dependency, maier1983theory}. Il calcolo di una copertura
minimale è un'operazione standard, e la sua applicazione alle RFD è stata studiata come forma di
semplificazione degli insiemi di regole.

Questa letteratura fornisce il termine di paragone naturale del lavoro, ed è importante
chiarire perché non lo risolve. La nozione classica di ridondanza è \emph{dichiarativa}: una
dipendenza è superflua se la sua rimozione non cambia l'insieme dei vincoli soddisfatti. Nel
sistema studiato il ruolo di una regola è invece \emph{operativo}, e dipende da come il programma
la consuma: una regola partecipa alla generazione di un valore e alla verifica di coerenza, e la
scorciatoia per le regole a soglia di destra nulla mostra che due regole equivalenti dal punto di
vista logico possono avere costi ed effetti diversi. Ne segue che la ridondanza classica non
implica la sicurezza della rimozione, e che il problema studiato non è un caso particolare del
calcolo di una copertura minimale.

Il \autoref{chap:theory} formalizza questa distinzione. La conseguenza per la valutazione
sperimentale è che il termine di paragone del lavoro non è una copertura minimale --- che non
garantirebbe l'equivalenza del comportamento --- ma l'\emph{equivalenza osservata} del sistema
prima e dopo la potatura, misurata per confronto degli output. La scelta è più onerosa e più
debole sul piano teorico, ma è l'unica che risponda alla domanda posta: se una regola possa essere
rimossa senza che il sistema se ne accorga.

% =======================================================================================
\section{Misure di similarità e distanza}\label{sec:similarity}

Il confronto fra tuple avviene attributo per attributo, e il risultato complessivo è una
aggregazione dei confronti elementari. Nel sistema studiato l'aggregazione è una media dei
contributi normalizzati sull'ampiezza del lato sinistro della regola. La normalizzazione ha
conseguenze non ovvie: a parità di scostamento assoluto, un contributo è tanto minore quanto più
grande è la soglia di tolleranza dell'attributo, sicché attributi con scale diverse non sono
confrontabili direttamente.

Per gli attributi numerici il confronto elementare è la differenza assoluta, eventualmente
normalizzata sull'intervallo dei valori osservati. Per gli attributi testuali si impiegano misure
di edit distance, fra cui la distanza di Levenshtein, che conta il numero minimo di inserimenti,
cancellazioni e sostituzioni necessari a trasformare una stringa nell'altra
\cite{levenshtein1966binary}; la sua versione normalizzata sulla lunghezza della stringa è
impiegata per rendere confrontabili valori di lunghezza diversa, ed è quella adottata per la
famiglia Physician, i cui attributi testuali designano la stessa entità con forme diverse. Anche
in questo caso la soglia di tolleranza della regola agisce da parametro di accettazione.

La scelta della misura non è neutrale, come mostra il difetto D1 discusso nella
\autoref{sec:d1}: nel sistema studiato l'attributo categorico non riceve la misura prevista, e il
confronto ricade su una distanza di stringa applicata a etichette lunghe e prive di relazione. La
circostanza è rilevante per l'interpretazione dei risultati, perché rende il comportamento del
sistema dipendente da una misura che non è quella dichiarata.

% =======================================================================================
\section{Semantic pruning: definizione operativa e collocazione}\label{sec:semantic-pruning}

Con \emph{semantic pruning} si intende, in questo lavoro, la selezione di un sottoinsieme di un
insieme di metadati --- nel caso in esame, le RFD --- guidata dal \emph{significato} che quei
metadati hanno per il sistema che li consuma, e non da criteri sintattici o statistici sul file.

La distinzione è sostanziale. Un criterio sintattico rimuove regole in base alla loro forma
(duplicati, dimensione del lato sinistro, valore delle soglie); un criterio statistico le rimuove
in base alla loro frequenza o alla loro resa osservata. Un criterio semantico richiede invece di
sapere \emph{che cosa il sistema fa} con ciascuna regola, e quindi di derivare le condizioni di
rimozione dal comportamento effettivo del codice. La qualificazione ``semantico'' non è quindi un
sinonimo di ``euristico'': indica la fonte dell'informazione impiegata, che è la semantica
operativa del consumatore dei metadati e non una proprietà statistica del file.

\begin{definition}[Potatura sicura]\label{def:safe-pruning}
Sia $\mathcal{R}$ un insieme di regole e $\Pi$ il programma che le consuma. Una rimozione
$r \in \mathcal{R}$ è \emph{sicura rispetto a una proprietà $P$} se il comportamento di $\Pi$
rispetto a $P$ è invariato quando $r$ è assente.
\end{definition}

La definizione è parametrica rispetto a $P$, e la scelta di $P$ è il nodo del problema. Nel
sistema studiato si possono isolare almeno due proprietà candidate --- la capacità di generare
candidati e il potere di veto --- e si dimostrerà che esse non sono simultaneamente preservabili
oltre una soglia quantificata (\autoref{sec:negative-theorem}).

Questa impostazione distingue il lavoro da tre alternative presenti in letteratura. La prima è la
potatura \emph{a posteriori}, che misura l'effetto di ogni rimozione eseguendo il sistema:
corretta ma computazionalmente proibitiva, poiché richiede una esecuzione per candidato. La
seconda è la potatura basata su un \emph{proxy} offline, che stima la rilevanza di una regola con
una grandezza calcolata sul dataset: economica, ma --- come si mostrerà nel \autoref{chap:results}
--- potenzialmente fuorviante, al punto da risultare \emph{anti-correlata} con la rilevanza
effettiva. La terza è la selezione guidata dalla semantica \emph{dei nomi} dei metadati, che
impiega modelli linguistici per inferire relazioni fra colonne a partire dalla loro
denominazione \cite{trummer2024llmcorrelations} o per guidare l'esplorazione di fonti dati
eterogenee \cite{freire2024llmdiscovery}: si tratta di una linea di ricerca affine per ambizione
--- impiegare il significato invece della forma --- ma distinta per oggetto, perché riguarda
l'\emph{interpretazione} dei metadati e non la loro \emph{selezione} rispetto al consumo da parte
di un algoritmo noto.

La metodologia adottata cerca una quarta via, intermedia fra le precedenti: derivare dal codice
condizioni \emph{sufficienti}, verificabili offline, la cui validità non dipende da una stima. È
questa la scelta che rende il risultato negativo del \autoref{chap:theory} possibile e
significativo: poiché le condizioni non sono euristiche, il loro limite è un limite del problema e
non dell'implementazione.


    \chapter{Analisi del sistema \renuver{}}\label{chap:renuver}


\noindent Questo capitolo analizza il sistema oggetto di studio. L'analisi non si basa sulla
documentazione ma sul \emph{comportamento effettivo} del codice: il sistema è distribuito come
archivio eseguibile senza sorgenti, e le proprietà su cui poggia l'intera metodologia sono state
derivate per disassemblaggio. Ne emergono l'architettura e il flusso di imputazione, i ruoli che
le regole vi svolgono --- generazione dei candidati, verifica di coerenza e rivalutazione delle
regole escluse dalla ricerca ordinaria --- le due procedure con cui il sistema tratta le tuple
incomplete, e la distribuzione del costo computazionale fra le sue fasi. Da questa analisi
discendono le condizioni esatte di potatura su cui si fonda la metodologia, che sono esposte nel
capitolo successivo insieme alle strategie che le impiegano.

% =======================================================================================
\section{Obiettivi del sistema}\label{sec:renuver-goals}

\renuver{} \cite{breve2022renuver} è un algoritmo di imputazione di valori mancanti basato su
\emph{Relaxed Functional Dependencies}. L'idea portante è che le regole di imputazione non
debbano essere fornite dall'esterno, ma estratte dal dataset stesso: il sistema riceve in
ingresso un insieme di RFD già estratte e le impiega per ricostruire le celle nulle.

Rispetto ai metodi statistici, l'approccio ha il vantaggio di produrre imputazioni motivate da
vincoli osservati nei dati, e di poter \emph{rifiutare} di imputare quando nessun candidato
soddisfa i vincoli. Quest'ultima proprietà è deliberata: è preferibile una cella non imputata a
una cella imputata in violazione delle dipendenze del dataset.

% =======================================================================================
\section{Architettura generale}\label{sec:architecture}

\begin{figure}[htbp]
\centering
\begin{tikzpicture}[
    font=\footnotesize,
    box/.style={draw, rounded corners=1pt, align=center, inner sep=4pt, minimum height=8mm},
    file/.style={draw, align=center, inner sep=4pt, minimum height=8mm, fill=black!6},
    ar/.style={-{Latex[length=2mm]}, thick}
]
\node[file] (ds) {dataset\\$D$};
\node[box, right=8mm of ds] (est) {estrazione\\delle RFD};
\node[file, right=8mm of est] (rfd) {file delle\\regole $R$};
\node[box, right=8mm of rfd] (sys) {sistema\\\renuver{} };
\node[file, right=8mm of sys] (out) {celle\\popolate};
\node[box, below=10mm of rfd, fill=black!14] (prune) {potatura esterna\\(riscrittura del file)};
\draw[ar] (ds) -- (est);
\draw[ar] (est) -- (rfd);
\draw[ar] (rfd) -- (sys);
\draw[ar] (sys) -- (out);
\draw[ar] (ds.south) |- ([yshift=-4mm]sys.south) -- (sys.south);
\draw[ar] (rfd.south) -- (prune.north);
\draw[ar] (prune.east) -| ([xshift=6mm]sys.south);
\end{tikzpicture}
\caption{Flusso di elaborazione del sistema studiato e punto di innesto della potatura. La
potatura riscrive il file delle regole prima che il sistema lo consumi: il dataset e l'eseguibile
restano invariati.}
\label{fig:pipeline}
\end{figure}

Il sistema è distribuito come archivio Java eseguibile. L'analisi ha stabilito i seguenti fatti,
rilevanti per il seguito.

\begin{itemize}[leftmargin=1.4em]
    \item Il codice applicativo consta di \num{13} classi nel \emph{package} di default, incluse
    in un archivio ``grasso'' che incorpora anche la libreria di collezioni \texttt{fastutil}
    (\num{11040} classi complessive). Le classi applicative sono compilate per la versione
    \num{52} del formato \texttt{class} (Java~8).
    \item L'interfaccia a riga di comando accetta sei argomenti: la tipizzazione delle colonne,
    il file del dataset, il separatore, il file delle RFD, una soglia numerica e un selettore
    di euristica.
    \item \textbf{Il sesto argomento è ignorato.} Il metodo di avvio non lo legge: il sistema
    esegue sempre la medesima procedura di riempimento. Le altre procedure presenti nel codice
    --- ricerca locale, riordino per numero di valori nulli, riordino per coinvolgimento nel
    lato sinistro --- non sono raggiungibili dalla riga di comando. La verifica è stata condotta
    per disassemblaggio e confermata empiricamente. Il \texttt{README} dell'archivio documenta
    peraltro i parametri come \texttt{args0}--\texttt{args4} e \texttt{args6} --- l'indice
    \texttt{args5} non compare --- e precisa che ``in questa versione è supportata soltanto
    l'euristica sequenziale \texttt{0}'': la documentazione dichiara quindi un limite di
    \emph{versione}, mentre il codice mostra che l'argomento non viene mai letto. Le due
    circostanze sono distinte e la seconda è più forte della prima, perché esclude che il
    parametro possa essere onorato da una versione futura senza modifiche al codice.
    \item La quinta soglia ha un \emph{doppio ruolo}: filtra le RFD in ingresso e funge da
    valore iniziale della soglia di accettazione del punteggio di similarità. I due ruoli non
    sono separabili dalla riga di comando.
\end{itemize}

% =======================================================================================
\section{Flusso di imputazione e procedure di riempimento}\label{sec:workflow}\label{sec:two-branches}

Il flusso si articola in tre fasi.

\begin{enumerate}[leftmargin=1.6em]
    \item \textbf{Costruzione degli indici.} Le RFD vengono raggruppate in \emph{cluster},
    dove un cluster è l'insieme delle regole che condividono lo stesso attributo di destra e la
    stessa soglia di destra. Contestualmente si individuano le tuple del dataset che presentano
    celle nulle.
    \item \textbf{Trattamento delle tuple incomplete.} Per ogni tupla con celle nulle si tenta
    di riempirle, secondo due procedure distinte a seconda che la tupla abbia una sola cella
    nulla o più di una (\autoref{sec:two-branches}).
    \item \textbf{Verifica e registrazione.} Ogni imputazione candidata è sottoposta a una
    verifica di coerenza rispetto all'insieme delle regole; solo se supera la verifica viene
    registrata.
\end{enumerate}

% =======================================================================================
\section{Dataset impiegati}\label{sec:datasets}

L'archivio comprende cinque famiglie di dataset, con caratteristiche molto diverse fra loro. La
\autoref{tab:families} ne riassume la struttura rilevante per il pruning.

\begin{table}[htbp]
\centering
\caption{Caratteristiche delle famiglie di dataset. $N$ è il numero di RFD alla soglia indicata;
``celle valutate'' è il numero di celle per cui è disponibile il valore di ground truth, che
--- come si dirà --- non coincide con il numero di celle nulle del dataset.}
\label{tab:families}
\footnotesize
\begin{tabular}{lrrrr}
\toprule
famiglia & dataset & $N$ (soglia) & celle valutate & esito dell'analisi \\
\midrule
Bridges    & 25 & \num{1830}--\num{18844} & 14--70   & utilizzabile \\
Glass      & 25 & \num{1511}--\num{27445} & 22--107  & richiede correzione di schema \\
Restaurant & 45 & 25--124               & fino a \num{2074} & utilizzabile con metrica testuale \\
Physician  & 29 & \num{3734}            & fino a \num{338507} & non utilizzabile \\
Nursery    &  0 & \num{60}--\num{1693} byte & ---      & regole presenti, dataset assenti \\
\bottomrule
\end{tabular}
\end{table}

Due circostanze vanno segnalate perché condizionano la valutazione sperimentale.

\paragraph{Le celle valutate non sono le celle nulle.} Il valore di ground truth è fornito da un
file separato, in formato \texttt{riga;attributo;valore}, che elenca soltanto le celle introdotte
artificialmente. Il dataset può contenere altre celle nulle, preesistenti, che il sistema
tratta ma che non vengono valutate. Di conseguenza copertura e accuratezza sono relative
all'insieme delle celle elencate, non a tutte le celle nulle del dataset. La distinzione è
necessaria per interpretare correttamente le metriche.

\paragraph{Il costo computazionale non è uniforme.} I dataset Physician non remodulati contengono
\num{103592} righe e fino a \num{121913} celle nulle: una singola esecuzione richiederebbe tempi
non compatibili con una campagna sperimentale. Le versioni \emph{remodulate} della stessa
famiglia, di dimensioni ridotte, presentano a loro volta difetti nel materiale di ground truth
(\autoref{chap:defects}).

% =======================================================================================
\section{Il doppio ruolo delle RFD}\label{sec:dual-role}

\begin{algorithm}[htbp]
\caption{Verifica dei veti su una tupla incompleta, ricostruita dal bytecode. Il ciclo esamina
\emph{tutte} le regole dell'insieme e, per ciascuna, scansiona il dataset alla ricerca di un
testimone che violi la soglia di destra.}
\label{alg:veto}
\begin{algorithmic}[1]
\State \textbf{input}: tupla $t$ con la cella $t[A]$ nulla, insieme di regole $R$, dataset $D$
\ForAll{$r \in R$ tale che $\RHS(r) = A$}
    \If{$\LHS(r) \cap \NullCols(D) = \emptyset$}
        \State \textbf{continue} \Comment{nessun confronto possibile: la regola non può vetare}
    \EndIf
    \State cerca in $D$ una tupla $s \neq t$ che concordi con $t$ su ogni colonna di $\LHS(r)$
           entro la soglia $\theta_r$
    \If{esiste $s$ e $t[A]$ differisce da $s[A]$ oltre $\theta_{\RHS}(r)$}
        \State \Return \textbf{incoerente} \Comment{veto esercitato da $r$}
    \EndIf
\EndFor
\State \Return \textbf{coerente}
\end{algorithmic}
\end{algorithm}

Questa sezione stabilisce il vincolo che determina l'intero problema. L'analisi del codice mostra
che una RFD partecipa a \emph{due} procedure con requisiti opposti.

\begin{description}[leftmargin=0em]
    \item[Ruolo generativo.] Data una cella nulla, il sistema cerca tuple del dataset
    sufficientemente simili alla tupla incompleta; il punteggio di similarità è calcolato
    rispetto alle regole di un cluster, e il miglior punteggio determina il valore da imputare.
    In questo ruolo una soglia \emph{larga} rende la regola più permissiva.

    \item[Ruolo di veto.] Prima di registrare un'imputazione, il sistema verifica che essa non
    violi alcuna regola. In questo ruolo una soglia \emph{stretta} rende la regola più severa.
\end{description}

Le due procedure non sono separabili nel formato del file: una RFD è una riga di testo, e nulla
indica se essa serva all'una, all'altra, o a entrambe. Ne segue l'osservazione che governa il
lavoro.

\begin{observation}[Il pruning esterno è intrinsecamente conservativo]\label{obs:conservative}
Poiché rimuovere una regola ne elimina simultaneamente il contributo generativo e quello di veto,
e poiché i due contributi non sono distinguibili dal formato, una strategia di potatura che operi
sul solo file delle regole deve assumere che ogni regola rimossa possa essere necessaria al veto.
\end{observation}

La \autoref{tab:veto-share} quantifica quanto questa assunzione sia stringente: la quasi totalità
delle regole è, in linea di principio, in grado di esercitare un veto.

\paragraph{Un terzo ruolo, esercitato da una parte delle regole.} All'avvio il sistema compie una
propria riduzione dell'insieme: le regole con soglia di destra \emph{esattamente nulla} il cui lato
sinistro si comporta come una chiave sul dataset vengono tolte dall'insieme di lavoro, raccolte in
una struttura separata e sottoposte a una procedura distinta, invocata quando la ricerca ordinaria
di un candidato non produce alcun valore accettabile. Il log dell'esecuzione registra l'operazione
(``\texttt{Removing key RFD\ldots}''), e il numero di regole rimosse non è marginale: su
\texttt{bridges\_14\_1} a soglia 9 l'insieme passa da \num{8961} a \num{8581} regole. La verifica
della proprietà di chiave è inoltre quadratica nella dimensione del dataset, ed è la ragione per
cui il preprocessing del sistema --- che non è la potatura studiata in questo lavoro, ma una fase
interna dell'eseguibile --- domina il tempo sulle configurazioni grandi: sulla medesima
configurazione \SI{24.5}{\second} di preprocessing contro \SI{7.1}{\second} di imputazione, e fino
a \SI{113}{\second} su una configurazione Physician.

La circostanza ha due conseguenze. La prima riguarda la descrizione del sistema: i ruoli
osservabili sono \emph{tre}, non due, e il terzo è esercitato soltanto da una parte delle regole.
La seconda riguarda la portata delle garanzie: la condizione \Vtwo{} e il teorema del
\autoref{chap:theory} sono enunciati rispetto all'insieme delle regole così come compare nel file,
mentre il sistema ne consuma un sottoinsieme determinato da un criterio proprio; la verifica
sperimentale di equivalenza (\autoref{sec:res-vetocover}) copre l'interazione fra i due
meccanismi, ma un trattamento formale che la includesse dovrebbe modellare anche quella
riduzione. La questione è ripresa nella \autoref{sec:theory-scope}.

\begin{table}[htbp]
\centering
\caption{Quota di RFD che possono esercitare un veto, per famiglia e soglia. Una RFD può
esercitare un veto solo se il suo lato sinistro coinvolge almeno una colonna che presenta celle
nulle nel dataset.}
\label{tab:veto-share}
\footnotesize
\begin{tabular}{llrrr}
\toprule
famiglia & soglia & $N$ & RFD che \emph{possono} fare veto & quota \\
\midrule
Bridges    & 3  & \num{1830}  & \num{1818} & \SI{99.3}{\percent} \\
Bridges    & 9  & \num{8961}  & \num{8945} & \SI{99.8}{\percent} \\
Bridges    & 15 & \num{18844} & \num{18828} & \SI{99.9}{\percent} \\
Glass      & 3  & \num{1511}  & \num{1476} & \SI{97.7}{\percent} \\
Glass      & 9  & \num{11648} & \num{11564} & \SI{99.3}{\percent} \\
Glass      & 15 & \num{27445} & \num{27327} & \SI{99.6}{\percent} \\
Restaurant & 3  & 25          & 23         & \SI{92.0}{\percent} \\
Restaurant & 6  & 124         & 116        & \SI{93.5}{\percent} \\
\bottomrule
\end{tabular}
\end{table}

% =======================================================================================

\paragraph{Le due procedure di riempimento.} Il trattamento di una tupla incompleta segue due
procedure diverse, la cui distinzione è essenziale per interpretare i risultati.

\paragraph{Tupla con una sola cella nulla.} Il sistema scorre i cluster della colonna interessata
e si arresta al primo che produce un candidato valido. Esiste inoltre una scorciatoia: se la
regola impiegata ha soglia di destra nulla, il valore della tupla candidata viene scritto
direttamente e la procedura termina. Ne segue che, in questo ramo, la \emph{posizione} di un
cluster nell'ordine di scansione è determinante.

\paragraph{Tupla con più celle nulle.} La procedura è iterativa: per un numero di round pari al
massimo numero di cluster fra le colonne coinvolte, e per ciascun round, ogni cella nulla riceve
i candidati del cluster corrispondente al round corrente. La scrittura del valore avviene senza
verificare se la cella sia già stata riempita in un round precedente; un round successivo
sovrascrive quindi il risultato del precedente, e un round i cui candidati fallissero tutti la
verifica di coerenza può riportare la cella allo stato nullo. In questo ramo, dunque, il valore
finale di una cella dipende dall'\emph{ultimo} cluster produttivo della sua colonna.

La distinzione ha una conseguenza quantitativa che sarà ripresa nel \autoref{chap:results}: le
due procedure non sono ugualmente rappresentate nei dataset. Su Bridges circa tre quarti delle
celle nulle appartengono a tuple con più di una cella nulla, mentre su Glass e Restaurant la
grande maggioranza appartiene a tuple con una sola cella nulla.

% =======================================================================================
\section{Costo computazionale}\label{sec:cost}

\begin{figure}[htbp]
\centering
\begin{tikzpicture}
\begin{axis}[
    ybar,
    bar width=11pt,
    width=0.78\textwidth,
    height=0.40\textwidth,
    ymin=0, ymax=104,
    ylabel={quota del tempo attribuito alle due procedure (\si{\percent})},
    xtick={1,2},
    xticklabels={\texttt{bridges\_70\_1} @15, \texttt{bridges\_14\_1} @9},
    legend style={font=\footnotesize, at={(0.5,-0.22)}, anchor=north, legend columns=2, draw=none, fill=none},
    nodes near coords,
    every node near coord/.append style={font=\scriptsize},
    grid=major,
    grid style={black!10},
]
\addplot coordinates {(1,0.8) (2,14.4)};
\addplot coordinates {(1,99.2) (2,85.6)};
\legend{generazione dei candidati, verifica dei veti}
\end{axis}
\end{tikzpicture}
\caption{Ripartizione del tempo di esecuzione fra le due procedure, misurata per strumentazione su
due configurazioni. La verifica dei veti domina in entrambe, e all'interno di essa il costo è
quasi interamente nella scansione del dataset.}
\label{fig:cost}
\end{figure}

Poiché il sistema è distribuito senza sorgenti, la determinazione di dove si concentri il tempo
di esecuzione richiede di ricompilarlo e strumentarlo. La procedura --- ricompilazione da sorgenti
decompilati, con verifica di equivalenza rispetto all'eseguibile originale --- è descritta
nella \autoref{sec:instrumentation}; qui se ne riportano i risultati, che orientano la scelta delle
strategie di ottimizzazione.

\begin{table}[htbp]
\centering
\caption{Ripartizione del tempo di esecuzione fra le due procedure principali, misurata per
strumentazione. La percentuale è calcolata sul tempo complessivamente attribuito alle due
procedure.}
\label{tab:cost}
\footnotesize
\begin{tabular}{lrrrr}
\toprule
configurazione & procedura & chiamate & tempo & quota \\
\midrule
\multirow{2}{*}{\texttt{bridges\_70\_1}, soglia 15}
  & generazione candidati & 531      & \SI{1.4}{\second}   & \SI{0.8}{\percent} \\
  & verifica dei veti     & 151      & \SI{158.4}{\second} & \SI{99.2}{\percent} \\
\midrule
\multirow{2}{*}{\texttt{bridges\_14\_1}, soglia 9}
  & generazione candidati & 293      & \SI{1.5}{\second}   & \SI{14.4}{\percent} \\
  & verifica dei veti     & 8        & \SI{9.1}{\second}   & \SI{85.6}{\percent} \\
\bottomrule
\end{tabular}
\end{table}

All'interno della verifica dei veti, la quasi totalità del tempo è assorbita dalla scansione del
dataset: su \texttt{bridges\_70\_1} a soglia 15, \num{354809} invocazioni della verifica di
violazione assorbono \SI{158.3}{\second} su \SI{158.4}{\second} totali, mentre il confronto dei
nomi di colonna --- che pure è il primo test eseguito --- costa \SI{0.16}{\second}.

Il risultato ha una conseguenza diretta e controintuitiva: \textbf{il tempo è nel veto, non nella
generazione}. Ne segue che le strategie di potatura producono un guadagno temporale proprio in
quanto riducono il numero di regole esaminate dalla verifica dei veti, e che qualunque
ottimizzazione mirata alla sola generazione dei candidati non potrebbe che incidere su una
frazione marginale del costo.



    \chapter{Metodologia di semantic pruning}\label{chap:methodology}


\noindent Questo capitolo presenta la metodologia sviluppata. Le condizioni esatte di potatura,
derivate dal comportamento del sistema, costituiscono il criterio con cui le strategie vanno
lette: da esse discende che cosa si può rimuovere con garanzia e che cosa soltanto a costo di un
rischio dichiarato. Il capitolo motiva l'approccio, organizza lo spazio delle strategie lungo le
dimensioni che ne determinano la natura, descrive quelle implementate con il loro regime di
garanzia, la loro riduzione e il loro costo, spiega perché la generazione dei candidati non
ammetta una condizione esatta della stessa natura di quella che governa il veto, e definisce il
contratto con cui la potatura si integra nel sistema. Le proprietà formali e il risultato negativo
che delimita il campo di applicabilità sono oggetto del capitolo successivo.

% =======================================================================================
\section{Motivazioni}\label{sec:why-pruning}

Le ragioni per potare l'insieme delle RFD sono tre, di natura diversa, e la loro compresenza
spiega perché il problema non si riduca a un'ottimizzazione di prestazioni.

\paragraph{Ragione economica.} Il costo dell'imputazione cresce con il numero di regole: la
verifica dei veti esamina l'intero insieme per ogni imputazione candidata, e la
\autoref{sec:cost} mostra che è questa la voce dominante del tempo di esecuzione. Ridurre
l'insieme riduce quel costo, e la riduzione è tanto più preziosa quanto più il sistema è
impiegato su dataset di dimensioni maggiori.

\paragraph{Ragione di qualità.} Una regola non è necessariamente benefica. Regole con soglie
larghe competono nella generazione dei candidati e possono prevalere su regole più precise,
producendo imputazioni peggiori di quelle che si otterrebbero senza di esse. La potatura può
quindi migliorare il risultato, non soltanto accelerarlo: è la circostanza che il
\autoref{chap:results} misura, e che rende il problema interessante dal punto di vista della
qualità dei dati e non solo dell'efficienza.

\paragraph{Ragione epistemica.} Un insieme di migliaia di regole è esso stesso un oggetto
difficile da interpretare: non è ispezionabile, non è confrontabile fra dataset e non consente di
attribuire un comportamento a una causa. La selezione di un sottoinsieme rilevante ha valore
conoscitivo indipendentemente dal guadagno prestazionale, perché è la condizione per cui il
comportamento del sistema diventa spiegabile.

\paragraph{La tensione che governa il problema.} Le tre ragioni convergono su un'unica
operazione --- rimuovere regole --- ma il sistema studiato assegna a ogni regola due ruoli
distinti, la generazione dei candidati e il veto sulle imputazioni incoerenti, governati da
requisiti opposti sulla stessa grandezza: una soglia stretta rende la regola più selettiva come
generatore e più severa come veto, mentre una soglia larga la rende più permissiva come generatore
e più debole come veto. Nulla, nel formato del file, distingue i due ruoli
(\autoref{sec:dual-role}). Ne segue che ogni strategia di rimozione deve dichiarare
\emph{quale} dei due ruoli è disposta a sacrificare, e che la nozione di ``potatura sicura''
(\autoref{def:safe-pruning}) richiede di specificare rispetto a che cosa la sicurezza è pretesa.
È questa precisazione a rendere possibile il risultato negativo del \autoref{chap:theory}.

% =======================================================================================
\section{Dal sistema alle condizioni esatte}\label{sec:exact-conditions}

\begin{algorithm}[htbp]
\caption{Ricerca dei copritori per la condizione \Vtwo{}. Il parametro $\mu$ limita la differenza
di ampiezza fra il lato sinistro della regola coprente e quello della regola coperta; la
\autoref{sec:maxextra} mostra che il risultato non dipende dal suo valore.}
\label{alg:cover}
\begin{algorithmic}[1]
\State \textbf{input}: insieme di regole $R$, dataset $D$, parametro $\mu$
\State $C \gets \emptyset$ \Comment{regole la cui rimozione è garantita}
\ForAll{$r \in R$ con $\LHS(r) \cap \NullCols(D) \neq \emptyset$}
    \ForAll{$r'' \in R \setminus \{r\}$ con $\RHS(r'') = \RHS(r)$ e
            $|\LHS(r)| - |\LHS(r'')| \le \mu$}
        \If{$\LHS(r'') \subseteq \LHS(r)$ \textbf{and} $\theta_{r''}(c) \ge \theta_r(c)\ \forall c \in \LHS(r'')$
            \textbf{and} $\theta_{\RHS}(r'') \le \theta_{\RHS}(r)$}
            \State $C \gets C \cup \{r\}$ \Comment{$r''$ copre $r$}
            \State \textbf{break}
        \EndIf
    \EndFor
\EndFor
\State \Return $C$
\end{algorithmic}
\end{algorithm}

Le condizioni che seguono non sono ipotizzate a priori: sono ricavate dal comportamento effettivo
del sistema, letto sul bytecode decompilato, e classificate in base al ruolo che una regola
\emph{può} svolgere. Esse costituiscono il fondamento della metodologia e la base del risultato
negativo del \autoref{chap:theory}.

\begin{definition}[Condizione \Vone{} --- veto impossibile]\label{def:v1}
Sia $\NullCols$ l'insieme delle colonne del dataset che presentano almeno una cella nulla. Una
RFD $r$ \emph{non può} esercitare alcun veto se
\[
    \LHS(r) \cap \NullCols = \emptyset .
\]
\end{definition}

La condizione è necessaria e sufficiente: la verifica di coerenza confronta valori, e il confronto
richiede che entrambi i valori siano presenti. Se nessuna colonna del lato sinistro presenta
valori mancanti, la verifica non può essere attivata.

\begin{definition}[Condizione \Vtwo{} --- veto coperto]\label{def:v2}
Una RFD $r$ è \emph{coperta} se esiste una RFD $r''$ tale che:
\begin{enumerate}[label=(\roman*),leftmargin=1.8em]
    \item $\RHS(r'') = \RHS(r)$ e $\LHS(r'') \subseteq \LHS(r)$, con $\LHS(r'') \neq \emptyset$;
    \item $\theta_{r''}(c) \ge \theta_r(c)$ per ogni $c \in \LHS(r'')$ (soglie \emph{più larghe});
    \item $\theta_{\RHS}(r'') \le \theta_{\RHS}(r)$ (soglia di destra \emph{più stretta}).
\end{enumerate}
\end{definition}

\begin{proposition}[Correttezza di \Vtwo{}]\label{prop:v2}
Se $r''$ copre $r$ secondo la Definizione~\ref{def:v2} e $r''$ sopravvive alla potatura, allora ogni veto
esercitato da $r$ è implicato da un veto di $r''$.
\end{proposition}

\begin{proof}
Sia $t$ una tupla che attiva il veto di $r$. Per definizione esistono una tupla $s$ e una colonna
$c \in \LHS(r)$ tali che $s$ e $t$ concordano entro $\theta_r(c)$ su $c$, e $s$ e $t$ differiscono
su $\RHS(r)$ di più di $\theta_{\RHS}(r)$.

Poiché $\LHS(r'') \subseteq \LHS(r)$, la tupla $s$ concorda con $t$ anche su ogni colonna di
$\LHS(r'')$: per la condizione (ii) le soglie di $r''$ sono più larghe, dunque l'accordo entro
$\theta_r(c)$ implica l'accordo entro $\theta_{r''}(c)$. La condizione (i) garantisce che
$\LHS(r'')$ non sia vuoto, sicché la verifica di $r''$ ha almeno una colonna su cui essere
valutata. Infine, per la condizione (iii), la differenza su $\RHS$ che supera $\theta_{\RHS}(r)$
supera a maggior ragione $\theta_{\RHS}(r'')$. Dunque $s$ e $t$ attivano anche il veto di $r''$.
\end{proof}

\begin{remark}
La direzione delle disuguaglianze è controintuitiva rispetto alla nozione classica di dominanza
fra dipendenze, dove si è soliti richiedere che la regola dominante sia \emph{più stretta}. Qui
occorre l'opposto: la regola coprente deve essere più permissiva sul lato sinistro e più severa
su quello destro, perché è la severità sul destro a implicare il rifiuto.
\end{remark}

\begin{definition}[Condizione \Vthree{} --- regola inerte]\label{def:v3}
Una RFD $r$ è \emph{provabilmente inerte} se
\[
    \LHS(r) \cap \NullCols = \emptyset \quad \wedge \quad \RHS(r) \notin \NullCols .
\]
\end{definition}

Una regola inerte non può né generare candidati --- la colonna di destra non ha celle nulle --- né
esercitare veti --- per la \Vone{}. La sua rimozione è quindi priva di effetti per costruzione.

\begin{table}[htbp]
\centering
\caption{Dimensione delle tre classi. Le classi \Vtwo{} e \Vthree{} sono disgiunte sui dati
esaminati: la loro unione è quindi la riduzione massima ottenibile con le condizioni esatte.}
\label{tab:conditions}
\footnotesize
\begin{tabular}{llrrrrr}
\toprule
famiglia & soglia & $N$ & \Vthree{} inerte & \Vtwo{} coperta & unione \\
\midrule
Bridges    & 3  & \num{1830}  & 2 (\SI{0.11}{\percent}) & 215  & \SI{11.9}{\percent} \\
Bridges    & 9  & \num{8961}  & 6 (\SI{0.07}{\percent}) & \num{1531} & \SI{17.2}{\percent} \\
Bridges    & 15 & \num{18844} & 6 (\SI{0.03}{\percent}) & \num{3927} & \SI{20.9}{\percent} \\
Glass      & 3  & \num{1511}  & 5 (\SI{0.33}{\percent}) & 40   & \SI{3.0}{\percent} \\
Glass      & 9  & \num{11648} & 18 (\SI{0.15}{\percent}) & 861 & \SI{7.5}{\percent} \\
Glass      & 15 & \num{27445} & 24 (\SI{0.09}{\percent}) & \num{2363} & \SI{8.7}{\percent} \\
Restaurant & 3  & 25          & 0 & 0   & \SI{0.0}{\percent} \\
Restaurant & 6  & 124         & 0 & 2   & \SI{1.6}{\percent} \\
\bottomrule
\end{tabular}
\end{table}

La \autoref{tab:conditions} mostra due fatti che orientano il resto del lavoro. Il primo è che la
classe \Vthree{} è trascurabile su tutte le famiglie, non soltanto su Bridges: la ragione è
strutturale, poiché richiede regole il cui lato sinistro impieghi esclusivamente colonne mai
nulle, e tali colonne sono poche. Il secondo è che la classe \Vtwo{} è l'unica leva
significativa, e cresce con la soglia fino a poco più di un quinto dell'insieme.

% =======================================================================================
\section{Lo spazio delle strategie}\label{sec:strategy-space}

\begin{figure}[htbp]
\centering
\begin{tikzpicture}[font=\footnotesize, every node/.style={align=center}]
\draw[black!25] (0,0) grid (10,6);
\node[anchor=south west, font=\scriptsize, black!60] at (0,6.1) {garanzia dimostrabile};
\node[anchor=south east, font=\scriptsize, black!60] at (10,6.1) {garanzia assente (euristica)};
\node[anchor=north west, font=\scriptsize, black!60] at (0,-0.1) {rimozione di regole};
\node[anchor=north east, font=\scriptsize, black!60] at (10,-0.1) {riordino, nessuna rimozione};
\node[draw, rounded corners=1pt, inner sep=4pt] at (2.5,4.4) {\str{duplicates}\\ \str{inert}\\ \str{veto-cover}};
\node[draw, rounded corners=1pt, inner sep=4pt] at (7.5,4.4) {\str{dom-sub-adaptive}\\ \str{dom-sub}, \str{coverage}\\ \str{score-argmin}\textsuperscript{*}};
\node[draw, rounded corners=1pt, inner sep=4pt] at (2.5,1.5) {\emph{nessuna strategia}};
\node[draw, rounded corners=1pt, inner sep=4pt] at (7.5,1.5) {\str{exact-priority}\\ \str{precision-priority}\\ \str{intra-cluster-priority}};
\node[anchor=north west, font=\scriptsize, black!60] at (0,-0.5) {\textsuperscript{*}esatta rispetto alla sola generazione dei candidati};
\end{tikzpicture}
\caption{Collocazione delle strategie rispetto a che cosa modificano e a che cosa garantiscono.
Il quadrante in basso a sinistra è vuoto per costruzione: una strategia che non rimuove regole non
ha bisogno di garanzie, e una che le rimuove senza garanzia appartiene al quadrante in alto a
destra.}
\label{fig:taxonomy}
\end{figure}

Le diciannove strategie implementate non sono un elenco di tentativi: si distribuiscono lungo tre
dimensioni indipendenti, che ne determinano la natura e le possibilità.

\begin{description}[leftmargin=1.4em]
    \item[Che cosa cambiano.] Le strategie di \emph{rimozione} riscrivono il file eliminando
    regole; le strategie di \emph{riordino} non eliminano nulla e modificano soltanto la posizione
    delle righe. La distinzione non è di intensità ma di natura: il riordino è l'unica classe che
    preserva entrambi i ruoli per costruzione, perché nessuna regola viene sottratta al sistema, e
    il suo effetto è quindi interamente mediato dal comportamento dell'algoritmo.
    \item[Che cosa garantiscono.] Una strategia è \emph{esatta} se la sua applicazione è
    dimostrabilmente priva di effetti --- la rimozione è giustificata da una proprietà che si
    verifica sul file delle regole e sullo schema dei valori mancanti --- ed \emph{euristica} se
    la rimozione è plausibile ma non garantita. La distinzione è mantenuta in ogni fase: nelle
    strategie, nella politica di tutela dei veti e nella presentazione dei risultati.
    \item[Che informazione richiedono.] Alcune strategie decidono guardando soltanto il file delle
    regole; altre richiedono un'analisi del dataset (il calcolo statico dell'argomento del minimo,
    la stima della copertura empirica). Il costo di questa analisi è un costo di preprocessing, che
    va confrontato con il risparmio ottenuto in esecuzione: la \autoref{sec:cost} mostra che il
    confronto non è sempre favorevole.
\end{description}

Dalla combinazione delle tre dimensioni discende il criterio con cui le strategie sono presentate
nel seguito: per ciascuna si dichiarano la famiglia, il regime di garanzia, l'informazione
richiesta, la riduzione osservata e l'effetto misurato sulla qualità. Le strategie che non
soddisfano il criterio di garanzia non sono presentate come alternative equivalenti ma come
strumenti con un profilo di rischio esplicito.

% =======================================================================================
% =======================================================================================
\section{Le strategie di potatura}\label{sec:strategies}

\begin{algorithm}[htbp]
\caption{Driver di potatura. Le rimozioni garantite sono esenti dal computo del budget: essendo
prive di costo in termini di veto, addebitarle costituirebbe una doppia contabilizzazione.}
\label{alg:driver}
\begin{algorithmic}[1]
\State \textbf{input}: regole $R$, dataset $D$, strategia $S$, politica $\pi$, frazione $F$
\State $R' \gets S(R, D)$ \Comment{trasformazione richiesta dalla strategia}
\State $G \gets \{ r \in R \setminus R' : r \in C \ \textbf{or}\ r \in I \}$
       \Comment{rimozioni garantite: $C$ da \Vtwo{}, $I$ da \Vthree{}}
\State $H \gets \{ r \in R \setminus (R' \cup G) : \LHS(r) \cap \NullCols(D) \neq \emptyset \}$
       \Comment{rimozioni che sacrificano un veto}
\If{$\pi = \str{protect}$}
    \State $H \gets \emptyset$
\ElsIf{$\pi = \str{budget:F}$}
    \State ordina $H$ per indice di rischio crescente
    \State $H \gets$ le $\lceil F \cdot |H| \rceil$ regole di $H$ a rischio minore
           \Comment{le altre sono ripristinate}
\EndIf
\State \textbf{output}: $R \setminus (G \cup H)$
\end{algorithmic}
\end{algorithm}

La \autoref{tab:strategies} riassume le strategie principali, con la garanzia associata, la
riduzione osservata e l'effetto misurato sulla qualità\footnote{Le riduzioni e gli effetti
provengono da esperimenti diversi, indicati nella colonna ``effetto'' e descritti nel
\autoref{chap:experiments}: l'ablazione impiega venti configurazioni comuni su Bridges,
l'esperimento di ordinamento quindici configurazioni alla soglia 9, la campagna principale le
cinquanta configurazioni di Bridges. I valori non sono quindi confrontabili fra righe diverse
della tabella.}.

\begin{table}[htbp]
\centering
\caption{Strategie di potatura implementate. ``Garanzia'' indica se la rimozione è dimostrabilmente
priva di effetti; ``riduzione'' è l'ordine di grandezza osservato sui dataset impiegati;
``effetto'' è la variazione di imputazioni corrette rispetto alla configurazione non potata,
misurata nell'esperimento indicato. Le strategie citate nel seguito del lavoro sono in grassetto.}
\label{tab:strategies}
\footnotesize
\begin{tabular}{@{}llccp{2.9cm}@{}}
\toprule
famiglia & strategia & garanzia & riduzione & effetto misurato \\
\midrule
\multirow{3}{*}{strutturale}
  & \str{duplicates}        & sì & \SI{0}{\percent} & nessuno \\
  & \str{veto-dominance}    & no  & --- & non impiegata \\
  & \str{dom-exact}         & no  & --- & non impiegata \\
\midrule
\multirow{3}{*}{esatta}
  & \textbf{\str{veto-cover}}      & \textbf{sì} & \SI{0}{\percent}--\SI{20.9}{\percent} & nessuno (\num{113}/\num{113} configurazioni identiche) \\
  & \str{inert}             & sì & \SI{0.03}{\percent}--\SI{0.33}{\percent} & nessuno, per costruzione \\
  & \str{score-argmin}      & sì\textsuperscript{*} & fino a \SI{98}{\percent}\textsuperscript{*} & nessuna perdita di copertura, ma rallentamento fino a \num{251}$\times$ \\
\midrule
\multirow{4}{*}{euristica}
  & \textbf{\str{dom-sub-adaptive}} & no & \SI{32}{\percent}--\SI{93}{\percent} & $-20$ corrette da sola (ablazione) \\
  & \str{dom-sub}, \str{dom-sub-strict} & no & \SI{67}{\percent}--\SI{90}{\percent} & non impiegate da sole \\
  & \str{used}, \str{empirical}     & no & variabile & non impiegate da sole \\
  & \str{coverage}, \str{lhs-cap}, \str{rhs-thr} & no & variabile & non impiegate da sole \\
\midrule
\multirow{3}{*}{di riordino}
  & \textbf{\str{exact-priority}}   & no & \SI{0}{\percent} & $+25$ corrette da solo (ablazione) \\
  & \textbf{\str{precision-priority}} & no & \SI{0}{\percent} & $+53$ corrette (ordinamento) \\
  & \str{intra-cluster-priority}    & no & \SI{0}{\percent} & nessuno (\num{15}/\num{15} configurazioni identiche) \\
\bottomrule
\end{tabular}

\vspace{0.4em}
\raggedright\footnotesize\textsuperscript{*}La garanzia di \str{score-argmin} riguarda la sola
generazione dei candidati, non il potere di veto: la precisazione è essenziale ed è discussa nel
\autoref{sec:score-argmin}.
\end{table}

\paragraph{La strategia garantita: \str{veto-cover}.} Elimina le regole il cui veto è implicato da
una regola sopravvissuta, secondo la condizione \Vtwo{} della \autoref{sec:exact-conditions}. È
l'unica strategia di rimozione che offra una garanzia dimostrabile su entrambi i ruoli, e la sua
riduzione è per questo limitata: cresce con la soglia e non supera il \SI{20.9}{\percent} sui
dataset esaminati. La \autoref{sec:res-vetocover} ne verifica l'equivalenza di output.

\paragraph{Le strategie esatte di portata limitata.} \str{inert} rimuove le regole inerti secondo
la condizione \Vthree{}: la garanzia è totale e la riduzione trascurabile, circostanza che
costituisce di per sé un risultato negativo sulla praticabilità di questa via.
\str{score-argmin} preserva la generazione calcolando staticamente, per ogni cella nulla,
l'argomento del minimo che il sistema sceglierebbe: può rimuovere fino al \SI{98}{\percent} delle
regole, ma non preserva il veto, e la \autoref{sec:negative-results} mostra che il suo effetto sul
tempo è opposto a quello atteso.

\paragraph{Le strategie euristiche.} \str{dom-sub-adaptive} rimuove regole che risultano dominate
rispetto a una relazione di sussunzione adattata ai dati, con una soglia di aggressività
regolabile. È la strategia che produce le riduzioni maggiori, ed è anche quella il cui effetto
sulla qualità dipende dal dataset: la combinazione impiegata nelle campagne principali
(\str{A5b}) la associa al riordino per priorità di precisione.

\paragraph{Le strategie di riordino.} Non eliminano regole e non sono soggette ad alcuna politica
di tutela. \str{exact-priority} dispone per prime le regole che producono il valore esatto,
\str{precision-priority} quelle con la precisione empirica più alta, \str{intra-cluster-priority}
interviene sull'ordine interno ai cluster. Il loro effetto è mediato dalle due procedure di
riempimento (\autoref{sec:two-branches}) ed è oggetto dell'esperimento di ordinamento
i cui esiti sono riportati nella \autoref{sec:negative-results}.

\paragraph{Componibilità delle strategie.} Le famiglie non sono intercambiabili e la loro
composizione non è libera. Il \emph{riordino} si compone con qualunque strategia di rimozione,
perché non sottrae regole e non consuma budget di veto; la \emph{rimozione garantita} si compone
con l'euristica, ma il suo effetto si somma a quello dell'euristica soltanto se le due agiscono su
regole diverse. \str{score-argmin} non si compone con le altre strategie di rimozione: la sua
garanzia riguarda la generazione e vale rispetto all'insieme di partenza, sicché una rimozione
successiva può invalidare l'argomento del minimo che essa ha calcolato staticamente. È la ragione
per cui la combinazione impiegata nelle campagne è formata da un riordino e da una sola strategia
di rimozione euristica, e non da più rimozioni in cascata.

\section{Perché la generazione non ammette una condizione esatta della stessa natura}
\label{sec:generation-limit}

La condizione \Vtwo{} sfrutta una proprietà logica del veto: se una regola ne implica un'altra, la
seconda è superflua. Per la generazione una proprietà analoga non sussiste, e la ragione è
strutturale: il valore imputato dipende dal \emph{miglior punteggio} fra i candidati di un cluster,
e il miglior punteggio dipende dall'insieme delle regole presenti. Rimuovere una regola può quindi
cambiare l'argomento del minimo anche quando quella regola non è mai essa stessa il minimo. La
generazione è perciò una funzione dell'\emph{insieme}, non un insieme di contributi indipendenti,
e nessuna proprietà locale di una singola regola può certificare che la sua assenza non alteri
l'esito.

Da qui la distinzione fra la condizione \str{score-argmin}, che preserva la generazione
\emph{calcolando} staticamente l'argomento del minimo sul dataset --- quindi pagando un costo di
preprocessing proporzionale alla dimensione dei dati --- e le strategie euristiche, che rinunciano
a questa proprietà in cambio di riduzioni maggiori. La stessa asimmetria fra i due ruoli è la
ragione per cui, come dimostra il \autoref{chap:theory}, una potatura che operi sul solo file
delle regole non possa superare la classe delle regole coperte e delle regole inerti senza
rinunciare a una garanzia.

% =======================================================================================
\section{Integrazione con il sistema: il contratto di potatura}\label{sec:integration}

\paragraph{Preprocessing esterno.} La potatura si inserisce fra l'estrazione delle RFD e
l'esecuzione del sistema: il file delle regole viene riscritto secondo la strategia scelta, e il
sistema viene eseguito sul file risultante. Il dataset non è modificato, e l'eseguibile non è
toccato. La scelta ha il pregio di non richiedere alcuna modifica al sistema --- che resta quello
originale, verificato per checksum --- ed è stata mantenuta come vincolo per l'intera durata del
lavoro; la \autoref{sec:instrumentation} mostra che il vincolo è stato rimosso per una parte
dell'analisi, ma le strategie di potatura operano tutte nella forma esterna.

\paragraph{Politica sui veti.}\label{sec:veto-policy}
Poiché la rimozione di una regola può sacrificarne il potere di veto, ogni strategia di rimozione
è accompagnata da una politica esplicita, scelta fra tre:

\begin{itemize}[leftmargin=1.4em]
    \item \str{protect}: ogni regola in grado di esercitare un veto viene ripristinata. La
    riduzione si annulla, ma la garanzia è totale;
    \item \str{budget:F}: è consentito rimuovere una frazione $F$ delle regole portatrici di veto,
    scelte secondo un ordinamento di rischio;
    \item \str{allow}: nessuna tutela. La riduzione è massima e il potere di veto è
    esplicitamente sacrificato.
\end{itemize}

Le rimozioni il cui veto è \emph{provato coperto} dalla condizione \Vtwo{} sono esenti dal
computo del budget: essendo prive di costo in termini di veto, addebitarle costituirebbe una
doppia contabilizzazione. Questa precisazione, insieme a un difetto di correttezza della prima
implementazione, è discussa nel \autoref{chap:defects}.

\paragraph{Invarianti del contratto.} Il contratto di potatura è soggetto a tre invarianti,
verificate automaticamente dall'ambiente di valutazione. Il file prodotto deve essere un
sottoinsieme del file di partenza con lo stesso formato, e il dataset non deve essere alterato; le
righe devono conservare contenuto e ordinamento relativo, così che una strategia di riordino sia
distinguibile da una di rimozione; se il file prodotto coincide con quello di partenza, il
confronto è \emph{vacuo} e va segnalato anziché conteggiato come esito nullo. La terza invariante
è quella che nella pratica si è rivelata più utile, perché un numero non trascurabile di
configurazioni produce una potatura vuota, e trattarle come evidenza di assenza di effetto
avrebbe alterato le conclusioni.


    \chapter{Teoria: che cosa si può potare con garanzia}\label{chap:theory}


\noindent Questo capitolo formula e dimostra il risultato negativo che delimita il campo di
applicabilità della potatura esterna. Dopo aver fissato che cosa si intende per comportamento del
sistema, per attivabilità di una regola e per rimozione certificata --- e dopo aver illustrato le
definizioni su un caso minimo --- il capitolo enuncia il teorema e lo dimostra, ne trae la
conseguenza sul target di riduzione separando ciò che è dimostrato da ciò che è misurato, mostra
che la classe certificata non è vuota e che il limite non dipende dal parametro di ricerca dei
copritori, e ne precisa infine la portata, indicando che cosa resta fuori dal risultato.

% =======================================================================================
\section{Impostazione formale}\label{sec:theory-setup}

\begin{definition}[Sistema e comportamento]\label{def:behaviour}
Sia $D$ un dataset con schema $\mathcal{A}$ e sia $R$ un insieme finito di RFD. Si indica con
$B(R, D)$ il \emph{comportamento} del sistema su $D$ con l'insieme di regole $R$: l'insieme delle
celle valorizzate e i valori loro assegnati. Due insiemi di regole $R$ e $R'$ sono
\emph{equivalenti} su $D$ se $B(R, D) = B(R', D)$.
\end{definition}

Il comportamento è l'osservabile del sistema: nella notazione dell'implementazione corrisponde ai
file di output delle celle popolate e dei risultati. La nozione di equivalenza è quella che gli
esperimenti del \autoref{chap:results} verificano per confronto di impronte.

\begin{definition}[Attivabilità di una regola]\label{def:activability}
Sia $\NullCols(D)$ l'insieme delle colonne di $D$ che presentano almeno una cella nulla. Una RFD
$r$ è \emph{attivabile in veto} su $D$ se $\LHS(r) \cap \NullCols(D) \neq \emptyset$; è
\emph{attivabile in generazione} su $D$ se $\RHS(r) \in \NullCols(D)$.
\end{definition}

Le due condizioni sono di natura diversa. La prima è necessaria e sufficiente, ed è la condizione
\Vone{} della \autoref{sec:exact-conditions}: la verifica di coerenza confronta valori, e il
confronto richiede che entrambi siano presenti. La seconda è necessaria ma non sufficiente: perché
una regola generi un candidato occorre anche che esista, nella colonna di destra, un cluster con
cui confrontare la tupla incompleta. Ai fini del teorema è però la sola necessità che serve,
perché è la possibilità di attivazione a escludere la certificabilità.

\begin{definition}[Potatura esterna]\label{def:external-pruning}
Una \emph{strategia di potatura esterna} è una funzione $\Pi$ che, applicata a una coppia
$(R, D)$, produce un insieme $\Pi(R, D) \subseteq R$ usando esclusivamente il file delle regole,
lo schema di $D$ e il profilo delle colonne nulle $\NullCols(D)$. In particolare $\Pi$ non accede
al comportamento del sistema né può osservare quali regole siano attivate in esecuzione.
\end{definition}

La restrizione è quella imposta dal \autoref{chap:methodology}: la potatura riscrive il file e non
modifica l'eseguibile. È una restrizione reale, non una comodità dimostrativa: è ciò che rende la
metodologia applicabile a un sistema distribuito senza sorgenti.

\begin{definition}[Rimozione certificata]\label{def:certified}
Una rimozione $r \in R \setminus \Pi(R, D)$ è \emph{certificata} se esiste una condizione $\Phi$,
decidibile con l'informazione ammessa dalla Definizione~\ref{def:external-pruning}, tale che per
ogni dataset $D'$ con lo schema e il profilo delle colonne nulle di $D$ valga
\[
    \Phi(r, R, D) \;\Longrightarrow\; \text{$r$ non è attivabile in alcuno dei due ruoli su $D'$,}
\]
oppure
\[
    \Phi(r, R, D) \;\Longrightarrow\; \text{il veto di $r$ su $D'$ è implicato da una regola di
    $\Pi(R, D)$.}
\]
Una strategia è certificata se tutte le sue rimozioni lo sono.
\end{definition}

La definizione merita una precisazione. La certificazione non è richiesta rispetto al solo dataset
in esame, ma rispetto a ogni dataset che con esso condivide schema e profilo di valori mancanti: è
questa quantificazione a dare forza al risultato, ed è anche la ragione per cui la potatura
esterna non può giovarsi di circostanze accidentali. Una regola la cui rimozione non altera il
comportamento \emph{su questo} dataset può benissimo alterarlo su un altro con lo stesso profilo di
nullità, e una strategia che operi senza eseguire il sistema non ha modo di distinguere i due
casi. La \autoref{sec:target-impossible} mostra che la distinzione è verificabile a posteriori, e
che le rimozioni osservate come innocue ricadono comunque nelle classi certificate.

% =======================================================================================
% =======================================================================================
\section{Un esempio minimo}\label{sec:theory-example}

Le tre condizioni si leggono meglio su un caso concreto. Si consideri uno schema con tre
attributi $A$, $B$, $C$, un profilo di valori mancanti $\NullCols = \{A, C\}$ e quattro regole; la
\autoref{tab:theory-example} le classifica.

\begin{table}[htbp]
\centering
\caption{Esempio minimo di classificazione delle regole. La tabella è illustrativa: serve a
rendere leggibili le definizioni della sezione precedente, non riporta misure.}
\label{tab:theory-example}
\footnotesize
\begin{tabular}{@{}llll@{}}
\toprule
regola & espressione & classe & rimozione certificabile \\
\midrule
$r_2$ & $(A \le 3) \to (C \le 0)$ & né inerte né coperta; è la copritrice di $r_1$ & no: è la regola che copre \\
$r_1$ & $(A \le 1) \to (C \le 0)$ & \Vtwo{}: coperta da $r_2$ & sì, purché $r_2$ sopravviva \\
$r_3$ & $(B \le 1) \to (C \le 0)$ & né inerte né coperta & no: può generare candidati \\
$r_4$ & $(B \le 1) \to (B \le 0)$ & \Vthree{}: inerte & sì, per costruzione \\
\bottomrule
\end{tabular}
\end{table}

Tre aspetti dell'esempio meritano di essere evidenziati, perché sono quelli che rendono la
condizione \Vtwo{} controintuitiva. Il primo è la direzione delle soglie: $r_2$ copre $r_1$ pur
essendo \emph{più permissiva} a sinistra ($3$ contro $1$) e non più debole a destra (la soglia di
destra è la stessa); nella nozione classica di dominanza fra dipendenze si richiederebbe l'opposto.
Il secondo è che $r_1$ è rimovibile solo \emph{finché} $r_2$ sopravvive: la condizione si valuta
sull'insieme corrente delle regole, non su quello di partenza, ed è la precisazione che il
\autoref{sec:d5} ha reso necessaria. Il terzo è che $r_3$ --- che pure ha un lato sinistro senza
valori mancanti, e quindi non può esercitare alcun veto --- \emph{non} è rimovibile senza
rinunciare alla garanzia, perché la sua colonna di destra contiene celle nulle e la regola può
quindi generare candidati: è l'asimmetria fra i due ruoli che il teorema generale formalizza.

\section{Il teorema negativo}\label{sec:negative-theorem}

\begin{theorem}[Limite della potatura esterna certificata]\label{thm:negative}
Siano $D$ un dataset, $R$ un insieme di RFD e $\Pi$ una strategia di potatura esterna in possesso
della certificazione della Definizione~\ref{def:certified}. Allora ogni regola
$r \in R \setminus \Pi(R, D)$ è provabilmente inerte secondo la condizione \Vthree{}, oppure il
suo veto è implicato da una regola sopravvissuta.
\end{theorem}

Il teorema non afferma che le regole fuori da queste due classi non possano essere rimosse senza
effetti: afferma che la loro rimozione non può essere \emph{certificata} con l'informazione
disponibile a una strategia esterna. La distinzione fra le due affermazioni è il contenuto
dell'ultima sezione del capitolo.

% =======================================================================================
\section{Dimostrazione}\label{sec:theory-proof}

\begin{proof}
Sia $r \in R \setminus \Pi(R, D)$ una regola rimossa, e si supponga che $r$ non sia inerte secondo
la Definizione~\ref{def:v3}, cioè che valga
\[
    \LHS(r) \cap \NullCols(D) \neq \emptyset \quad \vee \quad \RHS(r) \in \NullCols(D).
\]
Si distinguono i due casi.

\emph{Primo caso: $\RHS(r) \in \NullCols(D)$.} Per la Definizione~\ref{def:activability} la regola
è attivabile in generazione: esiste almeno una cella nulla nella colonna di destra, e la regola
partecipa quindi all'insieme dei candidati per quella cella. Il valore che il sistema scrive
dipende dall'argomento del minimo fra i candidati del cluster, e l'argomento del minimo dipende
dall'insieme delle regole presenti (\autoref{sec:generation-limit}); nessuna proprietà della
singola regola $r$, decidibile dal file, determina se la sua presenza cambi l'esito. La condizione
$\Phi$ non può quindi stabilire la prima delle due alternative della
Definizione~\ref{def:certified}.

\emph{Secondo caso: $\LHS(r) \cap \NullCols(D) \neq \emptyset$.} Per la condizione \Vone{}, che è
necessaria e sufficiente, la regola è attivabile in veto: esistono un dataset $D'$ con lo schema e
il profilo di nullità di $D$ e una coppia di tuple su cui la verifica di coerenza di $r$ viene
eseguita e fallisce. Se $r$ non ricade nel primo caso, la prima alternativa della
Definizione~\ref{def:certified} è esclusa anche qui, e la certificazione può venire soltanto dalla
seconda: il veto di $r$ deve essere implicato da una regola che sopravvive alla potatura.

Dunque ogni rimozione certificata è una regola inerte oppure una regola il cui veto è implicato da
una regola sopravvissuta.
\end{proof}

La dimostrazione poggia su due fatti stabiliti altrove: l'esattezza della condizione \Vone{}, che
discende dalla semantica del confronto fra valori, e la dipendenza della generazione
dall'argomento del minimo, che discende dall'analisi delle due procedure di riempimento. Il primo
è un fatto logico; il secondo è una proprietà dell'algoritmo, verificata sul codice e confermata
dal fatto, misurato nella \autoref{sec:res-vetocover}, che le strategie di rimozione che non
tutelano la generazione producono output diversi.

% =======================================================================================
\section{Conseguenza sul target di riduzione}\label{sec:target-impossible}

Il teorema si traduce in un vincolo quantitativo solo una volta stabilito quali regole, fra quelle
certificabili, siano \emph{riconosciute} come tali dalla metodologia. La condizione \Vtwo{} è una
condizione \emph{sufficiente} di copertura del veto: quando è soddisfatta, la regola è eliminabile
con garanzia. Non è detto che sia l'unica condizione sufficiente possibile, ed è opportuno
distinguere le due affermazioni.

\begin{corollary}[Incompatibilità del target]\label{cor:target}
Se la copertura del veto è decisa mediante la condizione \Vtwo{}, allora la riduzione ottenibile
da una potatura esterna certificata è limitata dalla dimensione dell'unione delle classi \Vtwo{} e
\Vthree{}.
\end{corollary}

\paragraph{La classe certificata non è vuota.} Un teorema di impossibilità va difeso dalla prima
obiezione che riceve, cioè che la classe su cui il risultato morde sia vuota. Non lo è: sui dataset
esaminati la condizione \Vthree{} individua poche regole --- \num{6} su Bridges a soglia 9 --- ma la
condizione \Vtwo{} ne individua \num{1531}, pari al \SI{17.2}{\percent} dell'insieme
(\autoref{tab:conditions}), e la riduzione ottenibile cresce con la soglia fino al
\SI{20.9}{\percent}. La condizione è inoltre \emph{decidibile} senza eseguire il sistema: la ricerca
dei copritori della \autoref{alg:cover} esamina al più tutte le coppie di regole con la stessa
colonna di destra, ed è quindi di ordine quadratico nella dimensione del file, con il parametro
$\mu$ a limitare il numero di candidati esaminati per ciascuna regola.

La \autoref{tab:conditions} fornisce la misura di quell'unione sui dataset impiegati: essa vale
dall'\SI{0.0}{\percent} al \SI{20.9}{\percent}, con il massimo raggiunto da Bridges alla soglia più
alta. Ne segue che un obiettivo di riduzione pari o superiore al \SI{50}{\percent} è incompatibile
con la garanzia, e che l'incompatibilità non dipende dalla strategia adottata ma dalla struttura
delle classi. La precisazione è necessaria perché il corollario è una proposizione \emph{misurata}
e non un teorema: la limitazione è dimostrata, la sua entità numerica è osservata su otto casi, e
nulla esclude che su altri dataset la classe \Vtwo{} sia più ampia. Ciò che il teorema esclude è
che una via diversa dalla copertura del veto possa essere percorsa \emph{con garanzia}; ciò che
resta aperto è quanto la copertura del veto possa valere su dati diversi da quelli esaminati.

% =======================================================================================
\section{Quanto è stretto il limite}\label{sec:maxextra}

La ricerca delle regole coprenti impiega un parametro che limita la differenza di ampiezza fra il
lato sinistro della regola coprente e quello della regola coperta. Il parametro è una scelta di
ricerca, non una condizione di correttezza: la \autoref{prop:v2} vale per qualunque valore. Si è
quindi verificato se il tetto di riduzione osservato fosse un artefatto del valore adottato. Lo
sweep completo --- valori da 3 a 20, su otto casi --- mostra che la riduzione è \emph{piatta}: su
\num{8939} rimozioni complessive, allargare la ricerca al massimo ne aggiunge \emph{una}. Il limite
è pertanto intrinseco al criterio, e il teorema negativo non è aggirabile per questa via. La
verifica è condotta da \texttt{pruning/verify\_veto\_sweep.py}, i cui risultati sono raccolti in
\texttt{results/results\_maxextra\_bridges15.json}.

% =======================================================================================
\section{Portata del risultato}\label{sec:theory-scope}

Tre restrizioni delimitano il teorema, e conviene dichiararle esplicitamente perché ciascuna
indica una direzione di lavoro.

\paragraph{Solo rimozione.} La Definizione~\ref{def:external-pruning} considera trasformazioni che
producono un sottoinsieme del file di partenza. Una trasformazione che \emph{riscriva} le regole
--- restringendo o allargando i lati sinistri, modificando le soglie, fondendo più regole in una
--- non è coperta dal risultato, e non è detto che vi ricada: la condizione \Vtwo{} mostra del
resto che una regola può essere eliminata quando il suo contenuto informativo è conservato da
un'altra, e nulla vieta di cercare la stessa conservazione per via sintattica anziché per
soppressione. È, questa, la direzione di sviluppo più promettente indicata nelle
\autoref{sec:future-work}.

\paragraph{Solo informazione esterna.} La certificazione richiesta è quella ottenibile dal file,
dallo schema e dal profilo dei valori mancanti. Una strategia che eseguisse il sistema, o che
disponesse di una sua descrizione simbolica, potrebbe in linea di principio certificare rimozioni
ulteriori: è il caso di \str{score-argmin}, che calcola staticamente l'argomento del minimo sul
dataset e preserva la generazione, ma non il veto, e resta quindi fuori dalla classe garantita
(\autoref{sec:score-argmin}).

\paragraph{Completezza della condizione di copertura.} Il teorema caratterizza le rimozioni
certificabili, non le rimozioni innocue. Stabilire se la condizione \Vtwo{} esaurisca le regole il
cui veto è implicato da un'altra regola è un problema aperto, e la sua soluzione non cambierebbe
la natura del risultato ma potrebbe ampliarne la portata quantitativa. La
\autoref{sec:res-vetocover} fornisce, su questo punto, un'evidenza indiretta: le rimozioni
effettivamente osservate come innocue ricadono tutte nelle classi certificate, e la copertura
verificata a posteriori coincide con quella prevista.

\paragraph{La riduzione interna del sistema.} La certificazione è definita rispetto all'insieme
delle regole contenute nel file, ma il sistema ne consuma un sottoinsieme determinato da un
criterio \emph{proprio}: all'avvio esso rimuove le regole a soglia di destra nulla il cui lato
sinistro si comporta come una chiave e le destina a una procedura distinta
(\autoref{sec:dual-role}). Due osservazioni ne seguono. La prima è che questa riduzione interna
non è modellata dal teorema: le classi \Vtwo{} e \Vthree{} sono calcolate sul file, e nulla
garantisce che una regola coprente sopravviva alla riduzione interna. La seconda è che la
garanzia resta comunque valida sui dati esaminati, perché la proprietà di chiave implica che il
veto della regola coprente e quello della regola coperta siano entrambi non esercitabili: se il
lato sinistro della coprente è una chiave, ogni suo soprainsieme --- e quindi anche il lato
sinistro della coperta --- lo è, e nessuna coppia di tuple distinte può attivare la verifica. La
riduzione interna non può dunque sottrarre un veto che il teorema presume implicato; resta il
fatto che il modello formale non la rappresenta, e che una sua inclusione renderebbe
l'enunciato più generale.


    \chapter{Ipotesi e criteri di falsificazione}\label{chap:hypotheses}


\noindent Questo capitolo documenta il processo che ha condotto dalle osservazioni iniziali alla
formulazione delle ipotesi. La sua inclusione risponde a una ragione metodologica: il lavoro non è
consistito nell'applicare una tecnica nota a un sistema noto, ma nel ricostruire il comportamento
di un sistema di cui non erano disponibili i sorgenti. Le ipotesi hanno quindi svolto una funzione
euristica, e la loro falsificazione è parte del risultato. Il capitolo ricostruisce il processo di
analisi, enuncia le ipotesi primarie con la loro derivazione e il criterio di falsificazione ---
senza anticiparne l'esito, che è discusso insieme ai risultati --- e raccoglie separatamente le
domande esplorative. Il protocollo di verifica e la metodologia statistica, che delle ipotesi
costituiscono il banco di prova, sono esposti nel capitolo sulla valutazione sperimentale.

% =======================================================================================
\section{Come sono nate le ipotesi}\label{sec:analysis-process}

L'analisi si è svolta in quattro passi, ciascuno dei quali ha determinato il successivo.

\begin{enumerate}[leftmargin=1.6em]
    \item \textbf{Disassemblaggio.} Il sistema è distribuito senza sorgenti. La ricostruzione
    della logica è avvenuta per disassemblaggio del bytecode, con particolare attenzione alle
    procedure che consumano l'insieme delle regole.
    \item \textbf{Individuazione del doppio ruolo.} È emerso che una RFD partecipa a due
    procedure con requisiti opposti (\autoref{sec:dual-role}). Questa osservazione ha trasformato
    il problema da ``selezionare le regole utili'' a ``selezionare le regole utili sapendo che
    \emph{utile} ha due significati incompatibili''.
    \item \textbf{Derivazione delle condizioni.} Dal comportamento delle due procedure sono state
    ricavate le condizioni \Vone{}, \Vtwo{}, \Vthree{} e il teorema negativo che ne segue
    (\autoref{chap:theory}).
    \item \textbf{Verifica sperimentale.} Le condizioni sufficienti sono state verificate su
    stati reali; le strategie euristiche sono state misurate su più famiglie di dataset, con
    particolare attenzione a \emph{dove} esse funzionano e dove no.
\end{enumerate}

Un elemento del processo merita di essere esplicitato, perché ha modificato il corso del lavoro.
L'analisi del passo 3 ha prodotto un risultato negativo, il teorema della
\autoref{sec:negative-theorem}, che esclude la possibilità di raggiungere l'obiettivo di
riduzione mantenendo la garanzia. Una parte non trascurabile del lavoro è consistita nel
\emph{verificare} che quel risultato negativo non fosse un artefatto
dell'implementazione, prima di accettarlo come limite del problema
(\autoref{sec:maxextra}).

\paragraph{Origine delle ipotesi.} Il lavoro è esplorativo: le ipotesi non precedono l'analisi ma
nascono da essa, e la circostanza va dichiarata perché incide sul valore probatorio dei test. La
\autoref{tab:hypotheses} distingue perciò, per ogni ipotesi, l'\emph{origine}: le ipotesi
\emph{derivate} discendono da una proprietà stabilita prima di osservare i dati sperimentali ---
una condizione dimostrata, una misura di struttura, l'analisi del costo --- e la loro previsione è
argomentabile a priori; le ipotesi \emph{esplorative} sono state formulate durante l'analisi, e il
loro esito va inteso come indicazione per ulteriori verifiche e non come conferma indipendente.
Nessuna delle ipotesi primarie è stata formulata dopo aver visto l'esito dell'esperimento che la
verifica; la distinzione è mantenuta anche nel capitolo dei risultati, dove gli esiti sono
riportati per ciascuna.

% =======================================================================================
\section{Ipotesi primarie}\label{sec:hypotheses-list}

Le ipotesi primarie sono sette e corrispondono alle domande che il lavoro pone: che cosa si può
potare con garanzia, che cosa si può preservare senza garanzia, se esista una leva di qualità, se
tale leva generalizzi, da che cosa dipenda il guadagno quando si manifesta, e che rapporto vi sia
fra riduzione e tempo. La \autoref{tab:hypotheses} le enuncia con la previsione e il criterio di
falsificazione; l'esito non è qui riportato.

\begin{table}[htbp]
\centering
\caption{Ipotesi primarie, previsione, criterio di falsificazione e sezione in cui ciascuna è
risolta. L'origine --- derivata da una proprietà stabilita prima della misura, oppure formulata
durante l'analisi --- è dichiarata nella derivazione che segue alla tabella.}
\label{tab:hypotheses}
\footnotesize
\begin{tabular}{@{}cp{3.3cm}p{3.3cm}p{2.3cm}c@{}}
\toprule
\# & ipotesi & previsione & falsificazione & risolta in \\
\midrule
H1 & \str{score-argmin} con veti protetti & riduzione $\ge \SI{30}{\percent}$ con output identico & riduzione inferiore alla soglia & \ref{sec:negative-results} \\
H6 & \str{veto-cover} su tutte le famiglie & output identico alla configurazione di partenza & una sola configurazione con output diverso & \ref{sec:res-vetocover} \\
H7 & verifica empirica di \Vtwo{} & zero discrepanze fra copertura prevista e verificata & una sola discrepanza & \ref{sec:res-vetocover} \\
H8 & \str{precision-priority} da solo & guadagno di qualità & nessun guadagno o perdita & \ref{sec:ablation} \\
H11 & \str{precision-priority} su Glass e Restaurant & guadagno anche fuori da Bridges & perdita su entrambe le famiglie & \ref{sec:res-boundary} \\
H14 & attribuzione del guadagno alla precisione dei cluster & un riordino casuale non produce lo stesso guadagno & il riordino casuale produce lo stesso guadagno & \ref{sec:negative-results} \\
H16 & proporzionalità fra riduzione e tempo & riduzioni maggiori producono tempi minori & una configurazione con riduzione maggiore e tempo non inferiore & \ref{sec:res-cost} \\
\bottomrule
\end{tabular}
\end{table}

\paragraph{H1 --- La strategia esatta lato generazione.} \str{score-argmin} preserva la generazione
dei candidati per costruzione, perché calcola staticamente l'argomento del minimo che il sistema
sceglierebbe. Ci si attende pertanto che, associata a una politica che tuteli integralmente il
potere di veto, essa produca un output identico a quello della configurazione di partenza, e che
la riduzione residua sia ancora significativa: la politica di tutela ripristina le regole
portatrici di veto, ma non quelle la cui rimozione è provata priva di costo dalla condizione
\Vtwo{}, e la classe \Vtwo{} vale fino a poco più di un quinto dell'insieme
(\autoref{tab:conditions}). La previsione è una riduzione non inferiore al \SI{30}{\percent}.
L'ipotesi è derivata: entrambe le premesse --- l'esattezza lato generazione e la dimensione della
classe coperta --- sono stabilite prima della misura.

\paragraph{H6 --- La classe garantita.} Se la condizione \Vtwo{} è corretta, la sua applicazione
deve essere invisibile al sistema: le regole rimosse hanno un veto implicato da una regola
sopravvissuta, e la proposizione della \autoref{prop:v2} assicura che nessun veto venga perso. La
previsione è che l'output coincida con quello della configurazione di partenza su \emph{tutte} le
configurazioni della griglia, non soltanto su quelle di una famiglia. È l'ipotesi con il criterio
di falsificazione più severo, perché una sola configurazione difforme la falsifica.

\paragraph{H7 --- Verifica indipendente della copertura.} La correttezza della condizione \Vtwo{}
è stabilita per via logica, ma la sua \emph{implementazione} potrebbe essere difettosa: la ricerca
dei copritori è un algoritmo, e un algoritmo può sbagliare. La previsione è che, controllando a
posteriori su stati reali del dataset che ogni regola rimossa abbia effettivamente un copritore
sopravvissuto che esercita il veto sugli stessi testimoni, non emerga alcuna discrepanza. La
verifica è indipendente perché non riusa il codice della condizione ma ricostruisce i testimoni.

\paragraph{H8 --- La leva di qualità.} L'analisi del sistema mostra che l'esito di
un'imputazione dipende dall'ordine in cui i cluster vengono tentati
(\autoref{sec:two-branches}), e che la posizione di una regola nell'ordine di scansione non
richiede la rimozione di alcunché. Ne segue la previsione che un riordino che disponga per prime
le regole con la precisione empirica più alta migliori il risultato senza toccare il potere di
veto. L'ipotesi è derivata dalla struttura delle due procedure, non da un'osservazione sui
risultati.

\paragraph{H11 --- Generalizzazione della leva di qualità.} Se il meccanismo del guadagno è la
composizione dei candidati e non una proprietà di Bridges, lo stesso riordino deve produrre un
guadagno anche sulle altre famiglie. La previsione è che il guadagno si manifesti su Glass e su
Restaurant; la falsificazione è una perdita su entrambe. L'ipotesi è deliberatamente severa: è la
generalizzazione a decidere se il guadagno osservato su una famiglia sia un fenomeno del sistema o
un accidente del dataset.

\paragraph{H14 --- Attribuzione del guadagno.} Il riordino per priorità di precisione impiega un
segnale --- la precisione empirica dei cluster --- e la spiegazione naturale del guadagno è che
quel segnale sia la causa. Se la spiegazione è corretta, allora un riordino casuale non deve
produrre lo stesso guadagno: la previsione è che il guadagno del riordino informato superi in modo
sistematico quello di permutazioni con seme fissato. Si tratta di un'ipotesi di \emph{confutazione
di una spiegazione} più che di verifica di una strategia, ed è classificata come esplorativa
perché formulata dopo aver osservato il guadagno di cui cerca la causa.

\paragraph{H16 --- Rapporto fra riduzione e tempo.} L'analisi del costo
(\autoref{sec:cost}) mostra che il tempo è assorbito dalla verifica dei veti, che esamina
l'insieme delle regole per ogni candidato: se il costo per candidato è proporzionale al numero di
regole, la riduzione dell'insieme deve tradursi in una riduzione del tempo. La struttura del
sistema suggerisce però una previsione opposta: il veto non si limita a costare, ma \emph{scarta}
i candidati, e un insieme di regole più povero può lasciare in gara più alternative e quindi
moltiplicare le verifiche. L'ipotesi assume la prima previsione; la seconda è la ragione per cui
il criterio di falsificazione è enunciato in forma stretta --- una sola configurazione con
riduzione maggiore e tempo non inferiore falsifica l'ipotesi --- e non come semplice assenza di
correlazione.

% =======================================================================================
\section{Domande esplorative}\label{sec:exploratory}

Le restanti ipotesi non hanno il carattere di previsione argomentabile a priori: riguardano
varianti di parametri, strategie alternative e proprietà di portata locale, e sono state
formulate nel corso dell'analisi. Sono raccolte separatamente, senza esito, perché il loro ruolo è
delimitare lo spazio delle spiegazioni e indicare che cosa non funziona, non confermare una
previsione. La \autoref{tab:exploratory} le elenca con la domanda a cui rispondono.

\begin{table}[htbp]
\centering
\caption{Domande esplorative, con la sezione in cui ciascuna è risolta. Sono formulate durante
l'analisi; il loro esito è riportato nella \autoref{sec:hypotheses-outcomes} e non ha valore di
conferma indipendente.}
\label{tab:exploratory}
\footnotesize
\begin{tabular}{@{}cp{7.6cm}c@{}}
\toprule
\# & domanda & risolta in \\
\midrule
H2  & \str{score-argmin} preserva la copertura anche su Glass e Restaurant? & \ref{sec:negative-results} \\
H3  & esiste una strategia di copertura esatta delle celle, e a costo nullo? & \ref{sec:negative-results} \\
H4  & la conversione delle regole di tipo \texttt{B} in tipo \texttt{D} produce un guadagno gratuito? & \ref{sec:d1} \\
H5  & il riordino delle righe del dataset produce un guadagno? & \ref{sec:negative-results} \\
H9  & riordino e potatura sommano i loro effetti? & \ref{sec:ablation} \\
H10 & una regolarizzazione della stima delle precisioni produce un guadagno stabile? & \ref{sec:negative-results} \\
H12 & la versione regolarizzata generalizza meglio delle altre? & \ref{sec:negative-results} \\
H13 & il riordino entro il cluster è una leva di qualità? & \ref{sec:negative-results} \\
H15 & l'ordinamento delle rimozioni sotto vincolo di budget è migliorabile? & \ref{sec:negative-results} \\
\bottomrule
\end{tabular}
\end{table}

% =======================================================================================


    \chapter{Valutazione sperimentale}\label{chap:experiments}


\noindent Questo capitolo descrive la valutazione sperimentale. Vi sono esposti il protocollo con
cui le ipotesi formulate nel capitolo precedente sono tradotte in confronti controllati e la
metodologia statistica adottata per sintetizzarli; l'infrastruttura materiale su cui le campagne
si fondano; i dataset, la griglia delle configurazioni e le metriche con cui i risultati sono
misurati; il termine di paragone banale che dà una scala ai valori assoluti; gli esperimenti
condotti, con il collegamento esplicito alle ipotesi che essi verificano; e lo studio di ablazione
che isola il contributo dei singoli componenti della strategia principale. L'effetto della
potatura sul tempo di esecuzione, che non è una proprietà delle configurazioni ma un esito, è
discusso insieme ai risultati.

% =======================================================================================
\section{Protocollo di verifica}\label{sec:verification}

Il protocollo di verifica si fonda su cinque scelte, adottate per evitare errori di
interpretazione che si erano effettivamente presentati nel corso del lavoro.

\paragraph{Confronto appaiato.} Ogni configurazione potata è confrontata con la medesima
configurazione non potata, sullo stesso dataset, la stessa esecuzione e la stessa soglia. Il
confronto fra medie aggregate su configurazioni diverse è stato evitato perché sensibile alla
composizione del campione.

\paragraph{Unità di analisi dichiarata.} L'unità di analisi è la \emph{configurazione}, cioè la
tripletta dataset, soglia e ripetizione, e non la cella nulla. La scelta è motivata dalla
struttura del sistema: le imputazioni di una stessa esecuzione non sono indipendenti, perché
condividono l'ordine delle regole e i cluster, e trattarle come osservazioni separate
moltiplicherebbe artificialmente la numerosità campionaria. Ne segue che la numerosità dei test è
quella delle configurazioni --- decine, non migliaia --- e che la potenza statistica è
conseguentemente limitata.

\paragraph{Controllo della vacuità.} Il confronto è privo di contenuto in due circostanze
distinte. La prima è che il file di regole prodotto dalla potatura sia identico a quello di
partenza: la potatura non ha fatto nulla e non vi è nulla da misurare. La seconda, meno evidente,
è che la configurazione di partenza non produca \emph{alcuna} imputazione: la potatura non è
vacua, ma il confronto lo è, perché tutte le grandezze di qualità valgono zero in entrambi i
bracci. Il sistema di valutazione segnala entrambe le circostanze, e le configurazioni
corrispondenti sono escluse dai test e contate a parte: trattarle come esiti nulli introdurrebbe
una massa di zeri artificiale, e le tratterebbe come evidenza di assenza di effetto.

\paragraph{Controllo casuale.} Per attribuire un effetto a un ordinamento --- e non al mero atto
di riordinare --- è necessario un termine di paragone casuale. Sono state impiegate permutazioni
con seme fissato, in numero di tre e indipendenti fra loro; il confronto è condotto sia contro
ciascuna permutazione, sia contro la loro media per configurazione.

\paragraph{Verifica indipendente delle condizioni.} Le condizioni teoriche sono verificate su
stati reali del dataset, non soltanto per implicazione logica: per ogni regola rimossa dalla
condizione \Vtwo{} si controlla che una regola coprente sopravvissuta eserciti effettivamente il
veto sugli stessi testimoni.

% =======================================================================================
\section{Metodologia statistica}\label{sec:statistics}

Il confronto fra configurazione potata e configurazione di partenza è un confronto appaiato su
variabili che non si distribuiscono in modo normale --- il numero di imputazioni corrette è una
conta con molti valori estremi e molte differenze nulle --- e il test adottato è quindi il test
dei ranghi con segno di Wilcoxon \cite{wilcoxon1945individual}.

\paragraph{Il test.} Per ogni confronto si considerano le differenze $\Delta_i = B_i - A_i$ fra la
variante e la configurazione di partenza sulla medesima configurazione, si scartano le differenze
nulle e si ordinano i valori assoluti. La statistica è la somma dei ranghi delle differenze
positive, $W^+$. Quando non vi sono ranghi medi e il numero di differenze non nulle non supera
venticinque, la distribuzione nulla è calcolata in forma \emph{esatta} per enumerazione; negli
altri casi si impiega l'approssimazione normale con correzione di continuità e correzione per i
ranghi medi. Il test è a due code. Il numero di differenze non nulle, e non il numero di
configurazioni, è la numerosità effettiva del test: le colonne $n$ e $n_{\text{test}}$ della
\autoref{tab:statistics} li riportano entrambi.

\paragraph{Dimensione dell'effetto.} La significatività non dice quanto grande sia l'effetto, e con
numerosità dell'ordine di alcune decine è possibile ottenere valori di $p$ piccoli per differenze
di una sola unità. Si riporta quindi la \emph{differenza mediana} --- robusta rispetto agli
estremi --- e la correlazione rank-biserial per campioni appaiati,
$r = (W^+ - W^-) / (W^+ + W^-)$, che è l'analogo non parametrico della $d$ di Cohen: valori
assoluti intorno a $0.1$, $0.3$ e $0.5$ indicano effetto piccolo, medio e grande.

\paragraph{Intervalli di confidenza.} Per la differenza mediana si riporta un intervallo di
confidenza al \SI{95}{\percent} ottenuto per ricampionamento delle coppie con reimmissione
(diecidimila replicazioni, seme fissato), con il metodo dei percentili. L'intervallo è ampio per
costruzione, data la numerosità, ed è riportato per rendere esplicita l'incertezza invece di
affidare la lettura alla sola soglia di significatività.

\paragraph{Confronti multipli.} Il lavoro verifica sette ipotesi primarie su più metriche e
decine di configurazioni, e la probabilità di almeno un falso positivo cresce con il numero di
test. I test che costituiscono la famiglia primaria --- quelli su cui poggiano le conclusioni del
\autoref{chap:results} --- sono cinque, dichiarati in
\texttt{pruning/stats\_effects.py}, e per essi si riporta il valore di $p$ corretto con la
procedura di Holm \cite{holm1979simple}, che controlla la probabilità di almeno un rifiuto errato
sull'intera famiglia. La correzione è conservativa e va letta come tale: il suo effetto è di
alzare la soglia, non di abbassarla.

\begin{table}[htbp]
\centering
\caption{Test della famiglia primaria. $\Delta$ è la differenza fra variante e configurazione di
partenza, $r$ la correlazione rank-biserial, l'intervallo è quello di confidenza al
\SI{95}{\percent} della differenza mediana. $p$ è il valore grezzo, $p_{\text{Holm}}$ quello
corretto per i cinque test della famiglia.}
\label{tab:statistics}
\scriptsize
\begin{tabular}{@{}lrrrll@{}}
\toprule
confronto & $n$ & $n_{\text{test}}$ & mediana $\Delta$ [IC \SI{95}{\percent}] & $r$ & $p$ \;/\; $p_{\text{Holm}}$ \\
\midrule
\str{A5b} su Bridges, corrette      & 50 & 39 & $+0.5$ $[0.0,\ +2.0]$ & $+0.510$ & $0.0053$ / $0.0159$ \\
\str{A5b} su Bridges, errori        & 50 & 35 & $-1.0$ $[-2.0,\ 0.0]$ & $-0.905$ & $2.5\cdot10^{-6}$ / $1.3\cdot10^{-5}$ \\
taratura conservativa, corrette     & 10 & 0  & $0.0$ $[0.0,\ 0.0]$ & $0.000$ & $1.0000$ / $1.0000$ \\
ordinamento, primo contro originale & 15 & 13 & $+2.0$ $[+1.0,\ +5.0]$ & $+1.000$ & $0.0016$ / $0.0064$ \\
ordinamento, primo contro casuale   & 15 & 8  & $0.0$ $[0.0,\ +0.7]$ & $+0.083$ & $0.8878$ / $1.0000$ \\
\bottomrule
\end{tabular}
\end{table}

\begin{figure}[htbp]
\centering
\begin{tikzpicture}
\begin{axis}[
    width=0.55\textwidth,
    height=0.38\textwidth,
    xmin=-3.4, xmax=6.4,
    ymin=0.4, ymax=5.6,
    ytick={1,2,3,4,5},
    yticklabels={{ordinamento: primo vs casuale ($p$ 0,89)},
                 {ordinamento: primo vs originale ($p$ 0,0016)},
                 {taratura conservativa ($p$ 1,00)},
                 {A5b, errori ($p$ $2{,}5\!\cdot\!10^{-6}$)},
                 {A5b, corrette ($p$ 0,0053)}},
    yticklabel style={font=\scriptsize, align=right, text width=4.3cm},
    xlabel={differenza mediana di imputazioni corrette, con intervallo di confidenza al \SI{95}{\percent}},
    xlabel style={font=\scriptsize},
    grid=major,
    grid style={black!10},
    axis y line=left,
    axis x line=bottom,
    clip=false,
]
\addplot[only marks, mark=*, mark size=2.6pt,
         error bars/.cd, x dir=both, x explicit, error bar style={thick}]
    coordinates {
        (0.0,1) +- (0.35,0)
        (2.0,2) +- (2.0,0)
        (0.0,3) +- (0.0,0)
        (-1.0,4) +- (1.0,0)
        (0.5,5) +- (1.5,0)
    };
\draw[dashed, black!45] (axis cs:0,0.4) -- (axis cs:0,5.6);
\end{axis}
\end{tikzpicture}
\caption{Test della famiglia primaria: differenza mediana di imputazioni corrette e intervallo di
confidenza al \SI{95}{\percent} ottenuto per ricampionamento. I valori di $p$ sono quelli grezzi;
quelli corretti con la procedura di Holm sono nella \autoref{tab:statistics}. La linea verticale
indica l'assenza di differenza.}
\label{fig:forest}
\end{figure}

La \autoref{tab:statistics} si legge in due modi. Sul piano della significatività, i due test
sulla potatura aggressiva e quello sull'ordinamento sopravvivono alla correzione per confronti
multipli; il confronto fra riordino informato e riordino casuale non è significativo, ed è il
risultato che falsifica l'ipotesi H14, non un test fallito. Sul piano della dimensione
dell'effetto, i due test che sopravvivono hanno valori di $r$ medio e grande, mentre gli altri
hanno effetti trascurabili o nulli: la distinzione conta perché una differenza mediana di mezza
imputazione, pur significativa su cinquanta configurazioni, non è un guadagno utilizzabile, ed è
la ragione per cui il \autoref{chap:results} riporta le somme delle differenze e non soltanto i
valori di $p$.

Due limiti della metodologia statistica vanno dichiarati. Il primo è che la famiglia primaria è
stata dichiarata \emph{dopo} la raccolta dei dati, sebbene prima del calcolo dei test: non vi è
quindi pre-registrazione, e la correzione di Holm protegge dalla molteplicità ma non dalla
selezione dei confronti. Il secondo è la potenza: con alcune decine di configurazioni, differenze
piccole rispetto alla variabilità fra dataset non sono distinguibili dal caso, e l'assenza di
significatività non va letta come assenza di effetto.

\section{Setup sperimentale}\label{sec:setup}

\paragraph{Infrastruttura.} La valutazione si fonda su due strumenti sviluppati per il lavoro: un
\emph{pruner} esterno, che riscrive il file delle regole secondo la strategia richiesta, e un
\emph{benchmark} che esegue il sistema nella configurazione di partenza e in quella potata,
raccoglie le metriche e le confronta. Il sistema oggetto di studio non è mai stato modificato
durante la valutazione esterna: l'eseguibile impiegato è quello originale, verificato per
checksum all'inizio e alla fine delle campagne.

\paragraph{Isolamento delle esecuzioni.} Ogni esecuzione avviene in una cartella di lavoro
dedicata, che viene ripulita prima dell'esecuzione successiva. Le esecuzioni sono condotte in
sequenza, mai in parallelo: la cartella di lavoro è condivisa e due esecuzioni simultanee
sovrascriverebbero reciprocamente gli output. La circostanza è stata verificata sperimentalmente
e ha richiesto di rieseguire alcune configurazioni.

\paragraph{Riproducibilità.}\label{sec:reproducibility}
Tutti i numeri presentati nel \autoref{chap:results} sono ricalcolabili dai file di risultati
prodotti dalle campagne, e la ricalcolabilità è verificata automaticamente: una procedura di
verifica ricalcola i valori asseriti nella documentazione e li confronta con quelli memorizzati,
segnalando ogni divergenza. La procedura copre \num{67} asserzioni numeriche, incluse quelle
relative alle condizioni teoriche. I dettagli sono nell'\autoref{app:reproducibility}.

% =======================================================================================
\section{Dataset e configurazioni}\label{sec:datasets-exp}

La valutazione impiega tre famiglie di dataset, scelte fra quelle disponibili nell'archivio
secondo criteri di utilizzabilità e di diversità strutturale
(\autoref{tab:families}). La \autoref{tab:grid} riassume la griglia effettivamente coperta.

\begin{table}[htbp]
\centering
\caption{Griglia sperimentale. ``Configurazioni'' è il numero di triplette
\emph{dataset, run, soglia} distinte.}
\label{tab:grid}
\footnotesize
\begin{tabular}{lrrrr}
\toprule
famiglia & dataset & run & soglie & configurazioni \\
\midrule
Bridges    & 5 & 1--5 & 3, 9, 15 & 60 \\
Glass      & 5 & 1--5 & 3, 9, 15 & 75 \\
Restaurant & 4 & 1--3 & 3, 6     & 22 \\
\midrule
totale     & 14 & & & \textbf{157} \\
\bottomrule
\end{tabular}
\end{table}

Le tre famiglie differiscono per caratteristiche che si riveleranno determinanti
nell'interpretazione dei risultati: il numero di regole disponibili, la distribuzione delle celle
nulle fra tuple con una sola cella nulla e tuple con più celle nulle
(\autoref{sec:two-branches}), e la quota di imputazioni che il sistema ottiene per via esatta
(\autoref{sec:mechanism}).

Due famiglie dell'archivio non sono state impiegate per l'intera griglia. La famiglia Nursery
comprende cinque file di regole --- le soglie \num{3}, \num{6}, \num{9}, \num{12} e \num{15}
--- ma nessun dataset e nessun file di tuple iniziali, quindi nessuna configurazione eseguibile:
è il residuo di una fase di estrazione di cui non è stato distribuito il dato. La famiglia
Physician comprende venticinque configurazioni piene di \num{103592} righe, di dimensioni tali da
rendere impraticabile una campagna, e quattro versioni remodulate il cui materiale di ground truth
presenta indici di riga eccedenti il dataset (\autoref{chap:defects}). Le versioni remodulate sono
state rese eseguibili escludendo le celle non localizzabili, ed è su di esse che si fonda la prova
sulla quarta famiglia i cui esiti sono riportati nella \autoref{sec:physician}.

% =======================================================================================
\section{Metriche}\label{sec:metrics}

Le metriche adottate sono le seguenti. Sia $C$ l'insieme delle celle per cui è disponibile il
valore di ground truth; si ricordi che tale insieme non coincide necessariamente con l'insieme
delle celle nulle del dataset (\autoref{sec:datasets}).

\begin{table}[htbp]
\centering
\caption{Metriche adottate. $C$ è l'insieme delle celle con valore di ground truth disponibile,
$I \subseteq C$ l'insieme delle celle imputate, $K \subseteq I$ quelle imputate correttamente.}
\label{tab:metrics}
\footnotesize
\begin{tabular}{@{}llp{4.4cm}@{}}
\toprule
metrica & definizione & che cosa falsifica \\
\midrule
copertura & $|I| / |C|$ & una potatura che, togliendo regole, fa rifiutare al sistema celle che
prima valorizzava \\
precisione & $|K| / |I|$ & una potatura che rende accettate imputazioni di qualità peggiore \\
accuratezza & $|K| / |C|$ & l'effetto complessivo; è il prodotto delle due precedenti \\
errori & $|I| - |K|$ & l'aumento del danno, che l'accuratezza può mascherare quando la copertura
cresce \\
quota esatta & frazione di $I$ ottenuta per via esatta & la composizione delle imputazioni, primo
fattore del meccanismo (\autoref{sec:mechanism}) \\
preprocessing & secondi per riscrivere il file delle regole & il costo dell'analisi, da
confrontare con il risparmio in esecuzione \\
esecuzione & secondi del sistema, a parità di configurazione & l'effetto sul costo, che
non segue la riduzione (\autoref{sec:res-cost}) \\
\bottomrule
\end{tabular}
\end{table}

La distinzione fra copertura e precisione è essenziale nel seguito. Un metodo basato su vincoli
può rifiutare di imputare, e una potatura che rimuova regole può indurre nuovi rifiuti: in quel
caso la copertura cala mentre la precisione può salire, e l'accuratezza può restare invariata. La
lettura congiunta delle due metriche è quindi necessaria, e le analisi che seguono la rispettano
sistematicamente.

Per i dataset a prevalenza testuale (Restaurant) il confronto esatto fra stringhe è troppo
severo: due ragioni sociali che differiscono soltanto per la forma societaria finale sono contate
come valori diversi pur designando la stessa entità. Per questa famiglia si adotta pertanto una
soglia di similarità normalizzata di Levenshtein pari a \num{0.8}, e le metriche testuali sono
riportate separatamente da quelle esatte.

% =======================================================================================
% =======================================================================================
\section{Termine di paragone banale}\label{sec:trivial-baseline}

Le metriche del capitolo sono differenze rispetto alla medesima configurazione non potata: dicono
se la potatura migliora o peggiora il risultato del sistema, non se il risultato del sistema sia
buono in assoluto. Per dare una scala, la \autoref{tab:trivial} confronta l'accuratezza del
sistema con quella di un imputatore banale che sostituisce a ogni cella la \emph{moda} della
colonna, calcolata sui valori osservati. La moda è un termine di paragone volutamente povero ---
non impiega alcuna dipendenza --- ed è calcolata sulle stesse celle valutate dal sistema.

\begin{table}[htbp]
\centering
\caption{Confronto con un imputatore banale (moda di colonna) sulle stesse celle valutate. La
copertura della moda è totale per costruzione. I valori si riferiscono alle configurazioni con
schema allineato: \num{6} configurazioni di Bridges (prima ripetizione, soglie 3 e 9) e le
\num{2} versioni eseguibili di Physician.}
\label{tab:trivial}
\footnotesize
\begin{tabular}{@{}lrrrrr@{}}
\toprule
famiglia & configurazioni & celle & accuratezza moda & accuratezza sistema & copertura sistema \\
\midrule
Bridges   & \num{6} & \num{224} & \num{0.536} & \num{0.545} & \SI{61.6}{\percent} \\
Physician & \num{2} & \num{40}  & \num{0.275} & \num{0.150} & \SI{35.0}{\percent} \\
\bottomrule
\end{tabular}
\end{table}

La lettura del confronto è duplice, e va data per intero. Su Bridges la moda raggiunge
un'accuratezza praticamente uguale a quella del sistema (\num{0.536} contro \num{0.545}) con
copertura totale: su questo dataset la qualità \emph{assoluta} delle imputazioni non è il punto di
forza del sistema, e il valore che il lavoro misura è quello che la potatura aggiunge o toglie a
parità di sistema. Su Physician il divario è più netto e in direzione opposta: la moda è migliore
del sistema (\num{0.275} contro \num{0.150}), perché il sistema copre il \SI{35}{\percent} delle
celle e, fra quelle che copre, non sempre indovina.

Due avvertenze delimitano la portata del confronto. La prima è che il perimetro della tesi non
comprende il confronto con metodi di imputazione alternativi (\autoref{sec:res-limits}): la moda
non è un concorrente del sistema ma una \emph{scala} per i suoi numeri. La seconda è che
l'accuratezza è misurata soltanto sulle celle introdotte artificialmente, quelle per cui esiste il
ground truth: il sistema può imputare anche celle nulle preesistenti, che non entrano in questa
misura (\autoref{sec:datasets}).

\section{Esperimenti condotti}\label{sec:experiments}

Le campagne condotte ammontano complessivamente a \num{1618} esecuzioni del sistema in
configurazione potata, oltre alle esecuzioni di riferimento. Esse non sono un insieme di prove
giustapposte: ciascuna risponde a una delle ipotesi enunciate nel \autoref{chap:hypotheses}, e la
\autoref{tab:experiments} esplicita il collegamento, perché è su di esso che poggia la lettura
degli esiti.

\begin{table}[htbp]
\centering
\caption{Esperimenti condotti, obiettivo e ipotesi verificate. La colonna ``bracci'' indica che
cosa viene confrontato; la colonna ``risultati'' la sezione in cui l'esito è riportato.}
\label{tab:experiments}
\footnotesize
\begin{tabular}{@{}clp{4.4cm}ll@{}}
\toprule
\# & obiettivo & bracci confrontati & ipotesi & risultati \\
\midrule
E1 & equivalenza della potatura garantita & \str{veto-cover} contro configurazione di partenza, su tre famiglie & H6 & \ref{sec:res-vetocover} \\
E2 & verifica indipendente della copertura & ricostruzione dei testimoni del veto per ogni regola rimossa & H7 & \ref{sec:res-vetocover} \\
E3 & effetto della potatura aggressiva & \str{A5b} contro configurazione di partenza, su tre famiglie e tre soglie & H8, H11 & \ref{sec:res-a5b}, \ref{sec:res-boundary} \\
E4 & attribuzione del guadagno & riordino per precisione, riordino inverso, tre permutazioni casuali & H14 & \ref{sec:negative-results} \\
E5 & ablazione dei componenti & ciascun componente di \str{A5b} da solo, la combinazione, la somma & H8 & \ref{sec:ablation} \\
E6 & proporzionalità fra riduzione e tempo & tempi di esecuzione appaiati, per strategia e configurazione & H16 & \ref{sec:res-cost} \\
E7 & strategie esatte alternative & \str{score-argmin} con tre politiche sui veti, su Bridges e Glass & H1, H2 & \ref{sec:score-argmin} \\
\bottomrule
\end{tabular}
\end{table}

Due precisazioni completano il disegno. La prima riguarda le configurazioni \emph{vacue}: le
combinazioni in cui la strategia non rimuove alcuna regola producono un file identico a quello di
partenza, e sono escluse dai test e conteggiate a parte. La seconda riguarda l'ordine di
esecuzione: le campagne sono state condotte in modo strettamente sequenziale, perché la cartella
di lavoro del sistema è condivisa fra esecuzioni e due esecuzioni simultanee sovrascrivono
reciprocamente gli output senza segnalare alcun errore. La circostanza è stata verificata
sperimentalmente, ed è la ragione per cui non si riportano misure di tempi concorrenti.

\section{Studio di ablazione}\label{sec:ablation}

La strategia \str{A5b} è una combinazione di due interventi: un riordino delle regole a soglia di
destra nulla e una potatura adattiva. Per stabilire quale dei due produca il guadagno, entrambi
sono stati misurati anche separatamente, sulle stesse configurazioni.

\begin{table}[htbp]
\centering
\caption{Ablazione di \str{A5b} su Bridges, configurazioni comuni. Le variazioni sono rispetto
alla medesima configurazione non potata.}
\label{tab:ablation}
\footnotesize
\begin{tabular}{lrrrr}
\toprule
componente & riduzione RFD & $\Delta$ corrette & $\Delta$ errori & esito \\
\midrule
\str{exact-priority} (solo)     & \SI{0}{\percent}   & $+25$ & $-53$ & non stabile fra soglie \\
\str{dom-sub-adaptive} (solo)   & \SI{81}{\percent}--\SI{87}{\percent} & $-20$ & $+20$ & danneggia \\
\textbf{\str{A5b} (combinazione)} & \SI{81}{\percent}--\SI{87}{\percent} & $\mathbf{+45}$ & $\mathbf{-32}$ & \textbf{migliora} \\
\midrule
somma dei due componenti & & $+5$ & & \\
\textbf{sinergia} & & $\mathbf{+40}$ & & \\
\bottomrule
\end{tabular}
\end{table}

Il risultato dell'ablazione è che \emph{nessuno dei due componenti funziona da solo}: il riordino
migliora il risultato a una soglia e lo peggiora all'altra, mentre la potatura adattiva lo
peggiora sistematicamente. La combinazione è invece stabile e positiva. La differenza fra il
guadagno osservato ($+45$) e la somma dei due componenti isolati ($+5$) è pari a $+40$, e
costituisce una \emph{sinergia} che va spiegata, non registrata.

La spiegazione è che i due interventi agiscono su fattori diversi del medesimo prodotto. La
potatura adattiva modifica la \emph{composizione} dei candidati --- quali regole concorrono a
determinare il valore --- e quindi la quota di imputazioni ottenute per via esatta; il riordino
modifica la \emph{sequenza} con cui i cluster vengono tentati e, nella procedura a più celle nulle,
determina quale cluster scrive per ultimo, quindi quante celle vengono effettivamente valorizzate.
Presi singolarmente, il primo peggiora la qualità delle imputazioni senza aumentarne il numero, il
secondo aumenta il numero di imputazioni senza migliorarne la composizione; applicati insieme,
producono un aumento di entrambi i fattori, e il guadagno risultante supera la somma dei guadagni
isolati. Il \autoref{sec:mechanism} verifica questa lettura sui tre regimi osservati --- Bridges,
Glass e Restaurant --- e ne ricava il criterio diagnostico.

Va sottolineato un limite del disegno di ablazione: le configurazioni comuni ai tre bracci sono
venti, e la lettura della sinergia poggia quindi su un campione più piccolo di quello della
campagna principale. L'ablazione identifica il meccanismo, non ne stabilisce la generalità.

% =======================================================================================


    \chapter{Analisi dei risultati}\label{chap:results}


\noindent Questo capitolo presenta e interpreta i risultati. L'ordine non è per metrica ma per
\emph{caso}: la potatura che non perde nulla, il confronto con il criterio di ridondanza della
teoria classica, la potatura che produce il guadagno e il confine oltre il quale il guadagno non
si realizza, e il meccanismo che rende conto dei casi osservati. Da quel meccanismo discende un
criterio diagnostico, di cui il capitolo tenta la verifica su una famiglia indipendente; seguono i
risultati negativi e l'analisi del costo, che mostra perché la riduzione delle regole non basti a
garantire un risparmio di tempo. La lettura d'insieme, gli esiti delle ipotesi e i limiti
dell'evidenza sono oggetto del capitolo di discussione.

% =======================================================================================
\section{La strategia che non perde nulla}\label{sec:res-vetocover}

La condizione \Vtwo{} è l'unico criterio di potatura che offra una garanzia. La sua applicazione
è stata verificata su \num{113} configurazioni --- \num{111} distribuite su tre famiglie e
\num{2} su una quarta famiglia indipendente (\autoref{sec:physician}) --- con la condizione
che l'output del sistema sia \emph{identico} a quello della configurazione di partenza --- non
soltanto equivalente nelle metriche aggregate, ma identico cella per cella.

\begin{table}[htbp]
\centering
\caption{Applicazione della condizione \Vtwo{} su tre famiglie. ``Identiche'' è il numero di
configurazioni in cui corrette, errori e copertura coincidono esattamente con la configurazione
di partenza.}
\label{tab:res-vetocover}
\footnotesize
\begin{tabular}{llrrrrr}
\toprule
famiglia & soglia & $n$ & $N$ & riduzione & $\Delta$ corrette & identiche \\
\midrule
Bridges    & 3  & 10 & \num{1830}  & \SI{11.7}{\percent} & 0 & \textbf{10/10} \\
Bridges    & 9  & 10 & \num{8961}  & \SI{17.1}{\percent} & 0 & \textbf{10/10} \\
Bridges    & 15 & 10 & \num{18844} & \SI{20.8}{\percent} & 0 & \textbf{10/10} \\
Glass      & 3  & 25 & \num{1511}  & \SI{2.6}{\percent}  & 0 & \textbf{25/25} \\
Glass      & 9  & 25 & \num{11648} & \SI{7.4}{\percent}  & 0 & \textbf{25/25} \\
Glass      & 15 & 25 & \num{27445} & \SI{8.6}{\percent}  & 0 & \textbf{25/25} \\
Restaurant & 3  &  8 & 25          & \SI{0.0}{\percent}  & 0 & \textbf{8/8} \\
Restaurant & 6  &  8 & 124         & \SI{1.6}{\percent}  & 0 & \textbf{8/8} \\
\midrule
\textbf{totale} & & \textbf{113} & & & \textbf{0} & \textbf{113/113} \\
\bottomrule
\end{tabular}
\end{table}

Il risultato va letto nei due sensi. Nel senso positivo, è la conferma che la condizione \Vtwo{}
è corretta: una potatura \emph{esatta} è possibile, e la sua applicazione non altera di una sola
cella il comportamento del sistema. Nel senso negativo, la \autoref{tab:res-vetocover} mostra
anche il prezzo: la riduzione è nulla su Restaurant a soglia 3, inferiore al \SI{3}{\percent} su
Glass a soglia 3, e non supera il \SI{20.8}{\percent} in alcun caso.

La verifica è stata condotta anche in forma indipendente, controllando su stati reali del dataset
--- e non soltanto per implicazione logica --- che per ogni regola rimossa esistesse una regola
coprente sopravvissuta in grado di esercitare il veto sugli stessi testimoni. La verifica ha
coperto \num{1242321} terne, con zero discrepanze.

% =======================================================================================
% =======================================================================================
\section{La riduzione classica, e perché non è sicura}\label{sec:res-classical}

Il termine di paragone più ovvio per una potatura è la riduzione di ridondanza della teoria
classica delle dipendenze funzionali: una regola $X \to A$ è superflua se un'altra regola
$Y \to A$ ha $Y$ sottoinsieme proprio di $X$, perché $X \to A$ si ottiene per augmentazione
(\autoref{sec:rule-reduction}). Il criterio è sintattico e ignora le soglie; è stato implementato
come strategia (\str{classical-cover}) e misurato sulle medesime configurazioni di Bridges a
soglia 9, con la medesima ricetta di accoppiamento delle altre campagne.

\begin{table}[htbp]
\centering
\caption{Riduzione classica (\str{classical-cover}) su Bridges a soglia 9, prima ripetizione:
l'insieme passa da \num{8961} a \num{230} regole in tutte e cinque le configurazioni
(\SI{97.4}{\percent}). Ogni riga confronta la configurazione di partenza con quella potata.}
\label{tab:res-classical}
\footnotesize
\begin{tabular}{@{}lrrrrr@{}}
\toprule
configurazione & copertura & accuratezza & corrette & errori & tempo (\si{\second}) \\
\midrule
\texttt{bridges\_14\_1} & $0.714 \to 0.571$ & $0.643 \to 0.500$ & $9 \to 7$   & $1 \to 3$   & $9.9 \to 0.3$ \\
\texttt{bridges\_28\_1} & $0.714 \to 0.643$ & $0.643 \to 0.429$ & $18 \to 12$ & $2 \to 6$   & $32.0 \to 0.6$ \\
\texttt{bridges\_42\_1} & $0.595 \to 0.595$ & $0.524 \to 0.429$ & $22 \to 18$ & $3 \to 7$   & $36.9 \to 0.7$ \\
\texttt{bridges\_56\_1} & $0.679 \to 0.571$ & $0.589 \to 0.321$ & $33 \to 18$ & $4 \to 10$  & $56.9 \to 1.0$ \\
\texttt{bridges\_70\_1} & $0.529 \to 0.500$ & $0.457 \to 0.357$ & $32 \to 25$ & $6 \to 12$  & $54.9 \to 1.0$ \\
\midrule
totale & & & \textbf{114} $\to$ \textbf{80} & \textbf{16} $\to$ \textbf{38} & \\
\bottomrule
\end{tabular}
\end{table}

Il risultato va nella direzione opposta a quella che il solo numero di regole rimosse
suggerirebbe. La riduzione è del \SI{97.4}{\percent} --- da \num{8961} a \num{230} regole, molto
oltre il \SI{20.9}{\percent} della classe garantita e oltre l'\SI{84}{\percent} della strategia
aggressiva --- e il tempo di esecuzione cala di oltre cinquanta volte; ma l'output cambia in
\emph{tutte} le configurazioni, la copertura scende, l'accuratezza mediana perde \num{14.3} punti,
le imputazioni corrette passano da \num{114} a \num{80} e gli errori da \num{16} a \num{38}. La
circostanza va letta con la numerosità che ha: con cinque configurazioni il test esatto dà
$p = \num{0.0625}$, e nessuna delle cinque migliora, ma la soglia convenzionale di significatività
non è raggiunta. Ciò che il confronto stabilisce non è una differenza statisticamente significativa
rispetto al criterio classico, ma che il criterio classico \emph{non è innocuo}: la sua applicazione
altera il comportamento del sistema in ogni caso osservato.

La spiegazione è quella argomentata nella \autoref{sec:rule-reduction}. Il criterio classico
rimuove regole che il sistema usava, perché la relazione di sussunzione che lo fonda non tiene
conto né della \emph{direzione} delle soglie sul lato destro né dell'\emph{ampiezza} di quelle sul
lato sinistro: la regola dichiarata superflua può generare candidati, e la regola che la sussume
può essere più severa e quindi non implicarne il veto. La condizione \Vtwo{} è l'analogo operativo
di quella relazione, ma richiede soglie più larghe a sinistra e non più larghe a destra, cioè la
direzione opposta.

La conseguenza per la lettura del lavoro è duplice. Sul piano teorico, il confronto è la verifica
empirica della distinzione fra ridondanza dichiarativa e sicurezza operativa. Sul piano della
valutazione, mostra perché la riduzione non possa essere l'unico indicatore di una strategia: un
criterio che rimuove il \SI{97}{\percent} delle regole peggiorando il risultato è, per gli
obiettivi di questo lavoro, peggiore di uno che ne rimuove il \SI{17}{\percent} senza toccarlo.

\section{La strategia che produce il guadagno}\label{sec:res-a5b}

La combinazione \str{A5b}, che abbina un riordino delle regole a soglia di destra nulla a una
potatura adattiva, è stata valutata sull'intera griglia di Bridges: cinque livelli di
incompletezza, cinque esecuzioni indipendenti, due soglie.

\begin{table}[htbp]
\centering
\caption{\str{A5b} su Bridges, griglia completa. Confronto appaiato con la medesima
configurazione non potata.}
\label{tab:res-a5b}
\footnotesize
\begin{tabular}{lrrrrrr}
\toprule
soglia & $n$ & riduzione RFD & $\Delta$ corrette & $\Delta$ errori & $\Delta$ precisione & accelerazione \\
\midrule
9  & 25 & \SI{81.2}{\percent} & $+32$ & $-40$ & $+0.050$ & \num{9.0}$\times$ \\
15 & 25 & \SI{87.2}{\percent} & $+37$ & $-32$ & $+0.038$ & \num{17.5}$\times$ \\
\midrule
\textbf{totale} & \textbf{50} & \textbf{\SI{84.2}{\percent}} & $\mathbf{+69}$ & $\mathbf{-72}$ & $+0.044$ & \num{11.0}$\times$ \\
\bottomrule
\end{tabular}
\end{table}

Il test di Wilcoxon signed-rank sulle corrette dà, sull'insieme delle \num{50} configurazioni,
$n = 39$, $W^+ = 589$, $W^- = 191$, $p = \num{0.0053}$; sugli errori, $p = 2.5\cdot 10^{-6}$. Separando
per soglia il quadro è più sfumato e va riportato: a soglia 9 il test dà $p = \num{0.0562}$
($n = 18$) e a soglia 15 $p = \num{0.0484}$ ($n = 21$), sicché \emph{nessuna delle due soglie da
sola} è conclusiva al livello del \SI{5}{\percent}, mentre l'insieme delle configurazioni lo è
con ampio margine. La ragione è che in una parte delle configurazioni la potatura è neutra, e il
test vede soltanto le differenze non nulle.

Il risultato va confrontato con quello della \autoref{sec:res-vetocover}. La strategia garantita
non altera nulla e riduce poco; \str{A5b} riduce molto e \emph{migliora} il risultato. Le due
strategie non sono alternative sullo stesso asse: rispondono a due domande diverse.

% =======================================================================================
\section{Il confine: dove il guadagno non si realizza}\label{sec:res-boundary}

La medesima strategia, applicata a Glass e Restaurant con la stessa taratura, produce un
danneggiamento sistematico.

\begin{table}[htbp]
\centering
\caption{Effetto della potatura aggressiva su tutte le famiglie. La taratura è la medesima
($\alpha = 0.85$).}
\label{tab:res-boundary}
\footnotesize
\begin{tabular}{llrrrr}
\toprule
famiglia & $n$ & riduzione RFD & $\Delta$ corrette & $\Delta$ copertura & esito \\
\midrule
Bridges    & 50 & \SI{84.2}{\percent} & $\mathbf{+69}$ & $-0.003$ & \textbf{migliora} \\
Glass      & 45 & \SI{82.8}{\percent} & $-54$          & $-0.034$ & danneggia \\
Restaurant & 18 & \SI{44.2}{\percent} & $-12$          & $-0.005$ & danneggia \\
\bottomrule
\end{tabular}
\end{table}

La causa è precisata dal segno della variazione di copertura: su Glass la potatura riduce le
imputazioni effettuate di circa un terzo, e il danno è quindi una \emph{perdita di copertura},
non un peggioramento della qualità di ciò che viene imputato. Su Bridges, invece, la copertura
non cala.

Una taratura più conservativa della medesima strategia --- ottenuta alzando la soglia di
stabilità del parametro di potatura --- elimina il danneggiamento su Glass e Restaurant, al
prezzo di una riduzione molto inferiore:

\begin{table}[htbp]
\centering
\caption{La medesima strategia con taratura conservativa ($\alpha = 0.99$) su Glass e Restaurant,
griglia completa. ``Identiche'' indica le configurazioni in cui corrette, errori e copertura
coincidono con la configurazione di partenza.}
\label{tab:res-gentle}
\footnotesize
\begin{tabular}{lrrrrr}
\toprule
famiglia & config. & identiche & $\Delta$ corr. & $\Delta$ err. & riduz.\ (mediana / massima) \\
\midrule
Glass      & 75 & \textbf{74/75} & 0   & $-1$ & \SI{0.1}{\percent} / \SI{52.9}{\percent} \\
Restaurant & 22 & \textbf{20/22} & $-1$ & $-2$ & \SI{0.0}{\percent} / \SI{45.2}{\percent} \\
\midrule
\textbf{totale} & \textbf{97} & \textbf{94/97} & $\mathbf{-1}$ & $\mathbf{-3}$ & \SI{0.0}{\percent} / \SI{52.9}{\percent} \\
\bottomrule
\end{tabular}
\end{table}

Il confronto fra la \autoref{tab:res-boundary} e la \autoref{tab:res-gentle} è istruttivo: la
stessa strategia passa da $-54$ a $-1$ imputazioni corrette su Glass al variare di un solo
parametro. La spiegazione di questa sensibilità è l'oggetto della sezione successiva.

Un elemento della \autoref{tab:res-gentle} va sottolineato perché limita la portata del
risultato: la riduzione \emph{mediana} è nulla. Su più della metà delle configurazioni la
strategia conservativa non rimuove alcuna regola. Il beneficio si concentra in \num{27}
configurazioni su \num{97}, dove è sostanzioso. Ciò che la taratura conservativa garantisce su
tutte le configurazioni non è la riduzione, ma l'assenza di perdite rilevanti.

% =======================================================================================
\section{Il meccanismo}\label{sec:mechanism}

\begin{figure}[htbp]
\centering
\begin{tikzpicture}
\begin{axis}[
    width=0.82\textwidth,
    height=0.46\textwidth,
    xlabel={quota di imputazioni esatte nella configurazione di partenza (\si{\percent})},
    ylabel={variazione di imputazioni corrette},
    xmin=0, xmax=112, ymin=-75, ymax=95,
    grid=major,
    grid style={black!12},
    axis lines=left,
    ytick={-60,-40,-20,0,20,40,60,80},
    xtick={0,20,40,60,80,100},
    legend style={font=\footnotesize, at={(0.02,0.97)}, anchor=north west, draw=none, fill=none},
]
\addplot[only marks, mark=*, mark size=3.6pt, color=black] coordinates {(100,-54)};
\addplot[only marks, mark=square*, mark size=3.0pt, color=black!60] coordinates {(19.6,69)};
\addplot[only marks, mark=triangle*, mark size=2.6pt, color=black!35] coordinates {(23.2,-12)};
\addlegendentry{Glass (821 imputazioni)}
\addlegendentry{Bridges (504)}
\addlegendentry{Restaurant (241)}
\draw[dashed, black!45] (axis cs:0,0) -- (axis cs:112,0);
\node[font=\scriptsize, anchor=west] at (axis cs:24,73) {Bridges};
\node[font=\scriptsize, anchor=west] at (axis cs:26,-18) {Restaurant};
\node[font=\scriptsize, anchor=east] at (axis cs:97,-49) {Glass};
\end{axis}
\end{tikzpicture}
\caption{Il meccanismo su tre famiglie. In ascissa la quota di imputazioni ottenute per via esatta
nella configurazione di partenza, in ordinata la variazione di imputazioni corrette prodotta dalla
potatura aggressiva; la dimensione del marcatore è proporzionale al numero di imputazioni della
configurazione di partenza. Il guadagno richiede che la quota esatta non sia satura: dove lo è
(Glass) la potatura può soltanto togliere copertura.}
\label{fig:mechanism}
\end{figure}

Le sezioni precedenti hanno stabilito \emph{che cosa} accade. Questa sezione stabilisce
\emph{perché}, e il risultato è il contributo interpretativo principale del lavoro.

Il sistema registra, per ogni imputazione effettuata, se essa sia stata ottenuta per via
\emph{esatta} --- la scorciatoia che si attiva quando la regola impiegata ha soglia di destra
nulla --- oppure per via approssimata. La distinzione era disponibile nei dati di tutte le
campagne, ma non era stata impiegata nell'interpretazione. La
\autoref{tab:mechanism} la mette in relazione con il numero di imputazioni.

\begin{table}[htbp]
\centering
\caption{Il meccanismo del guadagno. Per ciascuna famiglia si confrontano il numero di
imputazioni, la quota di quelle ottenute per via esatta e il numero di imputazioni corrette,
prima e dopo la potatura aggressiva. I valori sono aggregati sulla \emph{campagna completa} di
ciascuna famiglia: \num{50} configurazioni per Bridges, \num{45} per Glass, \num{18} per
Restaurant, alla taratura $\alpha = 0{,}85$.}
\label{tab:mechanism}
\footnotesize
\begin{tabular}{lrrrrr}
\toprule
& \multicolumn{2}{c}{\textbf{configurazione di partenza}} & \multicolumn{2}{c}{\textbf{dopo \str{A5b}}} & \\
\cmidrule(lr){2-3}\cmidrule(lr){4-5}
famiglia & imputazioni & quota esatta & imputazioni & quota esatta & $\Delta$ corrette \\
\midrule
Bridges (\num{50})    & \num{1309} & \SI{21.8}{\percent} & \num{1306} & \textbf{\SI{40.4}{\percent}} & $\mathbf{+69}$ \\
Glass (\num{45})      & 821        & \SI{100.0}{\percent} & 721       & \SI{100.0}{\percent}        & $-54$ \\
Restaurant (\num{18}) & 241        & \SI{23.2}{\percent}  & 229       & \textbf{\SI{49.8}{\percent}} & $-12$ \\
\bottomrule
\end{tabular}
\end{table}

Il numero di imputazioni corrette è il prodotto di due fattori: il \emph{numero} di imputazioni
effettuate e la \emph{quota} di quelle corrette, che a sua volta dipende dalla composizione dei
candidati e in particolare dalla quota di imputazioni ottenute per via esatta. La potatura
produce un guadagno soltanto se la composizione ha spazio per migliorare \emph{e} il numero di
imputazioni non diminuisce abbastanza da compensare il miglioramento. La
\autoref{tab:mechanism} mostra che i tre casi si distinguono esattamente su questo:

\begin{itemize}[leftmargin=1.4em]
    \item \textbf{Bridges}: la quota esatta cresce dal \SI{21.8}{\percent} al \SI{40.4}{\percent}
    e il numero di imputazioni resta sostanzialmente invariato (\num{1309} contro \num{1306}). Il
    miglioramento della composizione si trasferisce integralmente al risultato: \num{69}
    imputazioni corrette in più.
    \item \textbf{Glass}: la quota esatta è \emph{già} pari al \SI{100}{\percent}, perché ogni
    imputazione passa per la via esatta. Non vi è spazio per migliorare la composizione, e la
    potatura può soltanto togliere copertura: le imputazioni scendono da \num{821} a \num{721} e
    le corrette da \num{517} a \num{463}.
    \item \textbf{Restaurant}: la composizione ha spazio e migliora più che su Bridges (dal
    \SI{23.2}{\percent} al \SI{49.8}{\percent}), ma il numero di imputazioni cala da \num{241} a
    \num{229} e le imputazioni guadagnate in composizione sono meno di quelle perse in copertura:
    il saldo è di \num{12} imputazioni corrette in meno.
\end{itemize}

Sulle \num{20} configurazioni comuni all'ablazione, dove entrambi i bracci sono misurati sulle
stesse configurazioni, i due fattori crescono entrambi e la variazione è di $+45$
(\autoref{tab:ablation}): la differenza rispetto al $+69$ della campagna completa è una
differenza di campione, non di strategia, ed è la ragione per cui le due tabelle dichiarano
ciascuna la propria base.

Questo risultato spiega anche un'osservazione che era rimasta senza interpretazione: il riordino
\str{exact-priority} è risultato un \emph{no-op esatto} su Glass, con output identico in
\num{45} configurazioni su \num{45}. La spiegazione è immediata una volta nota la quota esatta:
se tutte le imputazioni passano già per la via esatta, non vi è alcuna regola da promuovere.

\paragraph{Perché la sinergia non è misteriosa.} La \autoref{tab:ablation} --- che si riferisce
alle \num{20} configurazioni comuni ai quattro bracci --- mostra che il guadagno della
combinazione ($+45$) supera di molto la somma dei componenti isolati ($+5$). Il meccanismo rende
conto della differenza: i due componenti agiscono su fattori \emph{diversi}. Il riordino per
priorità di precisione dispone per prime le regole che producono valori esatti, e migliora quindi
la composizione dei candidati; la potatura adattiva modifica quali regole concorrono alla
generazione, e quindi quante celle trovano un candidato accettabile. Ciascun componente vale
tanto più quanto più è alto il fattore su cui l'altro agisce, e la combinazione è per questo
super-additiva: è la stessa legge dei due fattori, letta sul campione in cui entrambi sono
misurati.

% =======================================================================================
\section{Il criterio diagnostico}\label{sec:diagnostic}

\begin{figure}[htbp]
\centering
\begin{tikzpicture}[
    font=\footnotesize,
    node distance=7mm,
    st/.style={draw, rounded corners=1pt, align=center, inner sep=4pt, minimum height=8mm},
    dec/.style={draw, diamond, aspect=2.6, align=center, inner sep=1pt, font=\scriptsize},
    ar/.style={-{Latex[length=2mm]}, thick}
]
\node[st] (m1) {misura la quota esatta\\della configurazione di partenza};
\node[dec, below=of m1] (d1) {quota $\approx \SI{100}{\percent}$?};
\node[st, right=14mm of d1] (no1) {\textbf{non potare}\\la leva non esiste};
\node[st, below=of d1] (m2) {applica la strategia e misura\\la copertura su un sottoinsieme};
\node[dec, below=of m2] (d2) {la copertura cala?};
\node[st, right=14mm of d2] (no2) {taratura conservativa\\o nessuna potatura};
\node[st, below=of d2, fill=black!10] (ok) {\textbf{adotta la potatura}};
\draw[ar] (m1) -- (d1);
\draw[ar] (d1) -- node[above, font=\scriptsize] {sì} (no1);
\draw[ar] (d1) -- node[right, font=\scriptsize] {no} (m2);
\draw[ar] (m2) -- (d2);
\draw[ar] (d2) -- node[above, font=\scriptsize] {sì} (no2);
\draw[ar] (d2) -- node[right, font=\scriptsize] {no} (ok);
\end{tikzpicture}
\caption{Procedura del criterio diagnostico. Il primo passo non richiede esecuzioni, perché la
quota esatta è registrata dal sistema; il secondo richiede di misurare la copertura, ed è il solo
costo del criterio.}
\label{fig:criterion}
\end{figure}

Dal meccanismo discende un criterio operativo, che sostituisce la soglia euristica adottata in
precedenza. Il criterio si fonda su una grandezza che il sistema \emph{registra già}: la quota di
imputazioni ottenute per via esatta nella configurazione di partenza, disponibile nei file di
risultato prima di qualunque potatura.

\begin{description}[leftmargin=0em]
    \item[Quota esatta $\approx \SI{100}{\percent}$.] La leva della qualità non esiste: non vi è
    margine per sostituire imputazioni approssimate con imputazioni esatte, e la potatura può
    soltanto togliere copertura. \emph{Non potare}. È il caso di Glass.
    \item[Quota esatta bassa.] La leva esiste, e la potatura conviene \emph{a condizione che la
    copertura non cali}. È il caso di Bridges, dove la copertura cresce, e di Restaurant, dove
    cala.
\end{description}

\paragraph{La procedura.} Il criterio si traduce in quattro passi, eseguibili senza modificare né
il sistema né il dataset.

\begin{enumerate}[leftmargin=1.6em]
    \item \textbf{Misurare la quota esatta} della configurazione di partenza, leggendola dai file
    di risultato che il sistema produce comunque: il costo di questo passo è nullo, perché non
    richiede esecuzioni aggiuntive.
    \item \textbf{Decidere se la leva esiste.} Una quota prossima al \SI{100}{\percent} la
    esclude, e la decisione è di non potare; una quota bassa lascia la leva aperta, e si prosegue.
    \item \textbf{Misurare la copertura su un sottoinsieme} di configurazioni, applicando la
    strategia e confrontando il numero di celle valorizzate con quello della configurazione di
    partenza. È il solo passo che richieda esecuzioni, ed è quindi il costo del criterio.
    \item \textbf{Adottare la potatura solo se la copertura non cala}; in caso contrario ricorrere
    alla taratura conservativa, che riduce il rischio rinunciando a parte della riduzione.
\end{enumerate}

\paragraph{Che cosa il criterio decide e che cosa no.} Il primo termine --- se la leva esista --- è
decidibile a costo nullo e prima di qualunque esecuzione. Il secondo --- se la copertura calerà
--- non è prevedibile senza eseguire il sistema, e il criterio non è quindi sufficiente a decidere
in autonomia: la sua funzione è di \emph{escludere} i casi in cui la potatura non può che
danneggiare, non di autorizzarla. La \autoref{tab:criterion} riassume le tre famiglie su cui il
criterio è stato derivato.

\begin{table}[htbp]
\centering
\caption{Applicazione del criterio alle tre famiglie su cui è stato derivato. Le grandezze a
cui il criterio si riferisce sono quelle della \autoref{tab:mechanism}, che non vengono qui
ripetute. La tabella documenta i casi da cui il criterio è stato ricavato e non costituisce una
sua validazione: la prova su una famiglia indipendente è riportata nella
\autoref{sec:physician}.}
\label{tab:criterion}
\footnotesize
\begin{tabular}{@{}lp{4.5cm}p{5.4cm}@{}}
\toprule
famiglia & decisione del criterio & esito osservato \\
\midrule
Glass & \emph{non potare}: la leva non esiste & coerente: la potatura toglie copertura senza poter migliorare la composizione \\
Bridges & \emph{potare se la copertura tiene} & coerente: la copertura resta invariata e il risultato migliora \\
Restaurant & \emph{misurare prima di adottare} & coerente: la copertura cala e il risultato peggiora \\
\bottomrule
\end{tabular}
\end{table}

\paragraph{Perché il criterio non è una soglia adattata ai dati.} La grandezza impiegata non è un
punteggio scelto per separare i casi favorevoli da quelli sfavorevoli, ma il termine che il
meccanismo indica come necessario: il guadagno è il prodotto della quota esatta e del numero di
imputazioni, e una quota esatta pari al \SI{100}{\percent} annulla il primo fattore per
definizione. La soglia al \SI{100}{\percent} non è quindi un parametro da tarare; la zona
intermedia --- quote basse ma non nulle --- resta però indeterminata, ed è il limite principale
del criterio.

% =======================================================================================
\section{La prova sulla quarta famiglia}\label{sec:physician}

Il criterio della sezione precedente è stato derivato su tre famiglie. Il test che servirebbe è su
una famiglia mai impiegata, e l'archivio ne contiene una: la famiglia Physician, il cui schema è
allineato ai propri file di regole. La prova è stata condotta su tutto il materiale eseguibile, e
il suo esito è \emph{parzialmente conclusivo}: conferma l'equivalenza della potatura garantita e
fornisce due casi coerenti con il criterio, ma su un numero di celle troppo piccolo per fondare
una conclusione statistica.

\paragraph{Perché le configurazioni piene non sono impiegabili.} I dataset non remodulati
contengono \num{103592} righe e fino a \num{121913} celle nulle. La ragione dell'impraticabilità
non è però il tempo di imputazione: è il \emph{preprocessing interno} del sistema, che all'avvio
verifica per ciascuna regola a soglia di destra nulla se il suo lato sinistro si comporti da
chiave, con un confronto a coppie di tuple quadratico nel numero di righe
(\autoref{sec:dual-role}). La sonda condotta sulla configurazione più piccola
(\texttt{Physician\_Compare\_National\_13467\_1}, \num{13467} celle) è stata interrotta dopo
\SI{1200}{\second} senza che il sistema avesse prodotto \emph{una sola} riga di risultato: il
tempo era ancora interamente assorbito dal preprocessing. La circostanza è coerente con il
materiale d'archivio, che per la medesima configurazione contiene un file di risultati parziale di
\num{1005} righe e un file di celle popolate vuoto.

\paragraph{Che cosa è stato possibile eseguire.} Delle quattro versioni remodulate, due --- 0.05 e
0.1 --- non sono eseguibili per il difetto D3 (\autoref{sec:d3}): il sistema si arresta
all'avvio su righe del dataset con un numero di campi inferiore allo schema. Le due versioni
eseguibili, 0.005 e 0.01, valutano \num{13} e \num{27} celle; escludendo dalle altre due le celle
non localizzabili, il materiale eseguibile salirebbe a \num{329} celle, ma la riparazione delle
righe malformate non è stata applicata perché il costo di preprocessing la rende non conveniente
rispetto al contributo ottenibile (\autoref{sec:d3}).

\paragraph{Verifica di equivalenza.} La potatura garantita è stata applicata a tutte le
configurazioni eseguibili, con la politica di tutela integrale dei veti. L'esito è netto: su
\num{2} configurazioni su \num{2} i file delle celle popolate e dei risultati sono
\emph{byte-identici} a quelli della configurazione di partenza, con riduzioni dell'insieme di
regole del \SI{6}{\percent} e del \SI{5}{\percent}. È la prima verifica della condizione \Vtwo{}
su una famiglia indipendente da quelle su cui è stata derivata, e si aggiunge alle \num{111}
configurazioni già verificate.

\paragraph{Il test del criterio.} La \autoref{tab:res-physician} riporta i due casi eseguibili.

\begin{table}[htbp]
\centering
\caption{Prova esplorativa sulla famiglia Physician. Confronto fra configurazione di partenza e
potatura aggressiva sulle due versioni eseguibili; ``esatte'' sono le imputazioni ottenute per via
esatta, ``corrette'' quelle conformi al ground truth.}
\label{tab:res-physician}
\footnotesize
\begin{tabular}{lrrrrrrr}
\toprule
versione & celle & imputate & esatte & quota esatta & corrette & copertura & riduzione \\
\midrule
0.005, partenza & 13 & 3 & 3 & \SI{100}{\percent} & 2 & \SI{23.1}{\percent} & --- \\
0.005, \str{A5b} & 13 & 3 & 3 & \SI{100}{\percent} & 2 & \SI{23.1}{\percent} & \SI{67}{\percent} \\
0.01, partenza  & 27 & 11 & 5 & \SI{45}{\percent} & 4 & \SI{40.7}{\percent} & --- \\
0.01, \str{A5b}  & 27 & 9 & 3 & \SI{33}{\percent} & 3 & \SI{33.3}{\percent} & \SI{74}{\percent} \\
\bottomrule
\end{tabular}
\end{table}

I due casi hanno esiti diversi e istruttivi. Sulla versione 0.005 la configurazione di partenza
ottiene \emph{tutte} le proprie imputazioni per via esatta: la leva della qualità non esiste, il
criterio prescrive di non potare, e la potatura aggressiva --- pur rimuovendo due terzi delle
regole --- non cambia nulla. Sulla versione 0.01 la quota esatta di partenza è del
\SI{45}{\percent}, la leva esiste, e la potatura \emph{peggiora} il risultato: le celle valorizzate
scendono da \num{11} a \num{9}, le imputazioni esatte da \num{5} a \num{3}, le corrette da \num{4}
a \num{3}. Il criterio prescrive in questo caso di misurare la copertura prima di adottare la
potatura, e la misura dice di non adottarla.

Entrambi gli esiti sono quindi coerenti con il criterio, uno per esclusione della leva e uno per
caduta della copertura. Il valore probatorio resta però limitato, per tre ragioni che vanno
dichiarate: le celle valutabili sono \num{40} in tutto; la soglia disponibile è una sola
(\num{3}, l'unica per cui l'archivio contenga i file di regole); e su questa famiglia il sistema
imputa una frazione minoritaria delle celle, cosicché le differenze osservate valgono poche unità.
La prova va letta come un \emph{controllo di coerenza} su dati nuovi, non come una validazione.

\section{Risultati negativi}\label{sec:negative-results}

Una parte rilevante delle strategie sperimentate non ha prodotto il guadagno atteso. I risultati
sono riportati perché delimitano lo spazio delle soluzioni.

\paragraph{Il segnale di precisione non è la causa del guadagno del riordino.}
Il riordino dei cluster per precisione empirica produce un guadagno misurabile
($+96$ imputazioni corrette su \num{30} configurazioni, $p = 2.5\cdot 10^{-6}$). L'ipotesi che tale
guadagno fosse attribuibile al segnale di precisione è stata sottoposta a un esperimento di
controllo: a parità di insieme di regole, si sono confrontati l'ordinamento per precisione,
l'ordinamento inverso e tre permutazioni casuali indipendenti.

\begin{table}[htbp]
\centering
\caption{Esperimento di controllo sull'ordinamento dei cluster. Quindici configurazioni, cinque
bracci, insieme di regole invariato.}
\label{tab:res-order}
\footnotesize
\begin{tabular}{lrrrr}
\toprule
braccio & corrette & $\Delta$ vs orig.\ & migl.\ / pegg.\ & Wilcoxon \\
\midrule
originale                      & 333 & ---    & --- & --- \\
riordino per precisione        & 386 & $+53$  & 13 / 0 & $p = \num{0.0016}$ \\
permutazione casuale (media di 3) & --- & $+52.7$ & 13 / 0 & $p = \num{0.0016}$ \\
riordino inverso               & 313 & $-20$  & 6 / 1 & $p = \num{0.2702}$ \\
\bottomrule
\end{tabular}
\end{table}

Il confronto decisivo è fra il riordino per precisione e la permutazione casuale: la differenza è
$+0.33$ imputazioni corrette, con $p = \num{0.8878}$, e una delle permutazioni casuali supera
l'ordinamento per precisione. Il guadagno \emph{non è attribuibile al segnale impiegato}. Il
riordino inverso, d'altra parte, peggiora il risultato rispetto all'originale: non è dunque vero
che qualunque permutazione funzioni, ma che l'ordine di partenza è prossimo al caso peggiore e
che qualunque ordine che non sia l'inverso lo migliora.

\paragraph{La regolarizzazione della stima non ha effetto.} La sostituzione della stima di
precisione con una stima contratta verso un valore a priori è stata sperimentata su quattro
valori del parametro di contrazione. I risultati sono \emph{identici} in tutte le
configurazioni per tutti i valori: il parametro non è una variabile da tarare.

\paragraph{Il riordino entro il cluster è inerte per costruzione.} La procedura che genera i
candidati percorre tutte le regole di un cluster e ne conserva il minimo, senza uscite
anticipate. L'ordine interno al cluster non può quindi influenzare il valore imputato. La
previsione è stata confermata empiricamente su \num{15} configurazioni su \num{15}, con zero
differenze non nulle al test.

\paragraph{L'ordinamento delle rimozioni sotto vincolo di budget è un vicolo cieco.} La politica
\str{budget:F} ordina le regole da rimuovere secondo un indice di rischio. Tale indice è risultato
\emph{anti-correlato} con la rilevanza effettiva delle regole: su \num{8961} regole, la
correlazione di rango con il numero di celle che ciascuna regola può servire è $\rho = -0.237$.
La causa è stata individuata: l'indice somma i valori mancanti sulle colonne del lato
\emph{sinistro}, mentre una regola può servire soltanto celle della propria colonna di
\emph{destra}. La correzione dell'indice produce un miglioramento strutturale netto --- a parità
di budget, rimuove $2.5$ volte più regole che non servono alcuna cella e distrugge $2.6$ volte
meno celle servibili --- ma il vantaggio \emph{non si traduce in un miglioramento dell'esito}:
su \num{13} configurazioni distribuite su tre famiglie, il saldo delle imputazioni corrette è
identico. La conclusione è che, sotto vincolo di budget, conta la \emph{dimensione} del budget e
non l'\emph{ordine} con cui si sceglie.

\paragraph{Una potatura esatta della generazione può rallentare il sistema.} La strategia
\str{score-argmin}, che preserva la generazione calcolando staticamente quali regole non sono mai
l'argomento del minimo, raggiunge riduzioni fino al \SI{98}{\percent} e preserva la copertura.
Rende però l'esecuzione fino a \num{251} volte più lenta. Il risultato mostra che la relazione fra
numero di regole e tempo di esecuzione non è monotona, e che la riduzione non è di per sé un
indicatore di efficienza.

% =======================================================================================
\section{Il costo: perché ridurre non basta}\label{sec:res-cost}\label{sec:score-argmin}

\begin{figure}[htbp]
\centering
\begin{tikzpicture}
\begin{axis}[
    width=0.80\textwidth,
    height=0.44\textwidth,
    xlabel={riduzione delle regole (\si{\percent})},
    ylabel={fattore di tempo (scala logaritmica)},
    xmin=0, xmax=105,
    ymode=log,
    ymin=0.004, ymax=900,
    ytick={0.01,0.1,1,10,100},
    yticklabels={$0.01$, $0.1$, $1$, $10$, $100$},
    grid=major,
    grid style={black!12},
    axis lines=left,
]
\addplot[only marks, mark=*, mark size=3.4pt, color=black]
    coordinates {(98.5,0.0072) (84.2,0.091) (17.1,0.855) (45.7,251)};
\draw[dashed, black!45] (axis cs:0,1) -- (axis cs:105,1);
\node[font=\scriptsize, anchor=west] at (axis cs:60,0.0072) {\str{score-argmin:allow}, Bridges @15};
\node[font=\scriptsize, anchor=west] at (axis cs:52,0.11) {\str{A5b}, Bridges};
\node[font=\scriptsize, anchor=west] at (axis cs:19,0.62) {\str{veto-cover:protect}, Bridges @9};
\node[font=\scriptsize, anchor=east] at (axis cs:44,330) {\str{score-argmin:budget:0.5}, Glass @9};
\end{axis}
\end{tikzpicture}
\caption{Riduzione ottenuta e fattore di tempo per quattro casi misurati. Il fattore è il rapporto
fra il tempo di esecuzione della configurazione potata e quello della configurazione di partenza:
valori inferiori a uno indicano un'accelerazione. La relazione fra le due grandezze non è
monotona, e il caso con la riduzione maggiore è anche il più veloce mentre un caso con riduzione
inferiore è il più lento.}
\label{fig:time}
\end{figure}

L'ipotesi H16 prevedeva che il tempo di esecuzione seguisse il numero di regole rimosse. I dati la
falsificano, e il modo in cui la falsificano è più interessante della previsione.

Il costo di \emph{preprocessing} della potatura non è invece in discussione: la riscrittura del
file delle regole è trascurabile rispetto all'esecuzione del sistema e cresce con il numero di
regole, attestandosi nell'ordine dei secondi sui file maggiori (\num{27445} regole). Il confronto
fra i due costi è quindi interamente a favore della potatura, e il problema non è il costo
dell'analisi ma l'effetto sull'esecuzione.

\str{score-argmin} rimuove fino al \SI{98}{\percent} delle regole e rende il sistema
\emph{drammaticamente più lento}: il rallentamento mediano è di \num{5.8} volte con la variante
\str{allow} e di \num{20.1} volte con \str{budget:0.5}, con un massimo di \num{251} volte sulla
configurazione \texttt{glass\_43\_2} a soglia 9, dove il tempo di esecuzione passa da
\SI{3.5}{\second} a \SI{873}{\second} mentre le regole scendono da \num{11648} a \num{6323}. Su
quattordici configurazioni, la variante \str{allow} è più veloce della configurazione di partenza
in \emph{una sola}.

\begin{table}[htbp]
\centering
\caption{Rallentamento del tempo di esecuzione prodotto da \str{score-argmin} su Glass, rispetto
alla configurazione di partenza. Il rallentamento è il rapporto fra il tempo della variante potata
e quello della variante non potata, sulla stessa configurazione.}
\label{tab:res-slowdown}
\footnotesize
\begin{tabular}{lrr}
\toprule
variante & rallentamento mediano & rallentamento massimo \\
\midrule
\str{score-argmin:allow}       & \num{5.8}$\times$  & \num{28}$\times$ \\
\str{score-argmin:budget:0.5}  & \num{20.1}$\times$ & \num{251}$\times$ \\
\bottomrule
\end{tabular}
\end{table}

La circostanza che esclude la spiegazione ovvia è un confronto interno alla tabella: la variante
\str{budget:0.5} \emph{conserva più regole} di \str{allow} --- \num{6323} contro \num{1098} sulla
configurazione citata --- ed è circa nove volte più lenta. Non è dunque la \emph{quantità} delle
regole a determinare il tempo, ma la loro \emph{composizione}: il tempo dipende da quali regole
sopravvivono, non da quante.

Una spiegazione candidata esiste ed è coerente con l'analisi del sistema, ma \emph{non è
verificata} e va presentata come tale. Il sistema dispone di una scorciatoia esatta che chiude una
cella in tempo costante quando la tupla presenta una sola cella nulla e la regola impiegata ha
soglia di destra nulla (\autoref{sec:two-branches}). Se l'insieme potato non contiene più le regole
che attivano quella scorciatoia per la maggior parte delle celle, ogni cella ricade nella ricerca
per similarità sull'intero dataset, il cui costo è lineare nel numero di righe anziché costante.
La potatura rimuoverebbe cioè, insieme alle regole dannose, la via rapida. La verifica richiede di
strumentare il percorso della scorciatoia, che non è stato possibile isolare con la
strumentazione disponibile, e la questione è registrata fra gli sviluppi
(\autoref{sec:future-work}).

La conseguenza pratica è netta ed è il motivo per cui la sezione è necessaria: \str{score-argmin}
non è una strategia utilizzabile, malgrado il criterio esatto che la fonda e la conferma
sperimentale della sua proprietà principale. Un guadagno di copertura pagato con un tempo di
esecuzione da sei a duecentocinquanta volte maggiore non è un compromesso proponibile. Il caso
mostra inoltre che la domanda del lavoro --- se la potatura migliori \emph{e} acceleri --- non
ammette una risposta unica: le due grandezze possono muoversi in direzioni opposte, e la
valutazione di una strategia deve riportarle entrambe.

% =======================================================================================


    \chapter{Discussione}\label{chap:discussion}


\noindent Questo capitolo raccoglie la lettura dei risultati e non introduce grandezze nuove: il
suo contenuto è di interpretazione, e le misure a cui si riferisce sono quelle presentate nel
capitolo precedente. Vi sono ricostruiti gli esiti delle ipotesi e il quadro complessivo che ne
emerge, dichiarati i limiti dell'evidenza raccolta, e collocato il lavoro rispetto alla letteratura
--- la riduzione classica di un insieme di regole e la selezione guidata dai metadati ---
indicando quali verifiche estenderebbero i risultati e quali il materiale disponibile non ha
consentito.

\section{Esiti delle ipotesi}\label{sec:hypotheses-outcomes}

La \autoref{tab:hypotheses-outcomes} raccoglie gli esiti. Le ipotesi sono quelle enunciate nella
\autoref{tab:hypotheses} e nella \autoref{tab:exploratory}; nessuna è stata modificata dopo la
misura.

\begin{table}[htbp]
\centering
\caption{Esiti delle ipotesi. ``Confermata'' indica che la previsione è stata confermata dai dati;
``falsificata'' che i dati la contraddicono. Per H1 il valore riportato è quello ottenuto dopo la
correzione del difetto D5 (\autoref{sec:d5}): con il codice precedente la riduzione misurata era
dello \SI{0.1}{\percent}.}
\label{tab:hypotheses-outcomes}
\footnotesize
\begin{tabular}{@{}cp{3.6cm}p{4.0cm}p{3.4cm}@{}}
\toprule
\# & ipotesi & previsione & esito \\
\midrule
H1 & \str{score-argmin} con veti protetti & riduzione $\ge \SI{30}{\percent}$ & \textbf{falsificata} (\SI{0.3}{\percent}) \\
H6 & \str{veto-cover} su tutte le famiglie & output identico & \textbf{confermata} \num{113}/\num{113} \\
H7 & verifica empirica di \Vtwo{} & zero discrepanze & \textbf{confermata} \\
H8 & \str{precision-priority} da solo & guadagno di qualità & \textbf{confermata} ($+96$ corrette su \num{30} configurazioni) \\
H11 & \str{precision-priority} altrove & guadagno su Glass e Restaurant & \textbf{falsificata} \\
H14 & guadagno attribuibile alla precisione & il casuale non produce lo stesso guadagno & \textbf{falsificata} \\
H16 & riduzione e tempo proporzionali & tempi minori & \textbf{falsificata} (fino a \num{251}$\times$) \\
\midrule
H2  & \str{score-argmin} altrove & copertura non cala & confermata \\
H3  & copertura esatta e gratuita & nessuna perdita & falsificata \\
H4  & conversione \texttt{B}$\to$\texttt{D} & guadagno gratuito & falsificata \\
H5  & riordino delle righe del dataset & guadagno & falsificata \\
H9  & riordino $+$ potatura & guadagno maggiore & sub-additiva \\
H10 & regolarizzazione $\alpha = 3$ & guadagno $\ge \SI{20}{\percent}$ & falsificata \\
H12 & la regolarizzata generalizza meglio & guadagno maggiore & falsificata \\
H13 & riordino entro il cluster & leva di qualità & falsificata \\
H15 & ordinamento del budget migliorabile & guadagno & falsificata \\
\bottomrule
\end{tabular}
\end{table}

\paragraph{Le ipotesi confermate riguardano la teoria.} H6, H7 e H8 sono state confermate senza
riserve, e riguardano la garanzia e la leva di qualità presa da sola: la condizione \Vtwo{}
sopravvive alla verifica indipendente, la sua applicazione non altera il comportamento del sistema
su alcuna configurazione, e il riordino per priorità di precisione produce un guadagno quando è
impiegato da solo. È il nucleo dimostrativo del lavoro.

\paragraph{Le ipotesi falsificate riguardano le euristiche e le attribuzioni.} H1, H11, H14 e H16
sono state falsificate, e la loro falsificazione ha un valore che va oltre il caso singolo. H1
mostra che la strategia esatta lato generazione non è utilizzabile con la politica di tutela,
perché la classe garantita è troppo piccola per produrre la riduzione attesa. H11 mostra che il
guadagno non generalizza. H14 mostra che il segnale che si riteneva responsabile del guadagno non
lo spiega: il riordino casuale produce lo stesso effetto, e l'interpretazione che era stata
formulata in una fase precedente del lavoro è stata corretta. H16 mostra che la riduzione non
implica il risparmio di tempo.

\paragraph{La lettura d'insieme.} Il quadro che emerge è coerente e si può riassumere in quattro
punti. \emph{Primo}: esiste una potatura esatta, ed è piccola --- la condizione \Vtwo{} è corretta
e verificata, ma la riduzione che consente è limitata, e il limite è intrinseco al criterio
(\autoref{sec:maxextra}). \emph{Secondo}: esiste una potatura che migliora il risultato, ed è
condizionata --- \str{A5b} riduce dell'\SI{84.2}{\percent} e migliora il risultato su Bridges,
mentre sugli altri dataset danneggia, e la condizione di applicabilità è ora nota e misurabile.
\emph{Terzo}: il guadagno non proviene dal segnale che si riteneva responsabile, e la spiegazione
corretta è il prodotto di due fattori, la quota di imputazioni ottenute per via esatta e il numero
di imputazioni effettuate. \emph{Quarto}: il tempo non segue la riduzione, e una strategia può
rimuovere quasi tutte le regole peggiorando insieme il risultato e la durata.

Va infine osservato che i tre interventi studiati --- la potatura esatta, la potatura aggressiva e
la potatura conservativa --- non sono ordinabili. La prima non altera mai il risultato ma riduce
poco; la seconda produce il guadagno maggiore ma soltanto su una famiglia; la terza non peggiora
quasi mai ma rinuncia, nella maggior parte delle configurazioni, a qualunque riduzione. La scelta
fra le tre dipende dall'obiettivo: garanzia, guadagno, o assenza di rischio.

% =======================================================================================
\section{Limiti dell'evidenza}\label{sec:res-limits}

I limiti che seguono sono dichiarati perché delimitano la portata delle conclusioni, e non perché
ridimensionino i risultati: ciascuno indica anche la verifica che servirebbe a superarlo.

\paragraph{Il criterio diagnostico è post-hoc.} Il meccanismo a due fattori e il criterio che ne
discende sono stati ricavati osservando i risultati su tre famiglie, non previsti prima della
misura. La prova su una quarta famiglia è stata condotta (\autoref{sec:physician}) e ha prodotto
due esiti coerenti con il criterio, ma su \num{40} celle con copertura non nulla, a una sola
soglia, e senza che gli esiti siano distinguibili dal caso: due delle quattro versioni non sono
eseguibili per righe malformate nel dataset (\autoref{sec:d3}). Il criterio va quindi inteso come
una spiegazione con evidenza preliminare, non come una legge verificata su dati indipendenti.

\paragraph{Una sola famiglia con guadagno pieno.} Il guadagno di qualità è stato osservato su
Bridges; su Glass e Restaurant la stessa strategia danneggia, e la taratura conservativa riduce il
danno senza produrre guadagno. La generalità del meccanismo è sostenuta dalla coerenza dei tre
casi con la stessa legge, non da una molteplicità di famiglie favorevoli.

\paragraph{Le misure di tempo provengono da una sola macchina.} Le campagne sono state condotte su
una macchina a otto core, con esecuzioni strettamente sequenziali, e i valori assoluti dipendono
dall'ambiente. I confronti sono appaiati --- la stessa configurazione con e senza potatura --- e
sono quindi robusti rispetto all'ambiente, mentre i fattori di accelerazione e di rallentamento
riportati nel capitolo vanno letti come ordini di grandezza.

\paragraph{La numerosità è quella delle configurazioni.} L'unità di analisi è la configurazione,
non la cella nulla, e i test impiegano quindi alcune decine di osservazioni
(\autoref{sec:verification}). La potenza statistica è conseguentemente limitata: le conclusioni
sulle strategie il cui effetto è piccolo rispetto alla variabilità fra configurazioni vanno
considerate indicative. I test sono inoltre numerosi rispetto alle dimensioni del campione, e il
capitolo riporta per questo i valori corretti per confronti multipli accanto a quelli grezzi
(\autoref{sec:statistics}).

\paragraph{Il perimetro non comprende il confronto con altri metodi di imputazione.} Il lavoro
studia l'effetto della selezione delle regole \emph{a parità di sistema}: il termine di paragone è
la stessa esecuzione senza potatura, non un algoritmo di imputazione alternativo. La scelta
discende dal vincolo dichiarato --- la potatura opera sul file delle regole di un sistema
specifico --- e va tenuta presente nel leggere i risultati: essi non dicono se il sistema studiato
imputi meglio o peggio di altri, ma che cosa si guadagni o si perda potando le sue regole.

\paragraph{Due famiglie dell'archivio non sono utilizzabili.} La famiglia Nursery non contiene
dataset e la famiglia Physician, nelle versioni piene, ha dimensioni incompatibili con una
campagna (\num{103592} righe e fino a \num{121913} celle nulle). La griglia effettiva è quindi
composta da tre famiglie, ed è la ragione per cui la prova su una famiglia indipendente non ha
potuto che essere esplorativa.

% =======================================================================================
\section{Collocazione e verifiche ulteriori}\label{sec:positioning}

\paragraph{Rispetto alla riduzione classica di un insieme di regole.} La teoria classica delle
dipendenze funzionali affronta la ridondanza in forma dichiarativa: una dipendenza è superflua se
la sua rimozione non cambia l'insieme dei vincoli soddisfatti, e il calcolo di una copertura
minimale è un'operazione standard (\autoref{sec:rule-reduction}). Il lavoro mostra che quella
nozione non trasferisce al caso operativo, e la ragione è precisa: nel sistema studiato il ruolo di
una regola è definito da come il programma la consuma, non dalla sua semantica dichiarativa, e due
regole logicamente equivalenti possono avere effetti diversi perché differiscono per la direzione
delle soglie. La condizione \Vtwo{} è l'analogo operativo della dominanza classica, ma richiede
che la regola coprente sia \emph{più permissiva} a sinistra e \emph{più severa} a destra: la
direzione opposta a quella che si otterrebbe ragionando sui soli vincoli. Il teorema negativo
stabilisce che con l'informazione disponibile a una strategia esterna non si può andare oltre
questa classe senza rinunciare alla garanzia.

La distinzione non è soltanto argomentata: è misurata. Applicando il criterio classico come
strategia di potatura (\str{classical-cover}) sulle configurazioni di Bridges a soglia 9, la
riduzione raggiunge il \SI{97.4}{\percent} --- molto oltre il tetto della classe garantita --- ma
l'output cambia in tutte le configurazioni, l'accuratezza mediana perde \num{14.3} punti e le
imputazioni corrette scendono da \num{114} a \num{80} (\autoref{sec:res-classical}). Il confronto
mostra insieme due cose: che la ridondanza dichiarativa non implica la sicurezza operativa, e che
il numero di regole rimosse, da solo, non è un indicatore di qualità di una strategia.

\paragraph{Rispetto alla selezione guidata dai metadati.} Una linea di ricerca recente impiega
modelli linguistici per inferire relazioni fra colonne dal loro nome o per guidare
l'esplorazione di fonti eterogenee (\autoref{sec:semantic-pruning}). L'ambizione è la stessa ---
impiegare il significato invece della forma --- ma l'oggetto è diverso: quelle tecniche
\emph{interpretano} i metadati per produrne di nuovi, mentre il lavoro qui presentato
\emph{seleziona} metadati esistenti rispetto al consumo che ne fa un algoritmo noto. La differenza
ha una conseguenza pratica: la garanzia di equivalenza che il \autoref{chap:theory} dimostra non
dipende dalla qualità di una stima, ed è quindi verificabile per confronto di output; una
selezione guidata da un modello linguistico richiederebbe invece una validazione statistica del
modello stesso.

\uninatodo{Manca la misura dell'esercizio effettivo dei veti.}

\uninatodo{Manca il valore marginale del pruning esterno rispetto alla riduzione interna.}

\paragraph{Che cosa renderebbe più forti le conclusioni.} Tre verifiche sono state identificate
come decisive e non sono state condotte, ciascuna per una ragione dichiarata. La prima è la
\emph{misura dell'esercizio effettivo dei veti}: il meccanismo a due fattori è oggi inferito dalla
composizione degli output, e la strumentazione già disponibile consentirebbe di contare le
verifiche effettivamente eseguite prima e dopo la potatura, sostituendo l'inferenza con una misura
(\autoref{sec:future-work}). La seconda è il \emph{valore marginale della potatura esterna}
rispetto alla riduzione che il sistema compie già da sé sulle regole a soglia nulla con lato
sinistro chiave (\autoref{sec:dual-role}): quantificarlo richiede di calcolare la sovrapposizione
fra le due classi, e direbbe quanto del lavoro è già svolto dall'eseguibile. La terza è la
\emph{validazione prospettica del criterio} su una famiglia con quota esatta intermedia e
copertura stabile: la prova sulla quarta famiglia ha fornito due casi coerenti ma su \num{40}
celle, e non è sufficiente (\autoref{sec:physician}).

\paragraph{Che cosa non è stato possibile verificare.} Le configurazioni piene della famiglia
Physician restano fuori portata perché il preprocessing interno del sistema è quadratico nel numero
di righe, e le due versioni remodulate riparabili non sono state ripristinate perché il costo di
preprocessing non è giustificato dal numero di celle che se ne ricaverebbero
(\autoref{sec:d3}). Il confronto con metodi di imputazione alternativi è fuori dal perimetro
dichiarato del lavoro, che misura l'effetto della selezione delle regole a parità di sistema
(\autoref{sec:res-limits}). Entrambe le circostanze sono limiti del materiale o del perimetro, non
dell'impostazione, e sono state dichiarate prima che un lettore le incontrasse.


    \chapter{Difetti e criticità del sistema}\label{chap:defects}


\noindent Questo capitolo raccoglie i difetti individuati nel corso dell'analisi. Essi non sono
un residuo accessorio del lavoro: la loro individuazione è stata resa necessaria dalla necessità
di stabilire se i risultati negativi fossero da attribuire alle strategie sviluppate o al
materiale di partenza, e la distinzione si è rivelata decisiva in più di un caso. Il capitolo
separa i difetti del sistema studiato da quelli del codice sviluppato in questo lavoro --- i primi
limiti del materiale, i secondi questione di attendibilità delle misure --- ne quantifica
l'impatto, ne indica il rimedio e ne trae le conseguenze generali per chi impiegasse il sistema su
altri dati.

% =======================================================================================
\section{Classificazione dei difetti}\label{sec:defects-taxonomy}

I difetti individuati appartengono a quattro categorie, che è opportuno tenere distinte perché
richiedono interventi di natura diversa. La distinzione che più conta, ai fini della lettura del
capitolo, è però un'altra: quella fra i difetti del \emph{sistema studiato} --- quattro, tutti nel
archivio --- e gli errori del \emph{codice sviluppato in questo lavoro} --- due,
entrambi individuati e corretti durante l'analisi. I primi limitano ciò che si può concludere sul
sistema; i secondi riguardano l'attendibilità delle misure e sono esposti perché il loro
trattamento fa parte del metodo.

\begin{table}[htbp]
\centering
\caption{Difetti individuati, per categoria e provenienza.}
\label{tab:defects}
\footnotesize
\begin{tabular}{@{}clll@{}}
\toprule
\# & difetto & categoria & provenienza \\
\midrule
D1 & ramo irraggiungibile nel calcolo della similarità & algoritmo & materiale d'archivio \\
D2 & disallineamento di schema nei dataset Glass & dati & materiale d'archivio \\
D3 & indici di ground truth non ricostruibili (Physician) & dati & materiale d'archivio \\
D4 & eseguibile di build contenente istruzioni di debug & artefatto & materiale d'archivio \\
D5 & computo errato del budget nella politica sui veti & codice sviluppato & \textbf{questo lavoro} \\
D6 & indice di rischio anti-correlato con la rilevanza & codice sviluppato & \textbf{questo lavoro} \\
\bottomrule
\end{tabular}
\end{table}

% =======================================================================================
\begin{table}[htbp]
\centering
\caption{Impatto dei difetti, come sono stati rilevati e stato attuale. Le classi sono quelle
della \autoref{tab:defects}.}
\label{tab:defects-impact}
\footnotesize
\begin{tabular}{@{}lp{4.1cm}p{3.1cm}p{2.1cm}@{}}
\toprule
difetto & impatto quantificato & come è stato rilevato & stato \\
\midrule
D1 & la conversione \texttt{B}$\to$\texttt{D} è peggiore in \num{25} configurazioni su \num{25},
con una perdita di \num{21.7} punti di accuratezza & confronto fra i due rami sul bytecode, poi
misura & non corretto: il vincolo è di non modificare l'eseguibile \\
D2 & la famiglia Glass non è utilizzabile come distribuita: lo schema dei dataset è posizionale,
quello dei file di regole usa i nomi & prima esecuzione fallita, poi ispezione delle intestazioni & corretto esternamente, con rinomina documentata \\
D3 & due configurazioni non eseguibili per righe malformate; \num{115} celle non localizzabili su
\num{404}; celle valutabili da \num{329} a \num{40} & esecuzione fallita per eccezione di
indice fuori intervallo; conteggio delle righe con campi mancanti & non riparato: costo di preprocessing non giustificato \\
D4 & una build alternativa con due istruzioni di diagnostica, con comportamento identico
sull'output & confronto fra i due eseguibili sulla stessa configurazione & documentato: esclude che l'archivio contenga una versione emendata di D1 \\
D5 & riduzione allo \SI{0.1}{\percent} invece del \SI{17.1}{\percent} su Bridges a soglia 9;
\SI{43.7}{\percent} delle rimozioni dichiarate esatte senza copritore sopravvissuto & confronto fra
riduzione dichiarata ed effettiva su una pipeline composita & corretto e verificato su \num{8} configurazioni su \num{8} \\
D6 & indice di rischio con $\rho = -0.237$ su \num{8961} regole; dopo la correzione, $2.5$ volte
più regole inutili rimosse e $2.6$ volte meno celle servibili distrutte & validazione dell'indice
contro il numero di celle che ciascuna regola può servire & corretto, senza guadagno sull'esito \\
\bottomrule
\end{tabular}
\end{table}

\section{Difetti del sistema studiato}\label{sec:defects-system}

I quattro difetti che seguono appartengono al materiale d'archivio del sistema studiato. Essi hanno
una caratteristica comune: incidono sulla \emph{possibilità} di interpretare i risultati, non
soltanto sulla loro qualità, e in due casi hanno reso inutilizzabile una parte della griglia
sperimentale. Ciascuno è presentato con il sintomo, l'evidenza che lo stabilisce, l'impatto sui
risultati e il rimedio adottato.

\subsection{D1 --- Ramo irraggiungibile nel calcolo della similarità}\label{sec:d1}

Il sistema calcola il contributo di similarità di un attributo con una procedura che distingue
fra attributi numerici e attributi categorici. Il disassemblaggio mostra che la diramazione verso
il caso categorico è \emph{irraggiungibile}: la condizione di selezione confronta riferimenti a
oggetti anziché il contenuto, e non può quindi risultare vera. In conseguenza di ciò, un attributo
categorico viene elaborato dalla procedura numerica, che tenta una conversione a numero, fallisce,
e ricade su una distanza di stringa applicata a etichette che possono essere lunghe e prive di
relazione.

\begin{observation}
Il difetto è confermato dall'analisi statica e non è stato corretto: il vincolo di non modificare
l'eseguibile è stato mantenuto per tutta la durata della valutazione esterna. La sua esistenza
non è un'ipotesi ma un fatto verificabile nel bytecode.
\end{observation}

\paragraph{Impatto sull'analisi.} Il difetto ha conseguenze sulla tipizzazione delle colonne
impiegata nelle campagne. Si è verificato sperimentalmente che la conversione di una colonna da
categoria a numerica, lungi dal costituire un miglioramento gratuito, peggiora il risultato:
su \num{25} configurazioni, la conversione produce una perdita media di \SI{21.7}{} punti di
accuratezza, con un peggioramento in \num{25} configurazioni su \num{25}. Il risultato ha
richiesto di rivedere una conclusione formulata in una fase precedente del lavoro, nella quale la
conversione era stata presentata come un guadagno.

% =======================================================================================
\subsection{D2 --- Disallineamento di schema nei dataset Glass}\label{sec:d2}

I dataset della famiglia Glass presentano nomi di colonna in forma posizionale
(\texttt{A}, \texttt{B}, \dots), mentre i corrispondenti file di regole impiegano nomi
descrittivi. Per le altre famiglie i due insiemi di nomi coincidono; per Glass no.

Il disallineamento non è innocuo. Il sistema risolve le colonne \emph{per posizione}, non per
nome: quando gli schemi non coincidono, la corrispondenza fra le colonne del dataset e quelle
delle regole è apparentemente corretta ma sostanzialmente errata, e il sistema scrive le
imputazioni nella colonna sbagliata. La condizione è di corruzione silenziosa: nessun errore
viene segnalato.

\paragraph{Intervento.} La correzione adottata è esterna e consiste nel rinominare
posizionalmente le colonne del dataset perché coincidano con quelle dei file di regole, e nel
riscrivere di conseguenza il materiale di ground truth. L'ordine delle colonne non è modificato,
sicché gli indici posizionali impiegati dal sistema restano validi.

\paragraph{Impatto sull'analisi.} Prima della correzione la famiglia Glass risultava
inutilizzabile, con copertura nulla. Dopo la correzione la copertura si attesta intorno al
\SI{30}{\percent}. Va osservato che il difetto ha un effetto anche sull'interpretazione delle
statistiche di training: i file di risultato prodotti prima della correzione registrano valori
non imputati per tutte le celle, e non sono quindi impiegabili come riferimento. La circostanza ha
richiesto di raccogliere statistiche ex novo su esecuzioni condotte con lo schema corretto.

% =======================================================================================
\subsection{D3 --- Righe malformate e indici non ricostruibili nei dataset Physician}
\label{sec:d3}

La famiglia Physician comprende quattro versioni ``remodulate'' di dimensioni contenute, pensate
per la sperimentazione. Il loro materiale presenta due difetti distinti, che una verifica
successiva ha permesso di separare.

\paragraph{Righe con un numero di campi inferiore allo schema.} I dataset delle versioni 0.05 e
0.1 contengono righe malformate: \num{7}, \num{11} o \num{1} campi invece di \num{13}. Il file
della versione 0.05 termina a metà di un record, e nella 0.1 una riga mostra due record
concatenati, con un valore testuale che risulta tagliato e ricomposto (``\texttt{PENNSYLVANIA
COLLEGSICAL THERAPY}''). È questa la causa diretta della non eseguibilità: il controllo di chiave
che il sistema esegue all'avvio accede ai campi di una tupla per indice, e su una riga corta esce
con \texttt{IndexOutOfBoundsException}. La diagnosi è confermata dal materiale d'archivio: per
entrambe le versioni il log, il file delle celle popolate e il file dei risultati spediti con
l'archivio sono \emph{vuoti}, mentre per le versioni 0.005 e 0.01 sono popolati.

\paragraph{Indici di ground truth eccedenti il dataset.} Il materiale di ground truth --- in
formato \texttt{riga;attributo;valore} --- contiene indici di riga che eccedono il numero di righe
del dataset corrispondente, ed è ben formato: il difetto riguarda il \emph{riferimento}, non la
sintassi.

\begin{table}[htbp]
\centering
\caption{Difetto D3. Righe malformate e voci di ground truth con indice fuori intervallo, per
versione remodulata. Le colonne ``valutabili'' indicano le celle che restano dopo l'esclusione
degli indici fuori intervallo.}
\label{tab:physician-gt}
\footnotesize
\begin{tabular}{lrrrrr}
\toprule
versione & righe & righe malformate & voci ground truth & fuori intervallo & valutabili \\
\midrule
0.005 & 103  & 0 & 13  & 0   & 13 \\
0.01  & 207  & 0 & 27  & 0   & 27 \\
0.05  & 914  & \textbf{1} & 135 & \textbf{14}  & \textbf{121} \\
0.1   & 1374 & \textbf{2} & 269 & \textbf{101} & \textbf{168} \\
\bottomrule
\end{tabular}
\end{table}

Gli indici fuori intervallo provengono dal dataset completo, prima della riduzione, e non sono
stati rimappati; le righe malformate indicano che la riduzione è avvenuta per \emph{troncamento}
del file, circostanza che spiega anche perché gli indici non corrispondano più. Escludendo le voci
non localizzabili, le celle valutabili passano da \num{40} a \num{329}: la correzione non richiede
di modificare il dataset, ma soltanto di dichiarare quali celle si escludono.

\paragraph{La ricostruzione della corrispondenza non è praticabile.} Si è verificato se gli indici
potessero essere ricostruiti confrontando le righe delle versioni ridotte con quelle del dataset
completo. La verifica ha dato esito negativo: su \num{200} righe esaminate, \num{29} non trovano
corrispondenza e \num{73} risultano ambigue. Le versioni remodulate non sono un sottocampionamento
del dataset completo ma una sua perturbazione, e l'informazione necessaria alla rimappatura non è
contenuta nei dati disponibili.

\paragraph{Riparazione possibile, non intrapresa.} Le righe malformate possono essere sostituite
da righe di soli valori mancanti, conservando la numerazione su cui il ground truth è indicizzato
e producendo righe che il sistema ignora in entrambi i ruoli; la procedura è implementata e
documentata (\texttt{pruning/repair\_physician\_datasets.py}). La riparazione non è stata
applicata alla campagna, perché il costo di preprocessing del sistema su queste configurazioni ---
il controllo di chiave è quadratico nel numero di righe, e vale decine di minuti per esecuzione
--- non è giustificato dal numero di celle che se ne ricaverebbero, e perché il contributo
scientifico delle due versioni sarebbe comunque marginale rispetto alle versioni già eseguibili.

\begin{observation}
La famiglia Physician è utilizzabile come quarta famiglia per la verifica di equivalenza, e lo è
stata (\autoref{sec:physician}); per la verifica di qualità fornisce \num{40} celle con copertura
non nulla, un campione insufficiente a trarre conclusioni. La limitazione è del materiale, non
dell'impostazione sperimentale, e va dichiarata come tale.
\end{observation}

% =======================================================================================
\subsection{D4 --- Eseguibile di build con istruzioni di debug}\label{sec:d4}

L'archivio contiene un secondo eseguibile, di dimensioni ridotte, la cui struttura interna è
danneggiata (la directory centrale dell'archivio è assente). L'estrazione delle voci è comunque
possibile, e il confronto con l'eseguibile principale mostra una sola differenza: la presenza di
due istruzioni di stampa a scopo di diagnostica. Il comportamento è per il resto identico, come
verificato eseguendo entrambi gli eseguibili sulla medesima configurazione e confrontando gli
output.

Non si tratta quindi di una versione con correzioni, ma di una build di sviluppo. La circostanza
è rilevante in senso negativo: esclude che l'archivio contenga una versione emendata del difetto
D1.

% =======================================================================================
\section{Errori rilevati e corretti nel codice sviluppato}\label{sec:defects-ours}

I due difetti che seguono non appartengono al sistema studiato ma agli strumenti sviluppati in
questo lavoro. La loro esposizione è deliberata e costituisce parte del contributo: entrambi sono
stati individuati da verifiche interne --- il confronto fra riduzione dichiarata e riduzione
effettiva, la validazione di un indice di rischio mai messo alla prova --- e in un caso la
correzione ha cambiato il valore di una misura riportata in una fase precedente dell'analisi. Un
lavoro che misuri l'effetto di una trasformazione su un sistema opaco deve poter esibire il
controllo di qualità dei propri strumenti, perché è da lì che provengono gli errori più difficili
da vedere.

\subsection{D5 --- Computo errato del budget nella politica sui veti}\label{sec:d5}

Questo difetto appartiene al codice sviluppato in questo lavoro, ed è esposto perché la sua
individuazione e correzione costituiscono parte del contributo.

\paragraph{Il difetto.} La politica \str{budget:F} limita la frazione di regole portatrici di veto
che possono essere rimosse. Il computo della frazione includeva però anche le rimozioni
\emph{provate prive di costo} dalla condizione \Vtwo{}: regole il cui potere di veto è implicato
da una regola sopravvissuta. Tali rimozioni non sacrificano alcun veto, e addebitarle costituiva
una doppia contabilizzazione.

\paragraph{Secondo difetto, nella stessa procedura.} La condizione \Vtwo{} veniva valutata
sull'insieme di regole \emph{in ingresso} anziché su quello \emph{sopravvissuto}. Se una strategia
successiva rimuove la regola coprente, la garanzia decade. Su una pipeline composita si è
misurato che il \SI{43.7}{\percent} delle rimozioni presentate come esatte non aveva più alcun
copritore sopravvissuto.

\paragraph{Un terzo effetto, nella politica di tutela.} La politica \str{protect}, che deve
ripristinare le regole portatrici di veto non altrimenti garantite, ripristinava anche le regole
la cui rimozione era dimostrabilmente priva di costo, annullando di fatto il risultato del
teorema. Sulla configurazione esaminata la riduzione scendeva allo \SI{0.1}{\percent} invece del
\SI{17.1}{\percent} che la condizione autorizza.

\paragraph{Correzione e verifica.} La condizione \Vtwo{} è stata riformulata in modo da essere
valutata sull'insieme corrente delle regole sopravvissute, e il computo del budget è stato
corretto in modo che le rimozioni garantite non lo consumino. Dopo la correzione,
\str{veto-cover} con tutela attiva raggiunge la medesima riduzione della variante senza tutela
(\SI{17.1}{\percent} su Bridges a soglia 9), con output identico alla configurazione di partenza
su \num{8} configurazioni su \num{8}.

\begin{remark}
Le campagne condotte prima della correzione non sono inficiate: la strategia \str{veto-cover} era
sempre stata impiegata \emph{da sola}, circostanza nella quale l'insieme corrente coincide con
quello in ingresso e i difetti non possono manifestarsi. Il difetto era latente, non realizzato.
\end{remark}

% =======================================================================================
\subsection{D6 --- Indice di rischio anti-correlato}\label{sec:d6}

La politica \str{budget:F} sceglie quali regole rimuovere ordinandole per un indice di rischio
definito come somma dei valori mancanti sulle colonne del lato sinistro. L'indice non era mai
stato validato. La validazione, condotta confrontando l'indice con il numero di celle che ciascuna
regola può effettivamente servire, ha dato una correlazione di rango $\rho = -0.237$ su
\num{8961} regole: l'indice non è semplicemente debole, è \emph{anti-correlato}.

La causa è stata individuata con precisione: l'indice somma i valori mancanti sulle colonne del
lato \emph{sinistro}, mentre una regola può servire soltanto celle della propria colonna di
\emph{destra}. Ne segue che le regole la cui colonna di destra non presenta mai valori mancanti
--- e che quindi non possono servire alcuna cella --- ricevono un indice basso e vengono rimosse
per prime. L'ordinamento promuoveva a bersaglio esattamente le regole innocue.

La correzione adotta un indice calcolato sulle celle nulle reali. A parità di budget, il nuovo
ordinamento rimuove $2.5$ volte più regole che non servono alcuna cella e distrugge $2.6$ volte
meno celle servibili. Come riportato nella \autoref{sec:negative-results}, il miglioramento
strutturale \emph{non} si traduce in un miglioramento dell'esito, e la circostanza ha portato a
concludere che sotto vincolo di budget conti la dimensione del budget e non l'ordinamento.

% =======================================================================================
\section{Che cosa insegnano questi difetti}\label{sec:system-behaviour}

Alcune caratteristiche del sistema non costituiscono difetti ma sono rilevanti per
l'interpretazione dei risultati, e sono raccolte qui insieme alle conseguenze generali del
catalogo.

\paragraph{Codice non raggiungibile.} Il disassemblaggio individua procedure che il codice
applicativo non invoca mai: una procedura di ricerca locale affetta da un errore di segno, due
procedure di riordino delle tuple basate su euristiche, un comparatore di soglie affetto da
troncamento nella conversione a intero, e un meccanismo di gestione di regole privilegiate di
fatto inerte. La circostanza è stata verificata empiricamente oltre che staticamente.

\begin{observation}
Il comparatore di soglie merita una nota, perché il suo troncamento lo rende non transitivo e
avrebbe potuto spiegare l'irregolarità dell'ordine di partenza osservata nella
\autoref{sec:negative-results}. La verifica sui sorgenti decompilati mostra però che nessuna
classe lo istanzia: è codice morto, e non può essere la causa di alcun comportamento osservato.
L'ipotesi è stata scartata su questa base.
\end{observation}

\paragraph{L'argomento di selezione dell'euristica è inerte.} Il sesto argomento della riga di
comando, che dovrebbe selezionare l'euristica di riempimento, non è letto. La circostanza è
stata verificata eseguendo il sistema con valori diversi dell'argomento e confrontando gli
output, che risultano identici. Tutte le osservazioni di questo lavoro valgono per la sola
procedura effettivamente eseguita.

\paragraph{Che cosa avrebbe potuto rilevare prima questi difetti.} Nessuno dei sei difetti è stato
individuato da un test automatico del materiale d'archivio, e per ciascuno esiste una verifica
elementare che lo avrebbe fatto prima. D1 e D4 sono difetti di codice: un test di copertura dei
rami avrebbe mostrato che il ramo categorico non è mai eseguito, e un confronto fra gli eseguibili
distribuiti avrebbe mostrato le due istruzioni di diagnostica. D2 e D3 sono difetti di dati: un
controllo di schema --- il numero di campi per riga e la corrispondenza fra intestazione e
riferimenti --- li avrebbe segnalati al momento della distribuzione, senza bisogno di eseguire
nulla. D5 e D6 sono difetti propri del codice sviluppato in questo lavoro, e in entrambi i casi il
controllo che li ha rivelati è stato un confronto fra ciò che la procedura \emph{dichiara} di fare
e ciò che \emph{fa}: nel primo caso fra riduzione annunciata e riduzione effettiva, nel secondo fra
l'indice di rischio e la grandezza che dovrebbe approssimare. La raccomandazione che se ne trae è
di introdurre questi quattro controlli --- copertura dei rami, confronto fra artefatti, validazione
dello schema, coerenza fra dichiarazione ed effetto --- prima di misurare qualunque cosa su un
sistema opaco.

\paragraph{Che cosa il catalogo insegna sul progetto di algoritmi a RFD.} I sei difetti, letti
insieme, indicano tre proprietà ricorrenti. La prima è la distanza fra documentazione e
comportamento: un sistema può esporre un parametro che non legge, eseguire una sola delle
procedure che dichiara di offrire, e contenere rami che il codice non può raggiungere; per questo
l'analisi del \autoref{chap:renuver} è stata condotta sul bytecode e non sui documenti. La seconda
è la concentrazione del costo: la quasi totalità del tempo è in una procedura sola
(\autoref{sec:cost}), e ciò rende la riduzione dell'insieme di regole l'unica leva temporale di
rilievo, ma anche la più delicata, perché quella stessa procedura è quella che il pruning può
indebolire. La terza riguarda chi misura: due dei sei difetti appartengono agli strumenti
sviluppati in questo lavoro, e la loro individuazione ha richiesto di confrontare in modo
sistematico ciò che una strategia \emph{dichiara} di fare con ciò che \emph{fa}
(\autoref{sec:defects-ours}). Ne segue una raccomandazione operativa per chi impiegasse il sistema
su altri dati: verificare l'effettiva equivalenza degli output prima di attribuire un effetto a
una strategia di potatura, e considerare le configurazioni in cui il file prodotto coincide con
quello di partenza come prove nulle anziché come conferme.


    \newpage\pagestyle{conclusions}

\unnumberedchapter{Conclusioni e sviluppi futuri}
\label{chap:conclusions}

\unnumberedsection{Sintesi del lavoro}

Il lavoro è partito da una domanda operativa: se sia possibile selezionare le regole rilevanti in
un insieme di metadati impiegato per l'imputazione di valori mancanti, riducendo il costo
computazionale senza degradare la qualità delle imputazioni. Il sistema studiato, \renuver{},
riceve le regole come input esterno, e la selezione è stata quindi affrontata nella forma di
\emph{preprocessing}: riscrivere il file delle regole, senza toccare l'eseguibile.

La risposta a quella domanda si è articolata in tre passaggi. Il primo è stato la ricostruzione
del comportamento del sistema, per il quale non erano disponibili i sorgenti: l'analisi del
bytecode ha mostrato che una regola partecipa a due procedure con requisiti opposti --- la
generazione dei candidati e la verifica dei veti --- e che nulla, nel formato del file, distingue
i due ruoli. Il secondo passaggio è stato la derivazione delle condizioni esatte di potatura che
ne discendono, e del risultato negativo che ne consegue. Il terzo è stato la misurazione
sistematica dell'effetto delle strategie sviluppate, condotta su tre famiglie di dataset per un
totale di \num{1618} esecuzioni in configurazione potata.

Alla domanda iniziale il lavoro risponde su tre piani distinti, che conviene richiamare in chiusura
senza ripetere l'analisi: la \emph{riduzione garantita} non può superare il \SI{20.9}{\percent}
sui dataset esaminati, e il limite è del criterio e non dell'implementazione
(\autoref{chap:theory}); la \emph{qualità} può migliorare, ma soltanto dove la composizione delle
imputazioni lascia spazio al miglioramento e dove la copertura non cala, come mostra il confronto
fra Bridges e le altre famiglie (\autoref{chap:results}); il \emph{tempo} non segue la riduzione,
perché dipende da quali regole sopravvivono e non da quante (\autoref{sec:res-cost}). Le tre
risposte sono riassunte nella tabella della \autoref{sec:research-question} e discusse nel
\autoref{chap:discussion}, dove sono raccolti anche gli esiti delle ipotesi e i limiti
dell'evidenza.

\unnumberedsection{Possibili sviluppi}\label{sec:future-work}

\paragraph{Riscrivere le regole invece di rimuoverle.}
Il teorema negativo delimita le trasformazioni che \emph{eliminano} regole operando sul solo file.
Una trasformazione che riscriva le regole --- restringendo o allargando i lati sinistri, correggendo
le soglie, fondendo più regole in una --- non ricade in quel limite, e la condizione \Vtwo{} mostra
del resto che una regola può essere eliminata quando il suo contenuto informativo è conservato da
un'altra: la stessa conservazione potrebbe essere cercata per via sintattica anziché per
soppressione. È la direzione che più probabilmente consente di superare il tetto del
\SI{20.9}{\percent} senza rinunciare alla garanzia.

\paragraph{Completezza della condizione di copertura.}
La condizione \Vtwo{} è sufficiente ma non necessaria: stabilire se esistano altre classi di
regole il cui veto è implicato da una regola sopravvissuta --- e se siano decidibili con
l'informazione disponibile --- amplierebbe la portata quantitativa del risultato senza cambiarne la
natura. Il problema è aperto e riguarda la teoria delle dipendenze rilassate, non
l'implementazione.

\paragraph{Verificare il meccanismo del costo.}
L'ipotesi che il rallentamento osservato dipenda dalla perdita della scorciatoia esatta non è
verificata. Isolare il percorso della scorciatoia con la strumentazione e contare quante celle lo
impiegano, prima e dopo la potatura, è un esperimento circoscritto che trasformerebbe una
spiegazione plausibile in un dato, e che consentirebbe di prevedere il tempo di una strategia
oltre che la sua qualità.

\paragraph{Ottimizzare il percorso dei veti.}
La caratterizzazione del costo indica una direzione precisa e diversa da quella seguita finora.
La verifica dei veti esegue \num{354809} scansioni del dataset su \num{71} tuple distinte: la
circostanza suggerisce che il verdetto per coppia (regola, tupla) sia memoizzabile. Un
\emph{indice} sulle colonne del dataset, o un ordinamento delle regole che incontri presto quelle
che violano, sono direzioni alternative. Si tratta di interventi che richiedono la modifica del
sistema, ora disponibile come possibilità ma non ancora sperimentata su questo fronte.

\paragraph{Un teorema lato generazione.}
La condizione \Vtwo{} copre il ruolo di veto. Per il ruolo generativo non si dispone di una
condizione sufficiente analoga: la strategia \str{score-argmin} è esatta rispetto alla generazione
ma calcola l'argomento del minimo sull'insieme \emph{originale}, e la sua esattezza decade quando
altre regole sono state rimosse. Una condizione che garantisca l'invarianza dell'argomento del
minimo \emph{dopo} altre rimozioni colmerebbe la lacuna e renderebbe componibili i due criteri.

\paragraph{Misurare l'esercizio effettivo dei veti.}
La politica di tutela dei veti impiega un indice di rischio che, come si è mostrato, è
anti-correlato con la rilevanza. La sua sostituzione con una misura del numero di veti
effettivamente esercitati a runtime --- ottenibile dalla strumentazione già disponibile ---
consentirebbe di sostituire la stima con un dato.

\paragraph{Validare il criterio diagnostico su una famiglia indipendente.}
Il criterio della \autoref{sec:diagnostic} è stato derivato su tre famiglie e non è stato
sottoposto a validazione prospettica. Una famiglia con quota esatta intermedia e poche regole
costituirebbe il test decisivo.

Il risultato di questo lavoro non è una strategia di potatura che funzioni ovunque: è la
determinazione di \emph{che cosa} si può garantire, di \emph{quanto} si può guadagnare
abbandonando la garanzia, e del \emph{meccanismo} che decide quale dei due regimi si applichi a un
dataset dato. In un ambito in cui la letteratura presenta prevalentemente risultati positivi su
configurazioni favorevoli, la determinazione di un limite dimostrabile e la sua quantificazione
hanno un valore che non dipende dal segno del risultato.


    % NOTA: \pagestyle e' locale al gruppo corrente. Posto qui varrebbe anche per l'ultima
    % pagina del capitolo 9, che viene stampata prima di \begin{appendices}; posto dentro
    % l'ambiente non varrebbe per l'ultima pagina dell'appendice. Il cambio di stile e' quindi
    % fatto in modo GLOBALE all'inizio della prima appendice, dopo la pagina di apertura.
    \begin{appendices}
        % =======================================================================================
% Appendice: riproducibilità
% =======================================================================================
\chapter{Riproducibilità dei risultati}\label{app:reproducibility}
\begingroup\globaldefs=1 \pagestyle{appendices}\endgroup

\noindent Questa appendice documenta come i risultati presentati nel
\autoref{chap:results} possano essere riprodotti e verificati.

% =======================================================================================
\section{Struttura del materiale}\label{sec:appendix-layout}

Il materiale è organizzato nelle seguenti cartelle.

\begin{description}[leftmargin=0em]
    \item[\texttt{RENUVER/}] l'archivio di partenza: eseguibile, dataset, file di regole,
    materiale di ground truth e file di risultato delle esecuzioni originali. L'eseguibile non è
    mai stato modificato; la verifica di integrità è effettuata per checksum. L'archivio include
    un \texttt{README} che documenta la struttura delle cartelle e i parametri della riga di
    comando: le discrepanze fra quel documento e il comportamento osservato del codice sono
    discusse nella \autoref{sec:architecture} e nel \autoref{chap:defects}.
    \item[\texttt{pruning/}] gli strumenti sviluppati: il pruner, il sistema di valutazione, gli
    script di verifica e quelli di orchestrazione delle campagne.
    \item[\texttt{results/}] i file di risultato di tutte le campagne, in formato JSON, con una
    voce per ogni esecuzione e i relativi indicatori; in \texttt{results/analisi/} i dati delle
    prove che non appartengono a una campagna.
    \item[\texttt{analisi/}] l'analisi del sistema: disassemblaggio commentato, modello dei dati,
    inventario del codice non raggiungibile.
    \item[\texttt{report.md}] la documentazione estesa del lavoro, con la derivazione completa
    delle condizioni teoriche e il dettaglio di tutte le campagne.
\end{description}

% =======================================================================================
\section{Verifica automatica dei risultati}\label{sec:appendix-audit}

I numeri asseriti nella documentazione sono ricalcolati dai file di risultato da una procedura
automatica di verifica, che li confronta con i valori asseriti e segnala ogni divergenza. La
procedura copre \num{67} asserzioni, fra cui le condizioni teoriche derivate dal bytecode, le
metriche aggregate di ciascuna campagna e i risultati dei test statistici.

La procedura è stata impiegata durante la stesura: ha individuato e consentito di correggere tre
valori non aggiornati, fra cui un risultato il cui segno si invertiva estendendo il campione da
\num{14} a \num{15} configurazioni, e che era stato inizialmente riportato con il segno errato.
La circostanza motiva l'adozione della procedura: un risultato che cambia segno aggiungendo una
configurazione su quindici non è robusto, e non deve essere presentato come se lo fosse.

Due procedure ulteriori completano la verifica dei numeri citati nella tesi, ed entrambe operano
sui soli file di risultato, senza rieseguire il sistema.

\begin{sloppypar}
\begin{itemize}[leftmargin=1.4em]
    \item la prima, \texttt{stats\_effects.py}, ricalcola i test della famiglia primaria con la
    medesima ricetta di appaiamento impiegata dalla campagna --- chiave (dataset, soglia,
    ripetizione), baseline e variante entrambe concluse con esito regolare --- e produce, oltre ai
    valori di $p$ già pubblicati, la dimensione dell'effetto, l'intervallo di confidenza bootstrap
    della differenza mediana e il valore di $p$ corretto con la procedura di Holm
    (\autoref{tab:statistics});
    \item la seconda, \texttt{make\_appendix\_tables.py}, genera le tabelle complete
    dell'\autoref{app:results} dai file di risultato; i totali che esse riportano coincidono con
    i valori discussi nel \autoref{chap:results}, e la generazione è quindi anche un controllo di
    coerenza fra i numeri citati e i dati grezzi.
\end{itemize}
\end{sloppypar}

Le due procedure sono deterministiche: a parità di file di ingresso producono lo stesso output, e
il ricampionamento impiegato per gli intervalli di confidenza usa un seme fissato. I dati
dell'esperimento di ordinamento, che non appartengono a una campagna ma a un controllo dedicato,
sono archiviati con gli altri in \texttt{results/analisi/}.

% =======================================================================================
\section{Verifica di equivalenza della ricompilazione}\label{sec:instrumentation}

La caratterizzazione del costo computazionale (\autoref{sec:cost}) ha richiesto di ricompilare il
sistema, distribuito senza sorgenti. La procedura adottata è la seguente.

\begin{enumerate}[leftmargin=1.6em]
    \item \textbf{Decompilazione} delle classi applicative dell'archivio.
    \item \textbf{Ricompilazione} per la versione \num{52} del formato \texttt{class}, contro le
    medesime librerie incorporate nell'archivio originale.
    \item \textbf{Ricostruzione} dell'archivio sostituendo \emph{soltanto} le classi applicative
    dentro una copia dell'originale: le restanti classi di libreria restano quelle di partenza.
    \item \textbf{Verifica di equivalenza} eseguendo l'archivio originale e quello ricostruito
    sulla medesima configurazione, in cartelle isolate, e confrontando gli output.
\end{enumerate}

L'esito della verifica è riportato nella \autoref{tab:equivalence}: i file di output coincidono
byte per byte, e le uniche differenze riguardano i tempi di esecuzione registrati, che variano
fisiologicamente fra esecuzioni.

\begin{table}[htbp]
\centering
\caption{Verifica di equivalenza della ricompilazione. Configurazione
\texttt{bridges\_70\_1}, soglia 15.}
\label{tab:equivalence}
\footnotesize
\begin{tabular}{lll}
\toprule
file & dimensione & esito del confronto \\
\midrule
\texttt{Populated} & \SI{1377}{\byte}  & \textbf{identico} (checksum \texttt{682428f09943}) \\
\texttt{Results}   & \SI{11522}{\byte} & \textbf{identico} (checksum \texttt{8eca4ca48f2a}) \\
\texttt{Logs}      & ---               & differisce \emph{solo} nelle righe di tempo \\
\bottomrule
\end{tabular}
\end{table}

Solo dopo questa verifica il sistema è stato strumentato, aggiungendo contatori e misure di tempo
mediante \emph{wrapper} sui metodi di interesse anziché modificandone i corpi, in modo che il
flusso di controllo restasse quello originale.

% =======================================================================================
\section{Comandi principali}\label{sec:appendix-commands}

Le campagne sono orchestrate da script che invocano il sistema di valutazione. A titolo
esemplificativo, la valutazione di una strategia su una famiglia di dataset richiede:

\begin{quote}
\texttt{python3 pruning/benchmark.py \textbackslash}\\
\texttt{\ \ --datasets <elenco> --runs <elenco> --thresholds <elenco> \textbackslash}\\
\texttt{\ \ --variants "<nome>=<strategia>:<politica>" \textbackslash}\\
\texttt{\ \ --out results/<nome>.json}
\end{quote}

\noindent Il sistema di valutazione provvede a eseguire la configurazione di riferimento, a
produrre il file di regole potato per ciascuna variante, a eseguire il sistema su entrambi e a
calcolare le metriche in forma appaiata. I risultati sono memorizzati incrementalmente, cosicché
una campagna interrotta può essere ripresa.

% =======================================================================================
\section{Ambiente di esecuzione}\label{sec:appendix-env}

Le campagne sono state condotte su una macchina a otto core, con esecuzioni strettamente
sequenziali. La sequenzialità non è una scelta di comodo: la cartella di lavoro del sistema è
condivisa fra esecuzioni, ed esecuzioni simultanee sovrascrivono reciprocamente gli output
producendo risultati errati senza segnalare alcun errore. La circostanza è stata verificata
sperimentalmente.

Le versioni delle librerie impiegate per la ricompilazione sono fissate e documentate, poiché
versioni diverse possono produrre differenze nel comportamento della generazione dei numeri
casuali e quindi nei risultati.

        % =======================================================================================
% Appendice B: tabelle complete dei risultati
% =======================================================================================
\chapter{Tabelle complete dei risultati}\label{app:results}

\noindent Questa appendice raccoglie le tabelle complete delle campagne, che il
\autoref{chap:results} riporta in forma aggregata. Le tabelle sono \emph{generate} dai file di
risultato delle campagne dalla procedura \texttt{pruning/make\_appendix\_tables.py}, e i totali che
esse riportano coincidono con i valori discussi nel corpo della tesi: la generazione è quindi
anche una verifica di coerenza fra i numeri citati e i dati grezzi, con la stessa logica della
procedura di verifica automatica descritta nell'appendice sulla riproducibilità.

% =======================================================================================
\section{Campagna principale su Bridges}\label{sec:app-bridges}

La \autoref{tab:app-bridges} riporta le \num{50} configurazioni della campagna principale, con la
medesima ricetta di accoppiamento impiegata nei test: per ogni configurazione la baseline e la
variante \str{A5b}, eseguite sullo stesso dataset, alla stessa soglia, con lo stesso seme. I
totali in fondo alla tabella --- \num{69} imputazioni corrette in più e \num{72} errori in meno
--- sono quelli discussi nella \autoref{sec:res-a5b}.

% FILE GENERATO da pruning/make_appendix_tables.py — non modificare a mano.
\footnotesize
\begin{longtable}{@{}llrrrrrr@{}}
\caption{Campagna principale su Bridges, configurazione per configurazione. ``RFD'' sono le regole in ingresso e in uscita dalla potatura, ``rid.''\ la riduzione percentuale, $\Delta$~corr.\ e $\Delta$~err.\ le variazioni rispetto alla medesima configurazione non potata.}\label{tab:app-bridges}\\
\toprule
dataset & soglia & run & RFD in & RFD out & rid.\ (\si{\percent}) & $\Delta$ corr. & $\Delta$ err. \\
\midrule
\endfirsthead
\toprule
dataset & soglia & run & RFD in & RFD out & rid.\ (\si{\percent}) & $\Delta$ corr. & $\Delta$ err. \\
\midrule
\endhead
\midrule
\multicolumn{5}{r}{totale} & & \textbf{+69} & \textbf{-72} \\
\bottomrule
\endlastfoot
bridges\_14 & 9 & 1 & 8961 & 1366 & 84,8 & +1 & -1 \\
bridges\_14 & 9 & 2 & 8961 & 1381 & 84,6 & +0 & +0 \\
bridges\_14 & 9 & 3 & 8961 & 1388 & 84,5 & +0 & +0 \\
bridges\_14 & 9 & 4 & 8961 & 1368 & 84,7 & +0 & +0 \\
bridges\_14 & 9 & 5 & 8961 & 1391 & 84,5 & +0 & +0 \\
bridges\_28 & 9 & 1 & 8961 & 1395 & 84,4 & +0 & -1 \\
bridges\_28 & 9 & 2 & 8961 & 1406 & 84,3 & +0 & +0 \\
bridges\_28 & 9 & 3 & 8961 & 1410 & 84,3 & -2 & -1 \\
bridges\_28 & 9 & 4 & 8961 & 1415 & 84,2 & -1 & +0 \\
bridges\_28 & 9 & 5 & 8961 & 1427 & 84,1 & -1 & +0 \\
bridges\_42 & 9 & 1 & 8961 & 1450 & 83,8 & +2 & -1 \\
bridges\_42 & 9 & 2 & 8961 & 1548 & 82,7 & +8 & -5 \\
bridges\_42 & 9 & 3 & 8961 & 1545 & 82,8 & +1 & -2 \\
bridges\_42 & 9 & 4 & 8961 & 1566 & 82,5 & +1 & -2 \\
bridges\_42 & 9 & 5 & 8961 & 1433 & 84,0 & -2 & +0 \\
bridges\_56 & 9 & 1 & 8961 & 1718 & 80,8 & +0 & -2 \\
bridges\_56 & 9 & 2 & 8961 & 1854 & 79,3 & -3 & +1 \\
bridges\_56 & 9 & 3 & 8961 & 1877 & 79,1 & +1 & -1 \\
bridges\_56 & 9 & 4 & 8961 & 1568 & 82,5 & -1 & -2 \\
bridges\_56 & 9 & 5 & 8961 & 1628 & 81,8 & +5 & -5 \\
bridges\_70 & 9 & 1 & 8961 & 2248 & 74,9 & +5 & -5 \\
bridges\_70 & 9 & 2 & 8961 & 2317 & 74,1 & +6 & -5 \\
bridges\_70 & 9 & 3 & 8961 & 2362 & 73,6 & +3 & -3 \\
bridges\_70 & 9 & 4 & 8961 & 2266 & 74,7 & +5 & -3 \\
bridges\_70 & 9 & 5 & 8961 & 2808 & 68,7 & +4 & -2 \\
bridges\_14 & 15 & 1 & 18844 & 1853 & 90,2 & +2 & -1 \\
bridges\_14 & 15 & 2 & 18844 & 1881 & 90,0 & +0 & +1 \\
bridges\_14 & 15 & 3 & 18844 & 1891 & 90,0 & +0 & +0 \\
bridges\_14 & 15 & 4 & 18844 & 1849 & 90,2 & +1 & +1 \\
bridges\_14 & 15 & 5 & 18844 & 1896 & 89,9 & -1 & +0 \\
bridges\_28 & 15 & 1 & 18844 & 1904 & 89,9 & +2 & -2 \\
bridges\_28 & 15 & 2 & 18844 & 1918 & 89,8 & +2 & -1 \\
bridges\_28 & 15 & 3 & 18844 & 1923 & 89,8 & -3 & -1 \\
bridges\_28 & 15 & 4 & 18844 & 1934 & 89,7 & -1 & +0 \\
bridges\_28 & 15 & 5 & 18844 & 1961 & 89,6 & -1 & +0 \\
bridges\_42 & 15 & 1 & 18844 & 1994 & 89,4 & +3 & +0 \\
bridges\_42 & 15 & 2 & 18844 & 2139 & 88,6 & +7 & -3 \\
bridges\_42 & 15 & 3 & 18844 & 2158 & 88,5 & +1 & +0 \\
bridges\_42 & 15 & 4 & 18844 & 2177 & 88,4 & +0 & -2 \\
bridges\_42 & 15 & 5 & 18844 & 1974 & 89,5 & -4 & -1 \\
bridges\_56 & 15 & 1 & 18844 & 2438 & 87,1 & +7 & -4 \\
bridges\_56 & 15 & 2 & 18844 & 2663 & 85,9 & -3 & +1 \\
bridges\_56 & 15 & 3 & 18844 & 2707 & 85,6 & -1 & -1 \\
bridges\_56 & 15 & 4 & 18844 & 2200 & 88,3 & -1 & -2 \\
bridges\_56 & 15 & 5 & 18844 & 2253 & 88,0 & +7 & -4 \\
bridges\_70 & 15 & 1 & 18844 & 3524 & 81,3 & +6 & -4 \\
bridges\_70 & 15 & 2 & 18844 & 3472 & 81,6 & +0 & +0 \\
bridges\_70 & 15 & 3 & 18844 & 3672 & 80,5 & +4 & -4 \\
bridges\_70 & 15 & 4 & 18844 & 3413 & 81,9 & +5 & -2 \\
bridges\_70 & 15 & 5 & 18844 & 4502 & 76,1 & +5 & -3 \\
\end{longtable}

\begin{table}[htbp]
\centering
\caption{Ablazione di \str{A5b} sulle 20 configurazioni comuni ai quattro bracci. ``priorità'' è il solo riordino, ``adattiva'' la sola potatura adattiva; la somma delle due variazioni è +5, la variazione della combinazione è +45.}
\label{tab:app-ablation}
\footnotesize
\begin{tabular}{@{}llrrrr@{}}
\toprule
dataset & soglia & $\Delta$ priorità & $\Delta$ adattiva & somma & $\Delta$ \str{A5b} \\
\midrule
bridges\_14 & 9 & +1 & +0 & +1 & +1 \\
bridges\_14 & 9 & +0 & +0 & +0 & +0 \\
bridges\_28 & 9 & +1 & -1 & +0 & +0 \\
bridges\_28 & 9 & +1 & -1 & +0 & +0 \\
bridges\_42 & 9 & +2 & -1 & +1 & +2 \\
bridges\_42 & 9 & +9 & -1 & +8 & +8 \\
bridges\_56 & 9 & +3 & -4 & -1 & +0 \\
bridges\_56 & 9 & +2 & -4 & -2 & -3 \\
bridges\_70 & 9 & +7 & -2 & +5 & +5 \\
bridges\_70 & 9 & +7 & +0 & +7 & +6 \\
bridges\_14 & 15 & +1 & +1 & +2 & +2 \\
bridges\_14 & 15 & +0 & +0 & +0 & +0 \\
bridges\_28 & 15 & +2 & +1 & +3 & +2 \\
bridges\_28 & 15 & +2 & +0 & +2 & +2 \\
bridges\_42 & 15 & +2 & -1 & +1 & +3 \\
bridges\_42 & 15 & +7 & +0 & +7 & +7 \\
bridges\_56 & 15 & +8 & +1 & +9 & +7 \\
bridges\_56 & 15 & +1 & -5 & -4 & -3 \\
bridges\_70 & 15 & -33 & -2 & -35 & +6 \\
bridges\_70 & 15 & +2 & -1 & +1 & +0 \\
\midrule
\multicolumn{2}{r}{totale} & +25 & -20 & +5 & +45 \\
\bottomrule
\end{tabular}
\end{table}

\scriptsize
\begin{longtable}{@{}llrrrr@{}}
\caption{Taratura conservativa della potatura aggressiva ($\alpha = 0{,}99$) sulle 97 configurazioni con baseline disponibile. Riduzione mediana 0,0\si{\percent}, massima 52,9\si{\percent}; variazioni complessive $\Delta$ corr.\ -1, $\Delta$ err.\ -3.}
\label{tab:app-a99}\\
\toprule
dataset & soglia & run & rid.\ (\si{\percent}) & $\Delta$ corr. & $\Delta$ err. \\
\midrule
\endfirsthead
\toprule
dataset & soglia & run & rid.\ (\si{\percent}) & $\Delta$ corr. & $\Delta$ err. \\
\midrule
\endhead
\bottomrule
\endlastfoot
glass\_107 & 3 & 1 & 0,0 & +0 & +0 \\
glass\_107 & 3 & 2 & 0,0 & +0 & +0 \\
glass\_107 & 3 & 3 & 0,0 & +0 & +0 \\
glass\_107 & 3 & 4 & 0,0 & +0 & +0 \\
glass\_107 & 3 & 5 & 0,0 & +0 & +0 \\
glass\_22 & 3 & 1 & 24,1 & +0 & +0 \\
glass\_22 & 3 & 2 & 15,0 & +0 & +0 \\
glass\_22 & 3 & 3 & 23,5 & +0 & +0 \\
glass\_22 & 3 & 4 & 32,3 & +0 & +0 \\
glass\_22 & 3 & 5 & 41,1 & +0 & +0 \\
glass\_43 & 3 & 1 & 1,1 & +0 & +0 \\
glass\_43 & 3 & 2 & 0,4 & +0 & +0 \\
glass\_43 & 3 & 3 & 0,9 & +0 & +0 \\
glass\_43 & 3 & 4 & 1,3 & +0 & +0 \\
glass\_43 & 3 & 5 & 1,1 & +0 & +0 \\
glass\_64 & 3 & 1 & 0,0 & +0 & +0 \\
glass\_64 & 3 & 2 & 0,3 & +0 & +0 \\
glass\_64 & 3 & 3 & 0,5 & +0 & +0 \\
glass\_64 & 3 & 4 & 0,0 & +0 & +0 \\
glass\_64 & 3 & 5 & 0,3 & +0 & +0 \\
glass\_86 & 3 & 1 & 0,1 & +0 & +0 \\
glass\_86 & 3 & 2 & 0,0 & +0 & +0 \\
glass\_86 & 3 & 3 & 0,0 & +0 & +0 \\
glass\_86 & 3 & 4 & 0,0 & +0 & +0 \\
glass\_86 & 3 & 5 & 0,1 & +0 & +0 \\
restaurant\_104 & 3 & 1 & 0,0 & +0 & +0 \\
restaurant\_104 & 3 & 2 & 4,0 & +0 & +0 \\
restaurant\_104 & 3 & 3 & 0,0 & +0 & +0 \\
restaurant\_206 & 3 & 1 & 0,0 & +0 & +0 \\
restaurant\_206 & 3 & 2 & 0,0 & +0 & +0 \\
restaurant\_206 & 3 & 3 & 0,0 & +0 & +0 \\
restaurant\_259 & 3 & 1 & 0,0 & +0 & +0 \\
restaurant\_259 & 3 & 2 & 0,0 & +0 & +0 \\
restaurant\_52 & 3 & 1 & 16,0 & +0 & +0 \\
restaurant\_52 & 3 & 2 & 24,0 & +0 & +0 \\
restaurant\_52 & 3 & 3 & 16,0 & +0 & +0 \\
restaurant\_104 & 6 & 1 & 5,6 & +0 & +0 \\
restaurant\_104 & 6 & 2 & 9,7 & +0 & +0 \\
restaurant\_104 & 6 & 3 & 0,0 & +0 & +0 \\
restaurant\_206 & 6 & 1 & 0,0 & +0 & -1 \\
restaurant\_206 & 6 & 2 & 0,0 & +0 & +0 \\
restaurant\_206 & 6 & 3 & 0,0 & +0 & +0 \\
restaurant\_259 & 6 & 1 & 0,0 & +0 & +0 \\
restaurant\_259 & 6 & 2 & 0,0 & -1 & -1 \\
restaurant\_52 & 6 & 1 & 37,9 & +0 & +0 \\
restaurant\_52 & 6 & 2 & 45,2 & +0 & +0 \\
restaurant\_52 & 6 & 3 & 45,2 & +0 & +0 \\
glass\_107 & 9 & 1 & 0,0 & +0 & +0 \\
glass\_107 & 9 & 2 & 0,0 & +0 & +0 \\
glass\_107 & 9 & 3 & 0,0 & +0 & +0 \\
glass\_107 & 9 & 4 & 0,0 & +0 & +0 \\
glass\_107 & 9 & 5 & 0,0 & +0 & +0 \\
glass\_22 & 9 & 1 & 31,7 & +0 & +0 \\
glass\_22 & 9 & 2 & 17,4 & +0 & +0 \\
glass\_22 & 9 & 3 & 24,8 & +0 & +0 \\
glass\_22 & 9 & 4 & 38,7 & +0 & +0 \\
glass\_22 & 9 & 5 & 52,9 & +0 & +0 \\
glass\_43 & 9 & 1 & 0,4 & +0 & +0 \\
glass\_43 & 9 & 2 & 0,2 & +0 & +0 \\
glass\_43 & 9 & 3 & 0,5 & +0 & +0 \\
glass\_43 & 9 & 4 & 0,6 & +0 & +0 \\
glass\_43 & 9 & 5 & 0,4 & +0 & +0 \\
glass\_64 & 9 & 1 & 0,0 & +0 & +0 \\
glass\_64 & 9 & 2 & 0,1 & +0 & +0 \\
glass\_64 & 9 & 3 & 0,2 & +0 & +0 \\
glass\_64 & 9 & 4 & 0,0 & +0 & +0 \\
glass\_64 & 9 & 5 & 0,0 & +0 & +0 \\
glass\_86 & 9 & 1 & 0,0 & +0 & +0 \\
glass\_86 & 9 & 2 & 0,0 & +0 & +0 \\
glass\_86 & 9 & 3 & 0,0 & +0 & +0 \\
glass\_86 & 9 & 4 & 0,0 & +0 & +0 \\
glass\_86 & 9 & 5 & 0,0 & +0 & +0 \\
glass\_107 & 15 & 1 & 0,0 & +0 & +0 \\
glass\_107 & 15 & 2 & 0,0 & +0 & +0 \\
glass\_107 & 15 & 3 & 0,0 & +0 & +0 \\
glass\_107 & 15 & 4 & 0,0 & +0 & +0 \\
glass\_107 & 15 & 5 & 0,0 & +0 & +0 \\
glass\_22 & 15 & 1 & 34,6 & +0 & +0 \\
glass\_22 & 15 & 2 & 20,3 & +0 & +0 \\
glass\_22 & 15 & 3 & 23,2 & +0 & +0 \\
glass\_22 & 15 & 4 & 36,0 & +0 & -1 \\
glass\_22 & 15 & 5 & 51,6 & +0 & +0 \\
glass\_43 & 15 & 1 & 0,3 & +0 & +0 \\
glass\_43 & 15 & 2 & 0,2 & +0 & +0 \\
glass\_43 & 15 & 3 & 0,5 & +0 & +0 \\
glass\_43 & 15 & 4 & 0,5 & +0 & +0 \\
glass\_43 & 15 & 5 & 0,2 & +0 & +0 \\
glass\_64 & 15 & 1 & 0,0 & +0 & +0 \\
glass\_64 & 15 & 2 & 0,1 & +0 & +0 \\
glass\_64 & 15 & 3 & 0,1 & +0 & +0 \\
glass\_64 & 15 & 4 & 0,0 & +0 & +0 \\
glass\_64 & 15 & 5 & 0,0 & +0 & +0 \\
glass\_86 & 15 & 1 & 0,0 & +0 & +0 \\
glass\_86 & 15 & 2 & 0,0 & +0 & +0 \\
glass\_86 & 15 & 3 & 0,0 & +0 & +0 \\
glass\_86 & 15 & 4 & 0,0 & +0 & +0 \\
glass\_86 & 15 & 5 & 0,0 & +0 & +0 \\
\end{longtable}


% =======================================================================================
\section{Studio di ablazione}\label{sec:app-ablation}

La \autoref{tab:app-ablation} riporta le venti configurazioni comuni ai quattro bracci
dell'ablazione. La tabella rende verificabile la lettura discussa nella \autoref{sec:ablation}: la
somma delle variazioni dei due componenti isolati è pari a $+5$, mentre la variazione della
combinazione è pari a $+45$, e la differenza ($+40$) è la sinergia che il \autoref{sec:mechanism}
spiega.

% =======================================================================================
\section{Taratura conservativa}\label{sec:app-a99}

La \autoref{tab:app-a99} riporta le \num{97} configurazioni della campagna con il parametro di
aggressività portato al valore conservativo. Vi compare il dato su cui poggia la conclusione della
\autoref{sec:res-boundary}: la riduzione mediana è nulla, la massima è pari al
\SI{52.9}{\percent}, e le variazioni complessive sono di una imputazione corretta in meno e di tre
errori in meno. È questa la ragione per cui la strategia, in quella taratura, viene descritta come
\emph{non dannosa} e non come vantaggiosa: l'effetto sulla qualità non è distinguibile da zero
(\autoref{tab:statistics}), mentre quello sulla riduzione è concentrato in poche configurazioni.

    \end{appendices}

    % NOTA: la spedizione dell'ultima pagina dell'appendice avviene al \cleardoublepage interno
    % di \printbibliography; se il cambio di stile precedesse quella spedizione, l'ultima pagina
    % dell'appendice porterebbe l'intestazione della bibliografia. Il \cleardoublepage esplicito
    % la spedisce prima del cambio.
    \newpage
    \cleardoublepage
    \pagestyle{backmatter}
    \printbibliography[heading=bibintoc]

\end{document}
